%!TEX TS-program = xelatex

  %%%%%%%%%%%%%%%%%%%%%%%%%%%%%%%%%%%%%%%%%
% The Legrand Orange Book
% LaTeX Template
% Version 1.4 (12/4/14)
%
% This template has been downloaded from:
% http://www.LaTeXTemplates.com
%
% Original author:
% Mathias Legrand (legrand.mathias@gmail.com)
%
% License:
% CC BY-NC-SA 3.0 (http://creativecommons.org/licenses/by-nc-sa/3.0/)
%
% Compiling this template:
% This template uses biber for its bibliography and makeindex for its index.
% When you first open the template, compile it from the command line with the 
% commands below to make sure your LaTeX distribution is configured correctly:
%
% 1) pdflatex main
% 2) makeindex main.idx -s StyleInd.ist
% 3) biber main
% 4) pdflatex main x 2
%
% After this, when you wish to update the bibliography/index use the appropriate
% command above and make sure to compile with pdflatex several times 
% afterwards to propagate your changes to the document.
%
% This template also uses a number of packages which may need to be
% updated to the newest versions for the template to compile. It is strongly
% recommended you update your LaTeX distribution if you have any
% compilation errors.
%
% Important note:
% Chapter heading images should have a 2:1 width:height ratio,
% e.g. 920px width and 460px height.
%
%%%%%%%%%%%%%%%%%%%%%%%%%%%%%%%%%%%%%%%%%

%-=-=-=-=-=-=-=-=-=-=-=-=-=-=-=-=-=-=-=-=-=-=-=-=
%
%	PACKAGES AND OTHER DOCUMENT CONFIGURATIONS
%
%-=-=-=-=-=-=-=-=-=-=-=-=-=-=-=-=-=-=-=-=-=-=-=-=

\documentclass[11pt,fleqn]{book} % Default font size and left-justified equations

\usepackage[top=3cm,bottom=3cm,left=3.2cm,right=3.2cm,headsep=10pt,a4paper]{geometry} % Page margins

\usepackage{xcolor} % Required for specifying colors by name
\definecolor{ocre}{RGB}{0,126,196} % Define the orange color used for highlighting throughout the book
\newcommand{\alert}[1]{{\color{stockholmPink}{#1}}}
\newcommand{\mathref}[1]{{\color{stockholmPurple}{#1}}}

% Define Stockholms Stad Colors
\definecolor{stockholmPink}{RGB}{196,0,104}
% HEX #c40064
\definecolor{stockholmLightPink}{RGB}{254,222,237}
% HEX #fedeed
\definecolor{stockholmGreen}{RGB}{40,157,147}
% HEX #00867f
\definecolor{stockholmLightGreen}{RGB}{182,215,211}
% HEX #d5f7f4
\definecolor{stockholmOrange}{RGB}{203,80,25}
% HEX #dd4a2c
\definecolor{stockholmLightOrange}{RGB}{233,189,155}
% HEX #ffd7d2
\definecolor{stockholmBlue}{RGB}{0,126,196}
% HEX #006ebf
\definecolor{stockholmLightBlue}{RGB}{172,199,233}
% HEX #d6edfc
\definecolor{stockholmPurple}{RGB}{104,55,136}
% HEX #5d237d
\definecolor{stockholmLightPurple}{RGB}{188,170,208}
% HEX #f1e6fc
\definecolor{stockholmYellow}{RGB}{252,191,10}
% HEX #fcbf0a
\definecolor{stockholmGrey}{RGB}{245,243,238}
% HEX #f5f3ee
\definecolor{stockholmDarkGrey}{RGB}{51,51,51}
% HEX #333333

\usepackage{multicol}

%-=-=-=-=-=-=-=-=-=-=-=-=-=-=-=-=-=-=-=-=-=-=-=-=
%	FONT SETTINGS
%-=-=-=-=-=-=-=-=-=-=-=-=-=-=-=-=-=-=-=-=-=-=-=-=

%\usepackage{avant} % Use the Avantgarde font for headings
\usepackage{times} % Use the Times font for headings
\usepackage{mathpazo}
%\usepackage{mathptmx} % Use the Adobe Times Roman as the default text font together with math symbols from the Sym­bol, Chancery and Com­puter Modern fonts
\usepackage{microtype} % Slightly tweak font spacing for aesthetics
\usepackage[utf8]{inputenc} % Required for including letters with accents
\usepackage[T1]{fontenc} % Use 8-bit encoding that has 256 glyphs
\usepackage{xspace} % Used for printing a trailing space better than using a tilde (~) using the \xspace command

%-=-=-=-=-=-=-=-=-=-=-=-=-=-=-=-=-=-=-=-=-=-=-=-=
%	BIBLIOGRAPHY
%-=-=-=-=-=-=-=-=-=-=-=-=-=-=-=-=-=-=-=-=-=-=-=-=
\usepackage[style=alphabetic,sorting=nyt,sortcites=true,autopunct=true,babel=hyphen,hyperref=true,abbreviate=false,backref=true,backend=biber]{biblatex}
\addbibresource{bibliography.bib} % BibTeX bibliography file
\defbibheading{bibempty}{}

%-=-=-=-=-=-=-=-=-=-=-=-=-=-=-=-=-=-=-=-=-=-=-=-=
%	INDEX
%-=-=-=-=-=-=-=-=-=-=-=-=-=-=-=-=-=-=-=-=-=-=-=-=
\usepackage{calc} % For simpler calculation - used for spacing the index letter headings correctly
\usepackage{makeidx} % Required to make an index
\makeindex % Tells LaTeX to create the files required for indexing

%-=-=-=-=-=-=-=-=-=-=-=-=-=-=-=-=-=-=-=-=-=-=-=-=
%	MATHEMATICS SETTINGS
%-=-=-=-=-=-=-=-=-=-=-=-=-=-=-=-=-=-=-=-=-=-=-=-=

\usepackage{
amsmath,
amssymb,
amsfonts,
amsthm,
cancel,
colortbl,
longtable,
pdflscape,
sagetex,
setspace,
mathtools,
pgfplots
}


%-=-=-=-=-=-=-=-=-=-=-=-=-=-=-=-=-=-=-=-=-=-=-=-=
%	TIKZ LIBRARIES
%-=-=-=-=-=-=-=-=-=-=-=-=-=-=-=-=-=-=-=-=-=-=-=-=

%-=-=-=-=-=-=-=-=-=-=-=-=-=-=-=-=-=-=-=-=-=-=-=-=
%	TIKZ LIBRARIES
%-=-=-=-=-=-=-=-=-=-=-=-=-=-=-=-=-=-=-=-=-=-=-=-=

\usetikzlibrary{
arrows,
backgrounds,
fit,
chains,
calc,
decorations.pathmorphing,
matrix,
mindmap,
positioning,
shapes,
shadows,
trees
}

\tikzset{ >=stealth', help lines/.style={dashed, thick}, axis/.style={<->}, important line/.style={thick}, connection/.style={thick, dotted},}

\tikzstyle{firstterm} = [circle, draw, fill=sthlmRed!20, text centered,minimum size=1cm]
\tikzstyle{secondterm} = [circle, draw, fill=sthlmBlue!20, text centered, node distance=0.25cm,minimum size=1cm]
\tikzstyle{terms} = [rectangle, draw, fill=sthlmDarkGrey!20, text centered, node distance=0.25cm,minimum size=1cm]
\tikzstyle{factoradd} = [rectangle, draw=none, fill=white, text centered, node distance=0.25cm,minimum size=1cm]
\tikzstyle{multiply} = [rectangle, draw, fill=sthlmDarkGrey!20, text centered, minimum height=0.75cm,minimum size=1cm]
\tikzstyle{add} = [rectangle, draw, fill=sthlmGreen!20, text centered, minimum height=0.75cm,minimum size=1cm]
\tikzstyle{line} = [draw, color=black!50, -latex']
% Define the style for the red dotted boxes
\tikzset{factordotted/.style={draw=sthlmPurple!50!white, line width=1pt,
                               dash pattern=on 1pt off 4pt on 6pt off 4pt,
                                inner sep=4mm, rectangle, rounded corners}}
%-=-=-=-=-=-=-=-=-=-=-=-=-=-=-=-=-=-=-=-=-=-=-=-=
%	SIMPLIFYING EXPRESSIONS
%-=-=-=-=-=-=-=-=-=-=-=-=-=-=-=-=-=-=-=-=-=-=-=-=

\tikzstyle{operation} = [rectangle, draw, fill=stockholmBlue!40, text centered, rounded corners]
\tikzstyle{function} = [rectangle, draw, fill=stockholmPurple!40, text centered, rounded corners]
\tikzstyle{notation} = [rectangle, draw, fill=stockholmGreen!40, text centered, rounded corners]
\tikzstyle{algorithm} = [rectangle, draw, fill=stockholmOrange!40, text centered, rounded corners]
\tikzstyle{delim} = [rectangle, draw, fill=stockholmYellow!40, text centered, rounded corners]
\tikzstyle{property} = [rectangle, draw, fill=stockholmPink!40, text centered, rounded corners]
\tikzstyle{hide} = [rectangle, draw, fill=stockholmGrey, text centered, rounded corners]
\tikzstyle{active} = [rectangle, draw, fill=stockholmDarkGrey, text=white, text centered, rounded corners]
\tikzstyle{line} = [draw, very thick, color=black!50, -latex']

\tikzset{functiondotted/.style={draw=stockholmPurple!50!white, line width=1pt,
                               dash pattern=on 1pt off 4pt on 6pt off 4pt,
                                inner sep=4mm, rectangle, rounded corners}}
								
\tikzset{productdotted/.style={draw=stockholmBlue!50!white, line width=1pt,
                               dash pattern=on 1pt off 4pt on 6pt off 4pt,
                                inner sep=4mm, rectangle, rounded corners}}
								
\tikzset{sumdotted/.style={draw=stockholmPink!50!white, line width=1pt,
                               dash pattern=on 1pt off 4pt on 6pt off 4pt,
                                inner sep=4mm, rectangle, rounded corners}}


%-=-=-=-=-=-=-=-=-=-=-=-=-=-=-=-=-=-=-=-=-=-=-=-=
%	MATH COMMANDS
%-=-=-=-=-=-=-=-=-=-=-=-=-=-=-=-=-=-=-=-=-=-=-=-=

\newcommand{\unit}[1]{\ensuremath{\, \mathrm{#1}}}
\newcommand{\limit}{\displaystyle\lim}
\newcommand{\impart}{\textrm{Im}}
\newcommand{\repart}{\textrm{Re}}
\newcommand{\integer}{\mathbb{Z}}
\newcommand{\real}{\mathbb{R}}
\newcommand{\degree}{^{\circ}}
\newcommand{\rad}{\unit{rad}}
\newcommand{\dd}{\mathrm{d}}
\newcommand{\dx}{\mathrm{d}x}
\newcommand{\dy}{\mathrm{d}y}
\newcommand{\dr}{\mathrm{d}r}
\newcommand{\dtheta}{\mathrm{d}\theta}
\newcommand{\dydx}{\dfrac{\mathrm{d}y}{\mathrm{d}x}}
\newcommand{\dydt}{\dfrac{\mathrm{d}y}{\mathrm{d}t}}
\newcommand{\dydu}{\dfrac{\mathrm{d}y}{\mathrm{d}u}}
\newcommand{\dudx}{\dfrac{\mathrm{d}u}{\mathrm{d}x}}
\newcommand{\dxdt}{\dfrac{\mathrm{d}x}{\mathrm{d}t}}
\newcommand{\vect}[1]{\boldsymbol{#1}}
\newcommand{\farg}[1]{{\color{stockholmBlue}{\left[\boldsymbol{#1}\right]}}}
\newcommand{\cBlue}[1]{{\color{stockholmBlue}{#1}}}
\newcommand{\cGreen}[1]{{\color{stockholmGreen}{#1}}}
\newcommand{\cYellow}[1]{{\color{stockholmYellow}{#1}}}
\newcommand{\cPurple}[1]{{\color{stockholmPurple}{#1}}}
\newcommand{\cOrange}[1]{{\color{stockholmOrange}{#1}}}
\newcommand{\chide}[1]{{\color{stockholmDarkGrey}{#1}}}
\newcommand{\cGrey}[1]{{\color{stockholmGrey}{#1}}}
\newcommand{\cDarkGrey}[1]{{\color{stockholmDarkGrey}{#1}}}
\newcommand{\cDarkBlue}[1]{{\color{stockholmDarkBlue}{#1}}}
\newcommand{\cLightBlue}[1]{{\color{stockholmLightBlue}{#1}}}
\newcommand{\cRed}[1]{{\color{stockholmPink}{#1}}}
\newcommand{\cLightRed}[1]{{\color{stockholmPink}{#1}}}
\newcommand{\sname}[1]{\texttt{#1}}
\newcommand{\inv}{^{-1}}
\newcommand{\cb}{$\Box$}
\newcommand{\lesssteps}{\cGreen{\textbf{Less Steps Solution:}}}
\newcommand{\solnsteps}{\cGreen{\textbf{Solution:}}}
\newcommand{\qdependlist}{\cDarkGrey{\textbf{Dependencies:}}}
\newcommand{\suchthat}{\;\ifnum\currentgrouptype=16 \middle\fi|\;}
% used in tables
\newcommand{\cellyes}{\cellcolor{stockholmGreen!30}Yes}  
\newcommand{\cellno}{\cellcolor{stockholmPink!30}No}  
\newcommand{\ihat}{\hat{\imath}}
\newcommand{\jhat}{\hat{\jmath}} 
%-=-=-=-=-=-=-=-=-=-=-=-=-=-=-=-=-=-=-=-=-=-=-=-=
%
%	IMPORT STRUCTURE
%
%-=-=-=-=-=-=-=-=-=-=-=-=-=-=-=-=-=-=-=-=-=-=-=-=

%%%%%%%%%%%%%%%%%%%%%%%%%%%%%%%%%%%%%%%%%
% The Legrand Orange Book
% Structural Definitions File
% Version 2.0 (9/2/15)
%
% Original author:
% Mathias Legrand (legrand.mathias@gmail.com) with modifications by:
% Vel (vel@latextemplates.com)
% 
% This file has been downloaded from:
% http://www.LaTeXTemplates.com
%
% License:
% CC BY-NC-SA 3.0 (http://creativecommons.org/licenses/by-nc-sa/3.0/)
%
%%%%%%%%%%%%%%%%%%%%%%%%%%%%%%%%%%%%%%%%%

%----------------------------------------------------------------------------------------
%	VARIOUS REQUIRED PACKAGES AND CONFIGURATIONS
%----------------------------------------------------------------------------------------

\usepackage[top=3cm,bottom=3cm,left=3cm,right=3cm,headsep=10pt,a4paper]{geometry} % Page margins

\usepackage{graphicx} % Required for including pictures
\graphicspath{{Pictures/}} % Specifies the directory where pictures are stored

\usepackage{lipsum} % Inserts dummy text

\usepackage{tikz} % Required for drawing custom shapes

\usepackage[english]{babel} % English language/hyphenation

\usepackage{enumitem} % Customize lists
\setlist{nolistsep} % Reduce spacing between bullet points and numbered lists

\usepackage{booktabs} % Required for nicer horizontal rules in tables

\usepackage{xcolor} % Required for specifying colors by name
\definecolor{ocre}{RGB}{0,126,196} % Define the orange color used for highlighting throughout the book

%-=-=-=-=-=-=-=-=-=-=-=-=-=-=-=-=-=-=-=-=-=-=-=-=
%	MATHEMATICS SETTINGS
%-=-=-=-=-=-=-=-=-=-=-=-=-=-=-=-=-=-=-=-=-=-=-=-=

\usepackage{
amsmath,
amssymb,
amsfonts,
amsthm,
cancel,
colortbl,
longtable,
pdflscape,
setspace,
mathtools,
pgfplots
}
%-=-=-=-=-=-=-=-=-=-=-=-=-=-=-=-=-=-=-=-=-=-=-=-=	
%	sthlm COLOR THEME
%-=-=-=-=-=-=-=-=-=-=-=-=-=-=-=-=-=-=-=-=-=-=-=-=

%== BLUE
	\definecolor{sthlmLightBlue}{RGB}{214,237,252} % HEX #d6edfc
		\newcommand{\csthlmLightBlue}[1]{{\color{sthlmLightBlue}{#1}}}
	\definecolor{sthlmBlue}{RGB}{0,110,191} % HEX #006ebf
		\newcommand{\csthlmBlue}[1]{{\color{sthlmBlue}{#1}}}
	%\definecolor{sthlmDarkBlue}{RGB}{} %HEX #
		%\newcommand{\csthlmDarkBlue}[1]{{\color{sthlmDarkBlue}{#1}}}

%== GREEN

	\definecolor{sthlmLightGreen}{RGB}{213,247,244} % HEX #0d5f7f4
		\newcommand{\csthlmLightGreen}[1]{{\color{sthlmLightGreen}{#1}}}
	\definecolor{sthlmGreen}{RGB}{0,134,127} % #00867f
		\newcommand{\csthlmGreen}[1]{{\color{sthlmGreen}{#1}}}

%== GREY
	%\definecolor{sthlmLightGrey}{RGB}{}
		%\newcommand{\csthlmLightGrey}[1]{{\color{sthlmLightGrey}{#1}}}
	\definecolor{sthlmGrey}{RGB}{245,243,238} % HEX #f5f3ee
		\newcommand{\csthlmGrey}[1]{{\color{sthlmGrey}{#1}}}
	\definecolor{sthlmDarkGrey}{RGB}{51,51,51} % HEX #333333
		\newcommand{\csthlmDarkGrey}[1]{{\color{sthlmDarkGrey}{#1}}} 

%== ORANGE
	\definecolor{sthlmLightOrange}{RGB}{255,215,210} % HEX #ffd7d2
		\newcommand{\csthlmLightOrange}[1]{{\color{sthlmLightOrange}{#1}}}
	\definecolor{sthlmOrange}{RGB}{221,74,44} % HEX #dd4a2c
		\newcommand{\csthlmOrange}[1]{{\color{sthlmOrange}{#1}}}

%== PURPLE	

	\definecolor{sthlmLightPurple}{RGB}{241,230,252} % HEX #f1e6fc
		\newcommand{\csthlmLightPurple}[1]{{\color{sthlmLightPurple}{#1}}}
	\definecolor{sthlmPurple}{RGB}{93,35,125} % HEX #5d237d
		\newcommand{\csthlmPurple}[1]{{\color{sthlmPurple}{#1}}}

%== RED
					
	\definecolor{sthlmLightRed}{RGB}{254,222,237} % HEX #c40064
		\newcommand{\csthlmLightRed}[1]{{\color{sthlmLightRed}{#1}}}
	\definecolor{sthlmRed}{RGB}{196,0,100} % HEX #c40064
		\newcommand{\csthlmRed}[1]{{\color{sthlmRed}{#1}}}
	%\definecolor{sthlmDarkRed}{RGB}{} % HEX #fedeed
		%\newcommand{\csthlmDarkRed}[1]{{\color{sthlmDarkRed}{#1}}}	

%== YELLOW
		
	%\definecolor{sthlmLightYellow}{RGB}{}
		%\newcommand{\csthlmLightYellow}[1]{{\color{sthlmLightYellow}{#1}}}
	\definecolor{sthlmYellow}{RGB}{252,191,10} % HEX #fcbf0a
		\newcommand{\csthlmYellow}[1]{{\color{sthlmYellow}{#1}}}

%-=-=-=-=-=-=-=-=-=-=-=-=-=-=-=-=-=-=-=-=-=-=-=-=	
%	ISSR COLORS
%-=-=-=-=-=-=-=-=-=-=-=-=-=-=-=-=-=-=-=-=-=-=-=-=

%== BLUE
	\definecolor{issrBlue}{RGB}{0,111,174} % HEX #0066cc
		\newcommand{\cissrBlue}[1]{{\color{issrBlue}{#1}}}
		
%== GREY
	\definecolor{issrGrey}{RGB}{167,169,172} % HEX #999999
		\newcommand{\cissrGrey}[1]{{\color{issrGrey}{#1}}}
		
%-=-=-=-=-=-=-=-=-=-=-=-=-=-=-=-=-=-=-=-=-=-=-=-=	
%	GLOBAL COLOR THEME
%-=-=-=-=-=-=-=-=-=-=-=-=-=-=-=-=-=-=-=-=-=-=-=-=

%== BLUE

	\newcommand{\cLightBlue}[1]{{\color{sthlmLightBlue}{#1}}}
	\newcommand{\cBlue}[1]{{\color{sthlmBlue}{#1}}}
	\newcommand{\cDarkBlue}[1]{{\color{sthlmDarkBlue}{#1}}}

%== GREEN

	\newcommand{\cLightGreen}[1]{{\color{}{#1}}}
	\newcommand{\cGreen}[1]{{\color{sthlmGreen}{#1}}}
	\newcommand{\cDarkGreen}[1]{{\color{}{#1}}}

%== GREY

	\newcommand{\cLightGrey}[1]{{\color{sthlmLightGrey}{#1}}}
	\newcommand{\cGrey}[1]{{\color{sthlmGrey}{#1}}}
	\newcommand{\cDarkGrey}[1]{{\color{sthlmDarkGrey}{#1}}}

%== ORANGE

	\newcommand{\cLightOrange}[1]{{\color{}{#1}}}
	\newcommand{\cOrange}[1]{{\color{}{#1}}}
	\newcommand{\cDarkOrange}[1]{{\color{}{#1}}}

%== PURPLE

	\newcommand{\cLightPurple}[1]{{\color{}{#1}}}
	\newcommand{\cPurple}[1]{{\color{sthlmPurple}{#1}}}
	\newcommand{\cDarkPurple}[1]{{\color{}{#1}}}

%== RED

	\newcommand{\cLightRed}[1]{{\color{sthlmLightRed}{#1}}}
	\newcommand{\cRed}[1]{{\color{sthlmRed}{#1}}}
	\newcommand{\cDarkRed}[1]{{\color{sthlmDarkRed}{#1}}}

%== YELLOW

	\newcommand{\cLightYellow}[1]{{\color{sthlmLightYellow}{#1}}}
	\newcommand{\cYellow}[1]{{\color{sthlmYellow}{#1}}}
	\newcommand{\cDarkYellow}[1]{{\color{}{#1}}}
	
\newcommand{\alert}[1]{{\color{sthlmRed}{#1}}}
\newcommand{\mathref}[1]{{\color{sthlmPurple}{#1}}}

\usepackage{multicol}
%----------------------------------------------------------------------------------------
%	FONTS
%----------------------------------------------------------------------------------------

\usepackage{avant} % Use the Avantgarde font for headings
%\usepackage{times} % Use the Times font for headings
\usepackage{mathptmx} % Use the Adobe Times Roman as the default text font together with math symbols from the Sym­bol, Chancery and Com­puter Modern fonts

\usepackage{microtype} % Slightly tweak font spacing for aesthetics
\usepackage[utf8]{inputenc} % Required for including letters with accents
\usepackage[T1]{fontenc} % Use 8-bit encoding that has 256 glyphs

%----------------------------------------------------------------------------------------
%	BIBLIOGRAPHY AND INDEX
%----------------------------------------------------------------------------------------

\usepackage[style=alphabetic,citestyle=numeric,sorting=nyt,sortcites=true,autopunct=true,babel=hyphen,hyperref=true,abbreviate=false,backref=true,backend=biber]{biblatex}
\addbibresource{bibliography.bib} % BibTeX bibliography file
\defbibheading{bibempty}{}

\usepackage{calc} % For simpler calculation - used for spacing the index letter headings correctly
\usepackage{makeidx} % Required to make an index
\makeindex % Tells LaTeX to create the files required for indexing

%----------------------------------------------------------------------------------------
%	MAIN TABLE OF CONTENTS
%----------------------------------------------------------------------------------------

\usepackage{titletoc} % Required for manipulating the table of contents

\contentsmargin{0cm} % Removes the default margin

% Part text styling
\titlecontents{part}[0cm]
{\addvspace{20pt}\centering\large\bfseries}
{}
{}
{}

% Chapter text styling
\titlecontents{chapter}[1.25cm] % Indentation
{\addvspace{12pt}\large\sffamily\bfseries} % Spacing and font options for chapters
{\color{ocre!60}\contentslabel[\Large\thecontentslabel]{1.25cm}\color{ocre}} % Chapter number
{\color{ocre}}  
{\color{ocre!60}\normalsize\;\titlerule*[.5pc]{.}\;\thecontentspage} % Page number

% Section text styling
\titlecontents{section}[1.25cm] % Indentation
{\addvspace{3pt}\sffamily\bfseries} % Spacing and font options for sections
{\contentslabel[\thecontentslabel]{1.25cm}} % Section number
{}
{\hfill\color{black}\thecontentspage} % Page number
[]

% Subsection text styling
\titlecontents{subsection}[1.25cm] % Indentation
{\addvspace{1pt}\sffamily\small} % Spacing and font options for subsections
{\contentslabel[\thecontentslabel]{1.25cm}} % Subsection number
{}
{\ \titlerule*[.5pc]{.}\;\thecontentspage} % Page number
[]

% List of figures
\titlecontents{figure}[0em]
{\addvspace{-5pt}\sffamily}
{\thecontentslabel\hspace*{1em}}
{}
{\ \titlerule*[.5pc]{.}\;\thecontentspage}
[]

% List of tables
\titlecontents{table}[0em]
{\addvspace{-5pt}\sffamily}
{\thecontentslabel\hspace*{1em}}
{}
{\ \titlerule*[.5pc]{.}\;\thecontentspage}
[]

%----------------------------------------------------------------------------------------
%	MINI TABLE OF CONTENTS IN PART HEADS
%----------------------------------------------------------------------------------------

% Chapter text styling
\titlecontents{lchapter}[0em] % Indenting
{\addvspace{15pt}\large\sffamily\bfseries} % Spacing and font options for chapters
{\color{ocre}\contentslabel[\Large\thecontentslabel]{1.25cm}\color{ocre}} % Chapter number
{}  
{\color{ocre}\normalsize\sffamily\bfseries\;\titlerule*[.5pc]{.}\;\thecontentspage} % Page number

% Section text styling
\titlecontents{lsection}[0em] % Indenting
{\sffamily\small} % Spacing and font options for sections
{\contentslabel[\thecontentslabel]{1.25cm}} % Section number
{}
{}

% Subsection text styling
\titlecontents{lsubsection}[.5em] % Indentation
{\normalfont\footnotesize\sffamily} % Font settings
{}
{}
{}

%----------------------------------------------------------------------------------------
%	PAGE HEADERS
%----------------------------------------------------------------------------------------

\usepackage{fancyhdr} % Required for header and footer configuration

\pagestyle{fancy}
\renewcommand{\chaptermark}[1]{\markboth{\sffamily\normalsize\bfseries\chaptername\ \thechapter.\ #1}{}} % Chapter text font settings
\renewcommand{\sectionmark}[1]{\markright{\sffamily\normalsize\thesection\hspace{5pt}#1}{}} % Section text font settings
\fancyhf{} \fancyhead[LE,RO]{\sffamily\normalsize\thepage} % Font setting for the page number in the header
\fancyhead[LO]{\rightmark} % Print the nearest section name on the left side of odd pages
\fancyhead[RE]{\leftmark} % Print the current chapter name on the right side of even pages
\renewcommand{\headrulewidth}{0.5pt} % Width of the rule under the header
\addtolength{\headheight}{2.5pt} % Increase the spacing around the header slightly
\renewcommand{\footrulewidth}{0pt} % Removes the rule in the footer
\fancypagestyle{plain}{\fancyhead{}\renewcommand{\headrulewidth}{0pt}} % Style for when a plain pagestyle is specified

% Removes the header from odd empty pages at the end of chapters
\makeatletter
\renewcommand{\cleardoublepage}{
\clearpage\ifodd\c@page\else
\hbox{}
\vspace*{\fill}
\thispagestyle{empty}
\newpage
\fi}

%----------------------------------------------------------------------------------------
%	THEOREM STYLES
%----------------------------------------------------------------------------------------

\usepackage{amsmath,amsfonts,amssymb,amsthm} % For math equations, theorems, symbols, etc

\newcommand{\intoo}[2]{\mathopen{]}#1\,;#2\mathclose{[}}
\newcommand{\ud}{\mathop{\mathrm{{}d}}\mathopen{}}
\newcommand{\intff}[2]{\mathopen{[}#1\,;#2\mathclose{]}}
\newtheorem{notation}{Notation}[chapter]

\newtheoremstyle{ocrenumbox}% % Theorem style name
{0pt}% Space above
{0pt}% Space below
{\normalfont}% % Body font
{}% Indent amount
{\small\bf\sffamily\color{ocre}}% % Theorem head font
{\;}% Punctuation after theorem head
{0.25em}% Space after theorem head
{\small\sffamily\color{ocre}\thmname{#1}\nobreakspace\thmnumber{\@ifnotempty{#1}{}\@upn{#2}}% Theorem text (e.g. Theorem 2.1)
\thmnote{\nobreakspace\the\thm@notefont\sffamily\bfseries\color{black}---\nobreakspace#3.}} % Optional theorem note
\renewcommand{\qedsymbol}{$\blacksquare$}% Optional qed square

\newtheoremstyle{blacknumexa}% Theorem style Example
{5pt}% Space above
{5pt}% Space below
{\normalfont}% Body font
{} % Indent amount
{\small\bf\sffamily\color{sthlmGreen}}% Theorem head font
{\;}% Punctuation after theorem head
{0.25em}% Space after theorem head
{\nobreakspace\thmname{#1}\nobreakspace\thmnumber{\@ifnotempty{#1}{}\@upn{#2}}% Theorem text (e.g. Theorem 2.1)
\thmnote{\nobreakspace\the\thm@notefont\sffamily\bfseries---\nobreakspace#3.}}% Optional theorem note

\newtheoremstyle{blacknumexe}% Theorem style Exercise
{5pt}% Space above
{5pt}% Space below
{\normalfont}% Body font
{} % Indent amount
{\small\bf\sffamily\color{ocre}}% Theorem head font
{\;}% Punctuation after theorem head
{0.25em}% Space after theorem head
{\nobreakspace\thmname{#1}\nobreakspace\thmnumber{\@ifnotempty{#1}{}\@upn{#2}}% Theorem text (e.g. Theorem 2.1)
\thmnote{\nobreakspace\the\thm@notefont\sffamily\bfseries---\nobreakspace#3.}}% Optional theorem note

\newtheoremstyle{blacknumessential}% Theorem style name
{5pt}% Space above
{5pt}% Space below
{\normalfont}% Body font
{} % Indent amount
{\small\bf\sffamily\color{sthlmRed}}% Theorem head font
{\;}% Punctuation after theorem head
{0.25em}% Space after theorem head
{\nobreakspace\thmname{#1}\nobreakspace\thmnumber{\@ifnotempty{#1}{}\@upn{#2}}% Theorem text (e.g. Theorem 2.1)
\thmnote{\nobreakspace\the\thm@notefont\sffamily\bfseries---\nobreakspace#3.}}% Optional theorem note

\newtheoremstyle{blacknumpro}% Theorem style name properties
{5pt}% Space above
{5pt}% Space below
{\normalfont}% Body font
{} % Indent amount
{\small\bf\sffamily\color{sthlmPurple}}% Theorem head font
{\;}% Punctuation after theorem head
{0.25em}% Space after theorem head
{\nobreakspace\thmname{#1}\nobreakspace\thmnumber{\@ifnotempty{#1}{}\@upn{#2}}% Theorem text (e.g. Theorem 2.1)
\thmnote{\nobreakspace\the\thm@notefont\sffamily\bfseries---\nobreakspace#3.}}% Optional theorem note

\newtheoremstyle{blacknumbox} % Theorem style name
{0pt}% Space above
{0pt}% Space below
{\normalfont}% Body font
{}% Indent amount
{\small\bf\sffamily}% Theorem head font
{\;}% Punctuation after theorem head
{0.25em}% Space after theorem head
{\small\sffamily\thmname{#1}\nobreakspace\thmnumber{\@ifnotempty{#1}{}\@upn{#2}}% Theorem text (e.g. Theorem 2.1)
\thmnote{\nobreakspace\the\thm@notefont\sffamily\bfseries---\nobreakspace#3.}}% Optional theorem note

% Non-boxed/non-framed environments
\newtheoremstyle{ocrenum}% % Theorem style name
{5pt}% Space above
{5pt}% Space below
{\normalfont}% % Body font
{}% Indent amount
{\small\bf\sffamily\color{ocre}}% % Theorem head font
{\;}% Punctuation after theorem head
{0.25em}% Space after theorem head
{\small\sffamily\color{ocre}\thmname{#1}\nobreakspace\thmnumber{\@ifnotempty{#1}{}\@upn{#2}}% Theorem text (e.g. Theorem 2.1)
\thmnote{\nobreakspace\the\thm@notefont\sffamily\bfseries\color{black}---\nobreakspace#3.}} % Optional theorem note
\renewcommand{\qedsymbol}{$\blacksquare$}% Optional qed square
\makeatother

% Defines the theorem text style for each type of theorem to one of the three styles above
\newcounter{dummy} 
\numberwithin{dummy}{section}
\theoremstyle{ocrenumbox}
\newtheorem{theoremeT}[dummy]{Theorem}
\newtheorem{aruleT}[dummy]{Rule}
\newtheorem{definitionT}[dummy]{Definition}
\newtheorem{problem}{Problem}[chapter]
\theoremstyle{blacknumexe}
\newtheorem{exerciseT}{Exercise}[chapter]
\theoremstyle{blacknumexa}
\newtheorem{exampleT}{Example}[chapter]
\theoremstyle{blacknumessential}
\newtheorem{essentialqT}{Essential Questions}[chapter]
\theoremstyle{blacknumpro}
\newtheorem{propertyT}[dummy]{Property}
\theoremstyle{blacknumbox}
\newtheorem{vocabulary}{Vocabulary}[chapter]
%\newtheorem{definitionT}{Definition}[section]
\newtheorem{corollaryT}[dummy]{Corollary}
\theoremstyle{ocrenum}
\newtheorem{proposition}[dummy]{Proposition}

%----------------------------------------------------------------------------------------
%	DEFINITION OF COLORED BOXES
%----------------------------------------------------------------------------------------

\RequirePackage[framemethod=default]{mdframed} % Required for creating the theorem, definition, exercise and corollary boxes

% Theorem box
\newmdenv[skipabove=7pt,
skipbelow=7pt,
backgroundcolor=black!5,
linecolor=ocre,
innerleftmargin=5pt,
innerrightmargin=5pt,
innertopmargin=5pt,
leftmargin=0cm,
rightmargin=0cm,
innerbottommargin=5pt]{tBox}

% Rule box
\newmdenv[skipabove=7pt,
skipbelow=7pt,
backgroundcolor=black!5,
linecolor=ocre,
innerleftmargin=5pt,
innerrightmargin=5pt,
innertopmargin=5pt,
leftmargin=0cm,
rightmargin=0cm,
innerbottommargin=5pt]{arBox}

% Definition box
\newmdenv[skipabove=7pt,
skipbelow=7pt,
backgroundcolor=sthlmDarkGrey!5,
linecolor=ocre,
innerleftmargin=5pt,
innerrightmargin=5pt,
innertopmargin=5pt,
leftmargin=0cm,
rightmargin=0cm,
innerbottommargin=5pt]{dBox}

% Property box
\newmdenv[skipabove=7pt,
skipbelow=7pt,
backgroundcolor=sthlmDarkGrey!5,
linecolor=sthlmPurple,
innerleftmargin=5pt,
innerrightmargin=5pt,
innertopmargin=5pt,
leftmargin=0cm,
rightmargin=0cm,
innerbottommargin=5pt]{prBox}

% Exercise box	  
\newmdenv[skipabove=7pt,
skipbelow=7pt,
rightline=false,
leftline=true,
topline=false,
bottomline=false,
backgroundcolor=ocre!5,
linecolor=ocre,
innerleftmargin=5pt,
innerrightmargin=5pt,
innertopmargin=5pt,
innerbottommargin=5pt,
leftmargin=0cm,
rightmargin=0cm,
linewidth=4pt]{eBox}	

% Definition box
%\newmdenv[skipabove=7pt,
%skipbelow=7pt,
%rightline=false,
%leftline=true,
%topline=false,
%bottomline=false,
%linecolor=ocre,
%innerleftmargin=5pt,
%innerrightmargin=5pt,
%innertopmargin=0pt,
%leftmargin=0cm,
%rightmargin=0cm,
%linewidth=4pt,
%innerbottommargin=0pt]{dBox}

% Example box
\newmdenv[skipabove=7pt,
skipbelow=7pt,
rightline=false,
leftline=true,
topline=false,
bottomline=false,
backgroundcolor=sthlmGreen!5,
linecolor=sthlmGreen,
innerleftmargin=5pt,
innerrightmargin=5pt,
innertopmargin=5pt,
innerbottommargin=5pt,
leftmargin=0cm,
rightmargin=0cm,
linewidth=4pt]{exaBox}

% Essential Question box
\newmdenv[skipabove=7pt,
skipbelow=7pt,
rightline=false,
leftline=false,
topline=true,
bottomline=true,
backgroundcolor=sthlmRed!5,
linecolor=sthlmRed,
innerleftmargin=5pt,
innerrightmargin=5pt,
innertopmargin=5pt,
innerbottommargin=5pt,
leftmargin=0cm,
rightmargin=0cm,
linewidth=4pt]{essentialqBox}			

% Corollary box
\newmdenv[skipabove=7pt,
skipbelow=7pt,
rightline=false,
leftline=true,
topline=false,
bottomline=false,
linecolor=gray,
backgroundcolor=black!5,
innerleftmargin=5pt,
innerrightmargin=5pt,
innertopmargin=5pt,
leftmargin=0cm,
rightmargin=0cm,
linewidth=4pt,
innerbottommargin=5pt]{cBox}

% Creates an environment for each type of theorem and assigns it a theorem text style from the "Theorem Styles" section above and a colored box from above
\newenvironment{theorem}{\begin{tBox}\begin{theoremeT}}{\end{theoremeT}\end{tBox}}
\newenvironment{arule}{\begin{arBox}\begin{aruleT}}{\end{aruleT}\end{arBox}}
\newenvironment{exercise}{\begin{eBox}\begin{exerciseT}}{\hfill{\color{ocre}\tiny\ensuremath{\blacksquare}}\end{exerciseT}\end{eBox}}
\newenvironment{definition}{\begin{dBox}\begin{definitionT}}{\end{definitionT}\end{dBox}}
\newenvironment{property}{\begin{prBox}\begin{propertyT}}{\end{propertyT}\end{prBox}}			  
%\newenvironment{definition}{\begin{dBox}\begin{definitionT}}{\end{definitionT}\end{dBox}}
%\newenvironment{example}{\begin{exBox}\begin{exampleT}}{\end{exampleT}\end{exBox}}	
\newenvironment{example}{\begin{exaBox}\begin{exampleT}}{\hfill{\color{sthlmGreen}\tiny\ensuremath{\blacksquare}}\end{exampleT}\end{exaBox}}
\newenvironment{essentialq}{\begin{essentialqBox}\begin{essentialqT}}{\end{essentialqT}\end{essentialqBox}}
%\newenvironment{example}{\begin{exampleT}}{\hfill{\tiny\ensuremath{\blacksquare}}\end{exampleT}}		
\newenvironment{corollary}{\begin{cBox}\begin{corollaryT}}{\end{corollaryT}\end{cBox}}	

%----------------------------------------------------------------------------------------
%	REMARK ENVIRONMENT
%----------------------------------------------------------------------------------------

\newenvironment{remark}{\par\vspace{10pt}\small % Vertical white space above the remark and smaller font size
\begin{list}{}{
\leftmargin=35pt % Indentation on the left
\rightmargin=25pt}\item\ignorespaces % Indentation on the right
\makebox[-2.5pt]{\begin{tikzpicture}[overlay]
\node[draw=ocre!60,line width=1pt,circle,fill=ocre!25,font=\sffamily\bfseries,inner sep=2pt,outer sep=0pt] at (-15pt,0pt){\textcolor{ocre}{R}};\end{tikzpicture}} % Orange R in a circle
\advance\baselineskip -1pt}{\end{list}\vskip5pt} % Tighter line spacing and white space after remark

\newcommand{\soln}[1]{\par\vspace{10pt}\makebox[15pt]{\begin{tikzpicture}[overlay]
\node[draw=sthlmGreen!70,line width=1pt,circle,fill=sthlmGreen!25,font=\sffamily\bfseries,inner sep=2pt,outer sep=0pt] at (-15pt,0pt){\textcolor{sthlmGreen}{S}};\end{tikzpicture}}\cGreen{\noindent\hrulefill\vskip 15pt}}

%\newcommand{\soln}[1]{\vskip10pt\cGreen{\textbf{Solution:}\hspace{5pt}\noindent\hrulefill \vskip10pt}}

\newcommand{\qdepend}[1]{\par\vspace{10pt}\makebox[15pt]{\begin{tikzpicture}[overlay]
\node[draw=sthlmDarkGrey!70,line width=1pt,circle,fill=sthlmDarkGrey!25,font=\sffamily\bfseries,inner sep=2pt,outer sep=0pt] at (-15pt,0pt){\textcolor{sthlmDarkGrey}{D}};\end{tikzpicture}}\cDarkGrey{\noindent\hrulefill\vskip 15pt}}


%\newenvironment{depend}{\par\vspace{10pt}\small % Vertical white space above the remark and smaller font size
%\begin{list}{}{
%\leftmargin=35pt % Indentation on the left
%\rightmargin=25pt}\item\ignorespaces % Indentation on the right
%\makebox[-2.5pt]{\begin{tikzpicture}[overlay]
%\node[draw=sthlmDarkGrey!60,line width=1pt,circle,fill=sthlmDarkGrey!25,font=\sffamily\bfseries,inner sep=2pt,outer sep=0pt] at (-15pt,0pt){\textcolor{sthlmDarkGrey}{D}};\end{tikzpicture}} % Orange R in a circle
%\advance\baselineskip -1pt}{\end{list}\vskip5pt} % Tighter line spacing and white space after remark


%\newenvironment{soln}{\par\vspace{10pt}\small % Vertical white space above the remark and smaller font size
%\begin{list}{}{
%\leftmargin=35pt % Indentation on the left
%\rightmargin=25pt}\item\ignorespaces % Indentation on the right
%\makebox[-2.5pt]{\begin{tikzpicture}[overlay]
%\node[draw=sthlmGreen!60,line width=1pt,circle,fill=sthlmGreen!25,font=\sffamily\bfseries,inner sep=2pt,outer sep=0pt] at (-15pt,0pt){\textcolor{sthlmGreen}{S}};\end{tikzpicture}}\cGreen{\noindent\hrulefill} % Orange R in a circle
%\advance\baselineskip -1pt}{\end{list}\vskip5pt} % Tighter line spacing and white space after remark

%----------------------------------------------------------------------------------------
%	SECTION NUMBERING IN THE MARGIN
%----------------------------------------------------------------------------------------

\makeatletter
\renewcommand{\@seccntformat}[1]{\llap{\textcolor{ocre}{\csname the#1\endcsname}\hspace{1em}}}                    
\renewcommand{\section}{\@startsection{section}{1}{\z@}
{-4ex \@plus -1ex \@minus -.4ex}
{1ex \@plus.2ex }
{\normalfont\large\sffamily\bfseries}}
\renewcommand{\subsection}{\@startsection {subsection}{2}{\z@}
{-3ex \@plus -0.1ex \@minus -.4ex}
{0.5ex \@plus.2ex }
{\normalfont\sffamily\bfseries}}
\renewcommand{\subsubsection}{\@startsection {subsubsection}{3}{\z@}
{-2ex \@plus -0.1ex \@minus -.2ex}
{.2ex \@plus.2ex }
{\normalfont\small\sffamily\bfseries}}                        
\renewcommand\paragraph{\@startsection{paragraph}{4}{\z@}
{-2ex \@plus-.2ex \@minus .2ex}
{.1ex}
{\normalfont\small\sffamily\bfseries}}

%----------------------------------------------------------------------------------------
%	PART HEADINGS
%----------------------------------------------------------------------------------------

% numbered part in the table of contents
\newcommand{\@mypartnumtocformat}[2]{%
\setlength\fboxsep{0pt}%
\noindent\colorbox{ocre!20}{\strut\parbox[c][.7cm]{\ecart}{\color{ocre!70}\Large\sffamily\bfseries\centering#1}}\hskip\esp\colorbox{ocre!40}{\strut\parbox[c][.7cm]{\linewidth-\ecart-\esp}{\Large\sffamily\centering#2}}}%
%%%%%%%%%%%%%%%%%%%%%%%%%%%%%%%%%%
% unnumbered part in the table of contents
\newcommand{\@myparttocformat}[1]{%
\setlength\fboxsep{0pt}%
\noindent\colorbox{ocre!40}{\strut\parbox[c][.7cm]{\linewidth}{\Large\sffamily\centering#1}}}%
%%%%%%%%%%%%%%%%%%%%%%%%%%%%%%%%%%
\newlength\esp
\setlength\esp{4pt}
\newlength\ecart
\setlength\ecart{1.2cm-\esp}
\newcommand{\thepartimage}{}%
\newcommand{\partimage}[1]{\renewcommand{\thepartimage}{#1}}%
\def\@part[#1]#2{%
\ifnum \c@secnumdepth >-2\relax%
\refstepcounter{part}%
\addcontentsline{toc}{part}{\texorpdfstring{\protect\@mypartnumtocformat{\thepart}{#1}}{\partname~\thepart\ ---\ #1}}
\else%
\addcontentsline{toc}{part}{\texorpdfstring{\protect\@myparttocformat{#1}}{#1}}%
\fi%
\startcontents%
\markboth{}{}%
{\thispagestyle{empty}%
\begin{tikzpicture}[remember picture,overlay]%
\node at (current page.north west){\begin{tikzpicture}[remember picture,overlay]%	
\fill[ocre!20](0cm,0cm) rectangle (\paperwidth,-\paperheight);
\node[anchor=north] at (4cm,-3.25cm){\color{ocre!40}\fontsize{220}{100}\sffamily\bfseries\@Roman\c@part}; 
\node[anchor=south east] at (\paperwidth-1cm,-\paperheight+1cm){\parbox[t][][t]{8.5cm}{
\printcontents{l}{0}{\setcounter{tocdepth}{1}}%
}};
\node[anchor=north east] at (\paperwidth-1.5cm,-3.25cm){\parbox[t][][t]{15cm}{\strut\raggedleft\color{white}\fontsize{30}{30}\sffamily\bfseries#2}};
\end{tikzpicture}};
\end{tikzpicture}}%
\@endpart}
\def\@spart#1{%
\startcontents%
\phantomsection
{\thispagestyle{empty}%
\begin{tikzpicture}[remember picture,overlay]%
\node at (current page.north west){\begin{tikzpicture}[remember picture,overlay]%	
\fill[ocre!20](0cm,0cm) rectangle (\paperwidth,-\paperheight);
\node[anchor=north east] at (\paperwidth-1.5cm,-3.25cm){\parbox[t][][t]{15cm}{\strut\raggedleft\color{white}\fontsize{30}{30}\sffamily\bfseries#1}};
\end{tikzpicture}};
\end{tikzpicture}}
\addcontentsline{toc}{part}{\texorpdfstring{%
\setlength\fboxsep{0pt}%
\noindent\protect\colorbox{ocre!40}{\strut\protect\parbox[c][.7cm]{\linewidth}{\Large\sffamily\protect\centering #1\quad\mbox{}}}}{#1}}%
\@endpart}
\def\@endpart{\vfil\newpage
\if@twoside
\if@openright
\null
\thispagestyle{empty}%
\newpage
\fi
\fi
\if@tempswa
\twocolumn
\fi}

%----------------------------------------------------------------------------------------
%	CHAPTER HEADINGS
%----------------------------------------------------------------------------------------

\newcommand{\thechapterimage}{}%
\newcommand{\chapterimage}[1]{\renewcommand{\thechapterimage}{#1}}%
\def\@makechapterhead#1{%
{\parindent \z@ \raggedright \normalfont
\ifnum \c@secnumdepth >\m@ne
\if@mainmatter
\begin{tikzpicture}[remember picture,overlay]
\node at (current page.north west)
{\begin{tikzpicture}[remember picture,overlay]
\node[anchor=north west,inner sep=0pt] at (0,0) {\includegraphics[width=\paperwidth]{\thechapterimage}};
\draw[anchor=west] (\Gm@lmargin,-9cm) node [line width=2pt,rounded corners=15pt,draw=ocre,fill=white,fill opacity=0.5,inner sep=15pt]{\strut\makebox[22cm]{}};
\draw[anchor=west] (\Gm@lmargin+.3cm,-9cm) node {\huge\sffamily\bfseries\color{black}\thechapter. #1\strut};
\end{tikzpicture}};
\end{tikzpicture}
\else
\begin{tikzpicture}[remember picture,overlay]
\node at (current page.north west)
{\begin{tikzpicture}[remember picture,overlay]
\node[anchor=north west,inner sep=0pt] at (0,0) {\includegraphics[width=\paperwidth]{\thechapterimage}};
\draw[anchor=west] (\Gm@lmargin,-9cm) node [line width=2pt,rounded corners=15pt,draw=ocre,fill=white,fill opacity=0.5,inner sep=15pt]{\strut\makebox[22cm]{}};
\draw[anchor=west] (\Gm@lmargin+.3cm,-9cm) node {\huge\sffamily\bfseries\color{black}#1\strut};
\end{tikzpicture}};
\end{tikzpicture}
\fi\fi\par\vspace*{270\p@}}}

%-------------------------------------------

\def\@makeschapterhead#1{%
\begin{tikzpicture}[remember picture,overlay]
\node at (current page.north west)
{\begin{tikzpicture}[remember picture,overlay]
\node[anchor=north west,inner sep=0pt] at (0,0) {\includegraphics[width=\paperwidth]{\thechapterimage}};
\draw[anchor=west] (\Gm@lmargin,-9cm) node [line width=2pt,rounded corners=15pt,draw=ocre,fill=white,fill opacity=0.5,inner sep=15pt]{\strut\makebox[22cm]{}};
\draw[anchor=west] (\Gm@lmargin+.3cm,-9cm) node {\huge\sffamily\bfseries\color{black}#1\strut};
\end{tikzpicture}};
\end{tikzpicture}
\par\vspace*{270\p@}}
\makeatother

%----------------------------------------------------------------------------------------
%	HYPERLINKS IN THE DOCUMENTS
%----------------------------------------------------------------------------------------

\usepackage{hyperref}
\hypersetup{hidelinks,backref=true,pagebackref=true,hyperindex=true,colorlinks=false,breaklinks=true,urlcolor= ocre,bookmarks=true,bookmarksopen=false,pdftitle={Title},pdfauthor={Author}}
\usepackage{bookmark}
\bookmarksetup{
open,
numbered,
addtohook={%
\ifnum\bookmarkget{level}=0 % chapter
\bookmarksetup{bold}%
\fi
\ifnum\bookmarkget{level}=-1 % part
\bookmarksetup{color=ocre,bold}%
\fi
}
} % Insert the commands.tex file which contains the majority of the structure behind the template

\begin{document}

%-=-=-=-=-=-=-=-=-=-=-=-=-=-=-=-=-=-=-=-=-=-=-=-=
%
%	TITLE PAGE
%
%-=-=-=-=-=-=-=-=-=-=-=-=-=-=-=-=-=-=-=-=-=-=-=-=

\begingroup
\thispagestyle{empty}
\AddToShipoutPicture*{\put(6,5){\includegraphics[scale=1]{Pictures/background}}} % Image background
\centering
\vspace*{9cm}
\par\normalfont\fontsize{35}{35}\sffamily\selectfont
{\textcolor{stockholmGrey}{Mathematics Notebook}} \par % Book title
\vspace*{1cm}
{\Huge {\color{stockholmGrey}{Mark H. Olson}}}\par % Author name
\vspace*{1cm}
{\Large {\color{stockholmGrey}{Updated: \today}}}\par % Author name
\endgroup

%%-=-=-=-=-=-=-=-=-=-=-=-=-=-=-=-=-=-=-=-=-=-=-=-=
%
%	COPYRIGHT PAGE
%
%-=-=-=-=-=-=-=-=-=-=-=-=-=-=-=-=-=-=-=-=-=-=-=-=

\newpage
~\vfill
\thispagestyle{empty}

\noindent Copyright \copyright\ 2014 Mark Hendry Olson\\ % Copyright notice

\noindent \textsc{Published by Mark Hendry Olson}\\ % Publisher

\noindent \textsc{hendryolson.com}\\ % URL

\noindent Licensed under the Creative Commons Attribution-NonCommercial 3.0 Unported License (the ``License''). You may not use this file except in compliance with the License. You may obtain a copy of the License at \url{http://creativecommons.org/licenses/by-nc/3.0}. Unless required by applicable law or agreed to in writing, software distributed under the License is distributed on an \textsc{``as is'' basis, without warranties or conditions of any kind}, either express or implied. See the License for the specific language governing permissions and limitations under the License.\\ % License information

\noindent \textit{First printing, March 2013} % Printing/edition date

%-=-=-=-=-=-=-=-=-=-=-=-=-=-=-=-=-=-=-=-=-=-=-=-=
%
%	TABLE OF CONTENTS
%
%-=-=-=-=-=-=-=-=-=-=-=-=-=-=-=-=-=-=-=-=-=-=-=-=

\chapterimage{chapter_head_1.pdf} % Table of contents heading image

\pagestyle{empty} % No headers

\tableofcontents % Print the table of contents itself

\cleardoublepage % Forces the first chapter to start on an odd page so it's on the right

\pagestyle{fancy} % Print headers again

\begin{sagesilent}
var('x')
#t=ZZ.random_element(999999)
t=2015
set_random_seed(t)

# =-=-=-=-=-=-=-=-=-=-=-=-=-=-=
# Question 20141120-211843
# =-=-=-=-=-=-=-=-=-=-=-=-=-=-=

a_20141120_211843=[(ZZ.random_element(2,14)) for i in range(13)]
b_20141120_211843=[(ZZ.random_element(2,14)) for i in range(13)]

for j in range(len(a_20141120_211843)):
  while a_20141120_211843[j]==a_20141120_211843[(j-1)]:
    a_20141120_211843[j]=ZZ.random_element(2,14)

# =-=-=-=-=-=-=-=-=-=-=-=-=-=-=
# Question 20141121-195608
# =-=-=-=-=-=-=-=-=-=-=-=-=-=-=
	
a_20141121_195608=[(ZZ.random_element(2,14)) for i in range(8)]
b_20141121_195608=[(ZZ.random_element(2,14)) for i in range(8)]

for j in range(len(a_20141121_195608)):
  while b_20141121_195608[j]==b_20141121_195608[(j-1)]:
    b_20141121_195608[j]=ZZ.random_element(2,14)
	
for k_20141121_195608 in range(len(a_20141121_195608)):
  while a_20141121_195608[k_20141121_195608]==b_20141121_195608[(k_20141121_195608-1)]:
    a_20141121_195608[k_20141121_195608]=ZZ.random_element(2,14)	

# =-=-=-=-=-=-=-=-=-=-=-=-=-=-=
# Question 20141121-191224
# =-=-=-=-=-=-=-=-=-=-=-=-=-=-=

a_20141121_191224=[(ZZ.random_element(2,14)) for i in range(13)]
b_20141121_191224=[(ZZ.random_element(2,14)) for i in range(13)]

for j in range(len(a_20141121_191224)):
  while a_20141121_191224[j]==a_20141121_191224[(j-1)]:
    a_20141121_191224[j]=ZZ.random_element(2,14)
	
# =-=-=-=-=-=-=-=-=-=-=-=-=-=-=
# Question 20141123-091420
# =-=-=-=-=-=-=-=-=-=-=-=-=-=-=

a_20141123_091420=[(-1)^(ZZ.random_element(1,3))*(ZZ.random_element(2,14)) for i in range(13)]
b_20141123_091420=[(-1)^(ZZ.random_element(1,3))*(ZZ.random_element(2,14)) for i in range(13)]
c_20141123_091420=[(-1)^(ZZ.random_element(1,3))*(ZZ.random_element(2,14)) for i in range(13)]

for j in range(len(a_20141123_091420)):
  while a_20141123_091420[j]==a_20141123_091420[(j-1)]:
    a_20141123_091420[j]=(-1)^(ZZ.random_element(1,3))*(ZZ.random_element(2,14))
\end{sagesilent}
%%-=-=-=-=-=-=-=-=-=-=-=-=-=-=-=-=-=-=-=-=-=-=-=-=
%
%	CHAPTER 
%
%-=-=-=-=-=-=-=-=-=-=-=-=-=-=-=-=-=-=-=-=-=-=-=-=

\chapterimage{chapter_head_2.pdf} % Chapter heading image

\chapter{Arithmetic Expressions}
%!TEX root = /Users/markholson/Dropbox/+Projects/LatexFiles/Book/main.tex
%-=-=-=-=-=-=-=-=-=-=-=-=-=-=-=-=-=-=-=-=-=-=-=-=
%
%	CHAPTER 
%
%-=-=-=-=-=-=-=-=-=-=-=-=-=-=-=-=-=-=-=-=-=-=-=-=

\chapterimage{chapter_head_2.pdf} % Chapter heading image

\chapter{Algebraic Expressions}

%-=-=-=-=-=-=-=-=-=-=-=-=-=-=-=-=-=-=-=-=-=-=-=-=
%	SECTION: OPERATION OF ADDITION
%-=-=-=-=-=-=-=-=-=-=-=-=-=-=-=-=-=-=-=-=-=-=-=-=
\section{Expressions}\index{Algebraic Expressions}

\begin{essentialq}\hfill \\

\begin{enumerate}
	\item What is an algebraic expression?
\end{enumerate}

\end{essentialq}


\begin{table}
\begin{tabular}{|l|c|c|c|c|c|c}
\hline
 									& Arithmetic 	& Polynomial	& Algebraic  	\\
\hline
Constant 							& \cellyes		& \cellyes 		& \cellyes		\\
\hline
Factorial 							& \cellyes 		& \cellyes 		& \cellyes		\\
\hline
Variable: parameter/coefficient 	& \cellyes 		& \cellyes 		& \cellyes		\\
\hline
Variable: unknown/indeterminate 	& \cellno 		& \cellyes 		& \cellyes		\\
\hline
Power with $\integer^{+}$ exponent 	& \cellno 		& \cellyes 		& \cellyes		 \\
\hline
Power with $\integer$ exponent 		& \cellno 		& \cellno 		& \cellyes		 \\
\hline
$n$-th root 						& \cellno 		& \cellno 		& \cellyes		\\
\hline
Power with $\mathbb{Q}$ exponent 	& \cellno 		& \cellno 		& \cellyes		 \\
\hline
\end{tabular}
\caption{Names of different types of expressions}\label{tab:tableofexpressions}
\end{table}


%-=-=-= DEFINITION
%\begin{definition}[Algebraic Expression]\index{Algebraic Expression}

%An algebraic expression is made up of one or more terms (operation of addition), which are themselves made up of factors (operation of multiplication).\\

%Algebraic expressions have three different types of multiplicative factors:\\

%\begin{enumerate}   
%	\item constant (coefficient) 
%	\begin{itemize}
%		\item real number 
%		\item parameter 
%	\end{itemize}
%	\item variable
%	\item rational power
%\end{enumerate}
%\hfill \cite{mathworld:algebraicexpression}
%\end{definition}

\section{Polynomial Expressions}\index{Polynomial Expressions}

%-=-=-= DEFINITION
\begin{definition}[Operation of Addition (OOA)]\index{Operation!Operation of Addition}
\begin{align}
\underbrace{\underbrace{a}_{\text{Augend}}+\underbrace{b}_{\text{Addend}}}_{\text{Sum}} \label{eq:ooa}
\end{align}
\end{definition}


%-=-=-= DEFINITION
\begin{definition}[Common Denominator (CD)]\index{Common Denominator}
\begin{subequations}
\begin{align}
\dfrac{a}{b} + \dfrac{c}{b} &= \dfrac{a+c}{b} \label{eq:cd1} \\
\dfrac{a+c}{b}&= \dfrac{a}{b} + \dfrac{c}{b} \label{eq:cd2}
\end{align}
\end{subequations}
\end{definition}

%-=-=-= RULE
\begin{arule}[Fraction Operation of Addition (FOOA)]\index{Fraction Operation of Addition}
\begin{subequations}
\begin{align}
\dfrac{a}{b} + \dfrac{c}{d} &= \dfrac{ad+bc}{bd} \label{eq:fooa1} \\
\dfrac{ad+bc}{bd} &= \dfrac{a}{b} + \dfrac{c}{d} \label{eq:fooa2}
\end{align}
\end{subequations}
\end{arule}

\newpage
\begin{figure}
\begin{tikzpicture}[scale=1, node distance=1.1cm, text width=5em, auto]
    % Place nodes
\node [property] (MId1) {MId};
\node [notation, below of=MId1] (ONeg1) {ONeg};
\node [operation, below of=ONeg1] (DOS1) {DOS};
\node [delim, below of=DOS1] (DELIM) {DELIM};
\node [property, below of=DELIM] (DPE) {DPE};
\node [notation, below of=DPE] (JTC) {JTC};
\node [property, below of=JTC] (CPM) {CPM};
\node [operation, below of=CPM] (OOM) {OOM};
\node [property, right of=OOM, node distance=2.8cm] (APM) {APM};
\node [algorithm, below of=OOM] (RF1) {RF};
\node [notation, below of=RF1] (CTJ) {CTJ};
\node [notation, below of=CTJ] (CPA) {CPA};
\node [property, below of=CPA] (DPF) {DPF};
\node [operation, below of=DPF] (OOA) {OOA};
\node [property, right of=OOA,node distance=2.8cm] (APA) {APA};
\node [algorithm, below of=OOA] (RF2) {RF};
\node [property, below of=RF2] (AId) {AId};
\node [operation, below of=AId] (DOS2) {DOS};
\node [notation, below of=DOS2] (ONeg2) {ONeg};
\node [property, below of=ONeg2] (MId2) {MId};

\path [line] (MId1) -- (ONeg1);
\path [line] (ONeg1) -- (DOS1);
\path [line] (DOS1) -- (DELIM);
\path [line] (DELIM) -- (DPE);
\path [line] (DPE) -- (JTC);
\path [line] (JTC) -- (CPM);
\path [line] (CPM) edge [out= 315, in= 150] (APM);
\path [line] (APM) -- (OOM);
\path [line] (OOM) -- (RF1);
\path [line] (RF1) -- (CTJ);
\path [line] (CTJ) -- (CPA);
\path [line] (CPA) -- (DPF);
\path [line] (DPF) edge [out= 315, in= 150] (APA);
\path [line] (APA) -- (OOA);
\path [line] (OOA) -- (RF2);
\path [line] (RF2) -- (AId);
\path [line] (AId) -- (DOS2);
\path [line] (DOS2) -- (ONeg2);
\path [line] (ONeg2) -- (MId2);

\node [function, right of=DELIM, node distance=8cm] (PoTF) {PoTF};
\node [function, below of=PoTF] (RTPo) {RTPo};
\node [function, below of=RTPo] (PoNE1) {PoNE};
\node [function, below of=PoNE1] (PoPr) {PoPr};
\node [function, below of=PoPr] (PoQ) {PoQ};
\node [function, below of=PoQ] (PoPrPo) {PoPrPo};
\node [function, below of=PoPrPo] (PoQPo) {PoQPo};
\node [function, below of=PoQPo] (PrCBPo) {PrCBPo};
\node [function, below of=PrCBPo] (QCBPo) {QCBPo};
\node [function, below of=QCBPo] (PoPo) {PoPo};
\node [function, below of=PoPo] (PoNE2) {PoNE};
\node [operation, below of=PoNE2] (OOE) {OOE};
\node [function, below of=OOE] (PoTR) {PoTR};


\path [line] (DELIM) edge [bend left=30] (PoTF);
\path [line] (PoTF) -- (RTPo);
\path [line] (RTPo) -- (PoNE1);
\path [line] (PoPr) -- (PoQ);
\path [line] (PoQ) -- (PoPrPo);
\path [line] (PoPrPo) -- (PoQPo);
\path [line] (PoQPo) -- (PrCBPo);
\path [line] (PrCBPo) -- (QCBPo);
\path [line] (QCBPo) -- (PoPo);
\path [line] (PoPo) -- (PoNE2);
\path [line] (PoNE2) -- (OOE);
\path [line] (OOE) -- (PoTR);

%\node[productdotted, fit =(DELIM)(JTC)(CPM)(OOM)(RF1)(CTJ)(APM)](product) {};
%\node[left of =product, text centered,](powers) {\textcolor{sthlmPurple}{\textbf{Multiplication}}};
%\node[sumdotted, fit =(CPA)(DPF)(OOA)(RF2)(AId)(ONeg2)(DOS2)(MId2)(APA)](sum2) {};
%\node[sumdotted, fit =(MId1)(ONeg1)(DOS1)](sum1) {};
%\node[left of =sum1, text centered,](powers) {\textcolor{sthlmPurple}{\textbf{Addition}}};
%\node[left of =sum2, text centered,](powers) {\textcolor{sthlmPurple}{\textbf{Addition}}};


\node[functiondotted, fit =(PoTF)(RTPo)(PoNE1)(PoPr)(PoQ)(PoPrPo)(PoQPo)(PrCBPo)(QCBPo)(PoPo)(PoNE2)(OOE)(PoTR)](powerfunctions) {};
\node[below of =PoTR, text centered,](powers) {\textcolor{sthlmPurple}{\textbf{Powers}}};

\path [line] (PoNE1) edge [out= 180, in= 0] (DELIM);

\path [line] (CPM) edge [out= 0, in= 180] (PoPr);

\path [line] (PoTR) edge [out= 180, in= 0] (APM);

\path [line] (AId) edge [out= 180, in= 180] (DELIM);

\end{tikzpicture}
\caption{Simplifying Expressions Workflow: \newline {\color{stockholmPink!40}\ensuremath{\blacksquare}} Property, {\color{stockholmBlue!40}\ensuremath{\blacksquare}} Operation, {\color{stockholmGreen!40}\ensuremath{\blacksquare}} Notation, {\color{stockholmPurple!40}\ensuremath{\blacksquare}} Powers, {\color{stockholmYellow!40}\ensuremath{\blacksquare}} Delimiters, {\color{stockholmOrange!40}\ensuremath{\blacksquare}} Process, {\color{stockholmDarkGrey!40}\ensuremath{\blacksquare}} Not Used}
\end{figure}


%-=-=-= DEFINITION
\begin{definition}[Operation of Multiplication (OOM)] \index{Operation!Operation of Multiplication}
\begin{align}
\underbrace{\underbrace{a}_{\text{Multiplicand}} \cdot \underbrace{b}_{\text{Multiplier}}}_{\text{Product}} \label{eq:oom}
\end{align}
\end{definition}



%-=-=-= DEFINITION
\begin{definition}[Operation of Exponentiation (OOE)]\index{Operation!Operation of Exponentiation}
\begin{align}
\underbrace{{\underbrace{b}_{base}}^{\overbrace{m}^{Exponent}}}_{Power} \label{eq:ooe}
\end{align}
\end{definition}


%-=-=-=-=-=-=-=-=-=-=-=-=-=-=-=-=-=-=-=-=-=-=-=-=
%	SECTION: ALGEBRAIC EXPRESSIONS
%-=-=-=-=-=-=-=-=-=-=-=-=-=-=-=-=-=-=-=-=-=-=-=-=






%-=-=-= DEFINITION
\begin{definition}[Juxtaposition to Center-Dot (JTC)]\index{Notation!Juxtapostion to Center-Dot}
\begin{align}
ab &= a \cdot b \label{eq:jtc}
\end{align}
\end{definition}

%-=-=-= DEFINITION
\begin{definition}[Center-Dot to Justapostion (CTJ)]\index{Notation!Center-Dot to Juxtaposition}
\begin{align}
a \cdot b &= ab \label{eq:ctj}
\end{align}
\end{definition}


%-=-=-= DEFINITION
\begin{definition}[Commutative Property of Multiplication (CPM)]\index{Property!Commutative Property of Multiplication}
\begin{align}
\alert{a} \cdot b &= b \cdot \alert{a} \label{eq:cpm}
\end{align}
\end{definition}



%-=-=-= DEFINITION
\begin{definition}[Multiplicative Inverse (MI)]\index{Multiplicative Inverse}
\begin{subequations}
\begin{align}
a \cdot \alert{\dfrac{1}{a}} &= 1 \label{eq:mi1} \\
a \cdot \alert{a^{-1}} &= 1 \label{eq:mi2} 
\end{align}
\end{subequations}
\end{definition}

%-=-=-= DEFINITION
\begin{definition}[Associative Property of Multiplication (APM)]\index{Property!Associative Property of Multiplication}
\begin{subequations}
\begin{align}
a\cdot b\cdot c &= (a\cdot b)\cdot c \label{eq:apm1} \\
a\cdot b\cdot c &= a\cdot (b\cdot c) \label{eq:apm2}
\end{align}
\end{subequations}
\end{definition}


%-=-=-=-=-=-=-=-=-=-=-=-=-=-=-=-=-=-=-=-=-=-=-=-=
%	SECTION: Powers
%-=-=-=-=-=-=-=-=-=-=-=-=-=-=-=-=-=-=-=-=-=-=-=-=

\subsubsection{Powers}\index{Powers}


\begin{arule}[Power of a Quotient of Powers (PoQPo)]\index{Powers!Power of a Quotient of Powers}
\begin{subequations}
\begin{align}
	\left(\dfrac{a^m}{b^n}\right)^k &= \dfrac{a^{m \cdot k}}{b^{n\cdot k}} \label{eq:poqpo1}\\
	\dfrac{a^{m \cdot k}}{b^{n\cdot k}} &= \left(\dfrac{a^m}{b^n}\right)^k \label{eq:poqpo2}
\end{align}
\end{subequations}
\end{arule}

\begin{arule}[Power of a Product of Powers (PoPrPo)]\index{Powers!Power of a Product of Powers}
\begin{subequations}
\begin{align}
	\left(a^m \cdot b^n\right)^k &= a^{m \cdot k} \cdot b^{n\cdot k} \label{eq:poprpo1}\\
	a^{m \cdot k} \cdot b^{n \cdot k} &= \left(a^m \cdot b^n\right)^k \label{eq:poprpo2}
\end{align}
\end{subequations}
\end{arule}

\begin{definition}[Power To Factor (PoTF)]\index{Powers!Power to Factor}
\begin{align}
a^{\alert{n}} &= a_1 \cdot a_2 \cdot \ldots \cdot a_{n-1} \cdot a_{\alert{n}} \label{eq:potf} 
\end{align}
\end{definition}

\begin{definition}[Factor To Power (FTPo)]\index{Powers!Factor to Power}
\begin{align}
a_1 \cdot a_2 \cdot \ldots \cdot a_{n-1} \cdot a_{\alert{n}} &= a^{\alert{n}} \label{eq:ftpo} 
\end{align}
\end{definition}

%-=-=-= DEFINITION
\begin{definition}[Power Inverse (PoI)]\index{Power Inverse}
\begin{subequations}
\begin{align}
\left(b^{m}\right)^{\frac{1}{m}} &= b  \label{eq:poi} 
\end{align}
\end{subequations}
\end{definition}

%-=-=-= DEFINITION
\begin{definition}[Power Inverse (PoId)]\index{Power Identity}
\begin{subequations}
\begin{align}
1&= b^0  \label{eq:poid1} \\
b^{0}&= 1  \label{eq:poid2} 
\end{align}
\end{subequations}
\end{definition}

\begin{notation}[Radical To Power (RTPo)]\index{Powers!Radical to Power}
\begin{align}
\sqrt[m]{b^n} &= b^{\frac{n}{m}}	 \label{eq:rtpo} 
\end{align}
\end{notation}

\begin{notation}[Power To Radical (PoTR)]\index{Powers!Power to Radical}
\begin{align}
b^{\frac{n}{m}} &= \sqrt[m]{b^n}	 \label{eq:potr} 
\end{align}
\end{notation}




%-=-=-=-=-=-=-=-=-=-=-=-=-=-=-=-=-=-=-=-=-=-=-=-=
%	SUBSECTION: Monomials of Like Terms
%-=-=-=-=-=-=-=-=-=-=-=-=-=-=-=-=-=-=-=-=-=-=-=-=

\subsection{Monomials of Like Terms}\index{Monomials of Like Terms}








%-=-=-= DEFINITION
\begin{definition}[Additive Inverse (AI)]\index{Additive Inverse}
\begin{subequations}
\begin{align}
a + \alert{(-a)} &= 0 \label{eq:ai}
\end{align}
\end{subequations}
\end{definition}




\subsection{Surds}\index{Surds}

%-=-=-= EXAMPLE
\begin{example}[id:20141108-085327] \label{20141108-085327} \index{Example!20141108-085327} \hfill \\

Simplify $2\sqrt{2}-\dfrac{\left(\sqrt{2}\right)^3}{3}-\left(2\left(-\sqrt{2}\right)-\dfrac{\left(-\sqrt{2}\right)^3}{3} \right)$

\soln

\solnsteps
\begin{align*}
2 \sqrt{2}- \dfrac{\alert{1}\left(\alert{1}\sqrt{2}\right)^{3}}{3}-\alert{1} \left(2 \left(-\alert{1}\sqrt{2} \right) - \dfrac{\alert{1} \left(-\alert{1} \sqrt{2}\right)^{3}}{3} \right) && \text{MId} \eqref{eq:mid1} \\
2 \sqrt{2} - \dfrac{1 \left(1\sqrt{2}\right)^{3}}{3}-1 \left(2 \left(\alert{\neg 1} \sqrt{2} \right) - \dfrac{1 \left(\alert{\neg 1} \sqrt{2}\right)^{3}}{3} \right) && \text{ONeg} \eqref{eq:oneg1} \\
2 \sqrt{2} + \dfrac{\neg 1 \left(1\sqrt{2}\right)^{3}}{3}+ \neg 1 \left(2 \left(\alert{\neg 1} \sqrt{2} \right) + \dfrac{\neg 1 \left(\alert{\neg 1} \sqrt{2}\right)^{3}}{3} \right) && \text{DOS} \eqref{eq:dos1} \\
2 \cdot \alert{2^{1/2}} + \dfrac{\neg 1 \left(1 \cdot \alert{2^{1/2}}\right)^{3}}{3} + \neg 1 \left(2 \left(\neg  1 \cdot \alert{2^{1/2}} \right) + \dfrac{\neg 1 \left(\neg  1  \cdot \alert{2^{1/2}}\right)^{3}}{3} \right) && \text{RTPo} \eqref{eq:rtpo} \\
2 \cdot 2^{1/2} + \dfrac{\neg 1 \left(1 \cdot 2^{1/2}\right)^{3}}{3} + \neg 1 \left(\alert{2 \cdot \neg  1 \cdot 2^{1/2}} + \dfrac{\neg 1 \left(\neg 1  \cdot 2^{1/2}\right)^{3}}{3} \right) && \text{JTC} \eqref{eq:jtc} \\
2 \cdot 2^{1/2} + \dfrac{\neg 1 \cdot \alert{1 \cdot 2^{3/2}}}{3} + \neg 1 \left(2 \cdot \neg  1 \cdot 2^{1/2} + \dfrac{\neg 1 \cdot \alert{\neg 1  \cdot 2^{3/2}}}{3} \right) && \text{PoPrPo} \eqref{eq:poprpo1} \\
2 \cdot 2^{1/2} + \dfrac{\neg 1 \cdot 1 \cdot \alert{2^{2/2} \cdot 2^{1/2}}}{3} + \neg 1 \left(2 \cdot \neg  1 \cdot 2^{1/2} + \dfrac{\neg 1 \cdot \neg 1  \cdot \alert{2^{2/2} \cdot 2^{1/2}}}{3} \right) && \text{PrCBPo} \eqref{eq:prcbpo2} \\
2 \cdot 2^{1/2} + \dfrac{\neg 1 \cdot 1 \cdot \alert{2} \cdot 2^{1/2}}{3} + \neg 1 \left(2 \cdot \neg  1 \cdot 2^{1/2} + \dfrac{\neg 1 \cdot \neg 1  \cdot \alert{2} \cdot 2^{1/2}}{3} \right) && \text{MId} \eqref{eq:mid2} \\
2 \cdot \alert{\sqrt{2}} + \dfrac{\neg 1 \cdot 1 \cdot 2 \cdot \alert{\sqrt{2}}}{3} + \neg 1 \left(2 \cdot \neg  1 \cdot \alert{\sqrt{2}} + \dfrac{\neg 1 \cdot \neg 1  \cdot 2 \cdot \alert{\sqrt{2}}}{3} \right) && \text{PoTR} \eqref{eq:potr} \\
2 \cdot \sqrt{2} + \dfrac{\neg 1 \cdot 1 \cdot 2 \cdot \sqrt{2}}{3} + \alert{\neg 1} \cdot 2 \cdot \neg  1 \cdot \sqrt{2} + \dfrac{\alert{\neg 1} \cdot \neg 1 \cdot \neg 1  \cdot 2 \cdot \sqrt{2}}{3}  && \text{DPE} \eqref{eq:dpe1} \\
2 \cdot \sqrt{2} + \dfrac{\neg 2 \cdot \sqrt{2}}{3} + 2 \cdot \sqrt{2} + \dfrac{\neg 2 \cdot \sqrt{2}}{3}  && \text{OOM} \eqref{eq:oom} \\
2\sqrt{2} + \dfrac{\neg 2\sqrt{2}}{3} + 2\sqrt{2} + \dfrac{\neg 2\sqrt{2}}{3} && \text{CTJ} \eqref{eq:ctj} \\
\left(2 + \dfrac{\neg 2}{3} + 2 + \dfrac{\neg 2}{3} \right)  \sqrt{2} && \text{DPF} \eqref{eq:dpf1} \\
\dfrac{8}{3}\sqrt{2} && \text{OOA} \eqref{eq:ooa} 
\end{align*}

\qdepend 

\qdependlist
example \ref{20141108-083108}-20141108-083108

\end{example}




%!TEX root = /Users/markholson/Dropbox/+Projects/LatexFiles/Book/main.tex
%-=-=-=-=-=-=-=-=-=-=-=-=-=-=-=-=-=-=-=-=-=-=-=-=
%
%	CHAPTER 
%
%-=-=-=-=-=-=-=-=-=-=-=-=-=-=-=-=-=-=-=-=-=-=-=-=

\chapterimage{chapter_head_2.pdf} % Chapter heading image

\chapter{Univariate Polynomial Expressions}

%-=-=-=-=-=-=-=-=-=-=-=-=-=-=-=-=-=-=-=-=-=-=-=-=
%	SECTION: OPERATION OF ADDITION
%-=-=-=-=-=-=-=-=-=-=-=-=-=-=-=-=-=-=-=-=-=-=-=-=
\section{Classification of Univariate Polynomial Expressions}\index{Classification of Univariate Polynomial Expressions}

%-=-=-= DEFINITION
\begin{definition}[Indeterminate]\index{Indeterminate}

\[
\cRed{x}
\]	
	
An indeterminate is a symbol that is treated as a variable, but does not stand for anything else but itself and is used as a placeholder. 

\begin{itemize}
	\item it does \textbf{not} designate a constant or a parameter
	\item it is \textbf{not} an unknown that could be solved for
	\item it is \textbf{not} a variable designating a function argument
\end{itemize}

\hfill \cite{wikipediate:indeterminate}
\end{definition}


%-=-=-= DEFINITION
\begin{definition}[Coefficient]\index{Coefficient}
\[	
	\cBlue{C}x^{k}
\]
A coefficient, $\cBlue{C}$ is a real number multiplicative factor.\\

\end{definition}

%-=-=-= DEFINITION
\begin{definition}[Univariate Monomial]\index{Univariate Monomial}

\[
 \cBlue{C_{k}}\cRed{x^{k}}
\]	
	
A univariate monomial is made up of two factors.  The first factor of a monomial, $\cBlue{C_{k}}$, is the \cBlue{coefficient}. The second factor of each monomial, $\cRed{x^{k}}$, is an indeterminate raised to a non-negative integer power $k$. \\
 
\end{definition}

%-=-=-= DEFINITION
\begin{definition}[Multiplicative Identity (MId)]\index{Multiplicative Identity}
\begin{subequations}
\begin{align}
\alert{1}a &= a \label{eq:mid1} \\
a &= \alert{1}a \label{eq:mid2} 
\end{align}
\end{subequations}
\end{definition}

\begin{remark}
If the coefficient of a univariate monomial is the multiplicative identity \ref{eq:mid1}, 1, then it is not shown in it's canonical form.
\begin{align*}
 \cBlue{C_{k}}\cRed{x^{k}}	& = \cBlue{C_{k}} \cRed{x^{k}} \\
 							& = \cBlue{1} \cRed{x^{k}} \\
                            & = \cRed{x^{k}}
\end{align*}
\end{remark}

%-=-=-= EXAMPLE
\begin{example}[id:20141120-202842] \label{20141120-202842}\index{Example!20141120-202842} \hfill \\

Express $1x^2$ in canonical form.

\soln

\solnsteps
\begin{align*}
x^2 && \text{MId} \eqref{eq:mid2} 
\end{align*}
\end{example}

%-=-=-= DEFINITION
\begin{definition}[Univariate Polynomial Expression]\index{Univariate Polynomial Expression}
\begin{align}
\sum\limits_{k=0}^n \cBlue{C_{k}}\cRed{x^{n}} & = \cBlue{C_{n}}\cRed{x^{n}}+\cBlue{C_{n-1}}\cRed{x^{n-1}}+ \cdots + \cBlue{C_{k}}\cRed{x^{k}} + \cdots + \cBlue{C_{2}}\cRed{x^{2}}+\cBlue{C_{1}}\cRed{x^{1}}+\underbrace{\cBlue{C_{0}}}_{\cBlue{C_{0}}\cRed{x^{0}}} 
\end{align}
A univariate polynomial in an indeterminate $x$ is an expression made up of one or more summands of the form $\cBlue{C_{k}}\cRed{x^{k}}$, which are called monomials.  The first factor of each monomial, $\cBlue{C_{k}}$, is a numerical factor called the \cBlue{coefficient} where $\cBlue{C_{k}} \in \real$. The second factor of each monomial, $\cRed{x^{k}}$, is an indeterminate raised to a non-negative integer power $i$.  

\hfill \cite{wikipedia:polynomial}
\end{definition}

%-=-=-= DEFINITION
\begin{definition}[Degree of the Indeterminate]\index{Degree of the Indeterminate}

\[
x^{\cRed{k}}
\]
The exponent of an indeterminate power, $\cRed{k}$ is called the degree of the indeterminate.  

\hfill \cite{wikipedia:polynomial}
\end{definition}

%-=-=-= DEFINITION
\begin{definition}[Degree of the Univarite Polynomial]\index{Degree of the Univariate Polynomial}

\[
C_{n}x^{\cRed{n}}+C_{n-1}x^{n-1}+ \cdots + C_{k}x^{k} + \cdots + C_{2}x^{2}+C_{1}x^{1}+\underbrace{C_{0}}_{C_{0}x^{0}} 
\]

The degree of the univariate polynomial is determined by the monomial with the largest degree of the indeterminate.
 
%\hfill \cite{mathworld:algebraicexpression}
\end{definition}

%-=-=-=-=-=-=-=-=-=-=-=-=-=-=-=-=-=-=-=-=-=-=-=-=
%	SECTION: 
%-=-=-=-=-=-=-=-=-=-=-=-=-=-=-=-=-=-=-=-=-=-=-=-=
\section{Degree -1 Univariate Polynomials}\index{Degree -1 Univariate Polynomials}

%-=-=-=-=-=-=-=-=-=-=-=-=-=-=-=-=
%	SUBSECTION: 
%-=-=-=-=-=-=-=-=-=-=-=-=-=-=-=-=
\subsection*{Monomials}

%-=-=-=-=-=-=-=-=-=-=-=-=-=-=-=-=-=-=-=-=-=-=-=-=
%	SECTION: 
%-=-=-=-=-=-=-=-=-=-=-=-=-=-=-=-=-=-=-=-=-=-=-=-=
\section{Degree 0 Univariate Polynomials}\index{Degree 0 Univariate Polynomials}

% TODO Define the PowerID

\[
\cDarkGrey{C_{n}x^{n}+C_{n-1}x^{n-1}+ \cdots + C_{k}x^{k} + \cdots + C_{2}x^{2}+C_{1}x^{1}+}\underbrace{\cBlue{C_{0}}}_{\cBlue{C_{0}}\cRed{x^{0}}} 
\]

Degree 0 univariate polynomial expressions are made up of univariate monomials, $\cBlue{C_{0}}$, called \alert{constants}.  The power identity is an indeterminate raised to a power of 0 has a value of 1.  Thus, $\cRed{x^{0}}=1$ and results in the monomial $\cBlue{C_{0}} \cdot 1$.  The canonical form of a product does not show the multiplicative identity factor, so what remains of this monomial product is only the coefficient factor $\cBlue{C_{0}}$ and from now on will be referred to as a \alert{constant}.  

%-=-=-= EXAMPLE
\begin{example}[id:20141121-093747] \label{20141121-093747}\index{Example!20141121-093747} \hfill \\

Express $13x^0$ in canonical form.

\soln

\solnsteps
\begin{align*}
13 % TODO Add PowerID here.
\end{align*}
\end{example}  

%-=-=-=-=-=-=-=-=-=-=-=-=-=-=-=-=
%	SUBSECTION: 
%-=-=-=-=-=-=-=-=-=-=-=-=-=-=-=-=
\subsection*{Monomials}

Degree 0 univariate polynomial expressions are usually a monomial in their canonical form if $\cBlue{C_{0}}$ is a non-zero real number.  The exception is if $\cBlue{C_{0}}=0$, the additive identity, then the result is the zero polynomial, which can be considered a degree -1 polynomial.\\   

The expression can be manipulated into its monomial canonical form by simplifying the expression.  Simplifying the expression can be defined as evaluating the expression by following order of operations, which is the same as evaluating an arithmetic expression.  

% TODO Add reference to evaluating arithmetic expressions.

%-=-=-=-=-=-=-=-=-=-=-=-=-=-=-=-=-=-=-=-=-=-=-=-=
%	SECTION: 
%-=-=-=-=-=-=-=-=-=-=-=-=-=-=-=-=-=-=-=-=-=-=-=-=
\section{Degree 1 Univariate Polynomials}\index{Degree 1 Univariate Polynomials}

\[
\cDarkGrey{C_{n}x^{n}+C_{n-1}x^{n-1}+ \cdots + C_{k}x^{k} + \cdots + C_{2}x^{2}}+\cRed{C_{1}x^{1}+\underbrace{C_{0}}_{C_{0}x^{0}}} 
\]

Degree 1 univariate polynomial expressions can be expressed with at most two different terms and consequently this expression in its canonical form has at most two monomial summands -- called a binomial.

%-=-=-= DEFINITION
\begin{definition}[Univariate Like Terms]\index{Like Terms}

\[
\cBlue{C_{1}}\cRed{x^{k}} = \cBlue{C_{2}}\cRed{x^{k}} 
\]
Two or more univariate monomials are defined as having like terms if each monomial has the same term, which will be the same indeterminate raised to the same positive integer power.
\end{definition}

Sometimes the word \alert{term} is used to describe monomials (including both the coefficient and the term), which may be confusing when trying to define like terms.  For this reason, we will refer to the summands of a polynomial as monomials. 

The monomials $5x^1$ and $3x^1$ can be described as having like terms because they share the common term $x^1$.  One could also say that $5x^1$ and $3x^1$ are like terms by definition and consequenlty giving the reader an impression that $5x^1$ and $3x^1$ are terms themselves.

\begin{remark}
A degree 1 indeterminate does not display the multiplicative identity in the exponent when its in canonical form. 
\end{remark}

%-=-=-= EXAMPLE
\begin{example}[id:20141120-202042] \label{20141120-202042}\index{Example!20141120-202042} \hfill \\

Express $5x^1$ in canonical form

\soln

\solnsteps
\begin{align*}
5x && \text{MId} \eqref{eq:mid2} 
\end{align*}
\end{example}

%-=-=-= EXAMPLE
\begin{example}[id:20141121-093439] \label{20141121-093439}\index{Example!20141121-093439} \hfill \\

Express $7x^1+5$ in canonical form.

\soln

\solnsteps
\begin{align*}
7x+5 && \text{MId} \eqref{eq:mid2} \\
\end{align*}
\end{example}

%-=-=-=-=-=-=-=-=-=-=-=-=-=-=-=-=
%	SUBSECTION: 
%-=-=-=-=-=-=-=-=-=-=-=-=-=-=-=-=
\subsection*{Monomials}

\begin{essentialq}\hfill \\

\begin{enumerate}
	\item How do we simplify univariate polynomial expressions?
\end{enumerate}

\end{essentialq}

%-=-=-= DEFINITION
\begin{definition}[Additive Identity (AId)]\index{Additive Identity}
\begin{subequations}
\begin{align}
a+\alert{0} &= a \label{eq:aid1} \\
a &= a+\alert{0} \label{eq:aid2} 
\end{align}
\end{subequations}
\end{definition}

If the constant monomial is 0, the additive identity, then the canonical form of a degree 1 univariate polynomial is a degree 1 monomial.

\subsubsection{Simplifying Univariate Monomial Expressions}

The definition of a univariate monomial expression is based on the expanded canonical form of some polynomial expression.  It might be that the original expression might not be in the expanded canonical form, so a process called \alert{simplifying by expanding} will be introduced to manipulate the expression such that it can be written in its expanded canonical form.

This process of simplifying by expanding polynomial expressions will be developed to the extent that it will be used to simplify multivariate polynomials.  We will start by simplifying univariate monomial expressions.  
  
%-=-=-= DEFINITION
\begin{definition}[Distributive Property Factoring (DPF)]\index{Property!Distributive Property Factoring}
\begin{subequations}
\begin{align}
b\alert{a} + c\alert{a} &= (b+c)\alert{a} \label{eq:dpf1} \\
\alert{a}b + \alert{a}c &= \alert{a}(b+c) \label{eq:dpf2}
\end{align}
\end{subequations}
\end{definition}

%-=-=-= EXAMPLE
\begin{example}[id:20141120-203846] \label{20141120-203846}\index{Example!20141120-203846} \hfill \\

Simplify $6x+7x$

\soln

\solnsteps

Notice that the indeterminate of each monomial is of degree 1; however, the exponent 1 is not shown.  The monomials $6x$ and $7x$ have a like term of $x$.

\begin{align*}
(6+7)x && \text{DPF} \eqref{eq:dpf1} \\
13x && \text{OOA} \eqref{eq:ooa} 
\end{align*}

Notice that the sum of two monomials that have like terms can be found by adding the coefficients of the monomials.  The distributive property in the factoring direction provides some insight to why we can add the coefficients of monomials that have like terms.\\

\soln

\lesssteps
\begin{align*}
13x && \text{OOA} \eqref{eq:ooa} 
\end{align*}

\end{example}

\begin{exercise}[id:20141120-211843]\label{20141120-211843}\index{Exercise!20141120-211843} \hfill \\

Simplify the following univariate polynomial expressions

\begin{multicols}{3}

\begin{enumerate}
	\item $\sage{a_20141120_211843[0]}x+\sage{b_20141120_211843[0]}x$
	\item $\sage{a_20141120_211843[1]}x+\sage{b_20141120_211843[1]}x$
	\item $\sage{a_20141120_211843[2]}x+\sage{b_20141120_211843[2]}x$
	\item $\sage{a_20141120_211843[3]}x+\sage{b_20141120_211843[3]}x$
	\item $\sage{a_20141120_211843[4]}x+\sage{b_20141120_211843[4]}x$
	\item $\sage{a_20141120_211843[5]}x+\sage{b_20141120_211843[5]}x$
	\item $\sage{a_20141120_211843[6]}x+\sage{b_20141120_211843[6]}x$
	\item $\sage{a_20141120_211843[7]}x+\sage{b_20141120_211843[7]}x$
	\item $\sage{a_20141120_211843[8]}x+\sage{b_20141120_211843[8]}x$
	\item $\sage{a_20141120_211843[9]}x+\sage{b_20141120_211843[9]}x$
	\item $\sage{a_20141120_211843[10]}x+\sage{b_20141120_211843[10]}x$
	\item $\sage{a_20141120_211843[11]}x+\sage{b_20141120_211843[11]}x$
\end{enumerate}

\end{multicols}	
Answers: \ref{ans20141120-211843}
\end{exercise}

% TODO section on units as indeterminates

%-=-=-= EXAMPLE
\begin{example}[id:20141027-075159]\label{20141027-075159}\index{Example!20141027-075159} \hfill \\

Simplify $7\unit{cm}+8\unit{cm}$

\soln

\solnsteps
\begin{align*}
(7+8)\alert{\unit{cm}}  &&\text{DPF}\eqref{eq:dpf1} \\
15\unit{cm}  &&\text{OOA} \eqref{eq:ooa}
\end{align*}
\end{example}

\begin{remark}
Remember, if a monomial does not have a coefficient factor, then it's implied that the coefficient factor is 1, the multiplicative identity, and consequently its not explicitly shown.
\end{remark}

%-=-=-= EXAMPLE
\begin{example}[id:20141121-185558] \label{20141121-185558}\index{Example!20141121-185558} \hfill \\

Simplify $x+5x$

\soln

\solnsteps

It can be useful when simplifying expressions to make the multiplicative identity (MId) factor explicit. 

\begin{align*}
1x+5x && \text{MId} \eqref{eq:mid1} \\
(1+5)x && \text{DPF} \eqref{eq:dpf1} \\
6x && \text{OOA} \eqref{eq:ooa} 
\end{align*}

\soln

\lesssteps
\begin{align*}
1x+5x && \text{MId} \eqref{eq:mid1} \\
6x && \text{OOA} \eqref{eq:ooa} 
\end{align*}

As one becomes more experienced, there is no reason to make the multiplicative identity coefficient explicit.

\soln

\lesssteps
\begin{align*}
6x && \text{OOA} \eqref{eq:ooa} 
\end{align*}

\end{example}


%-=-=-= EXERCISE
\begin{exercise}[id:20141121-195608] \label{20141121-195608}\index{Example!20141121-195608} \hfill \\

Simplify the following univariate polynomial expressions

\begin{multicols}{3}

\begin{enumerate}
	\item $x+\sage{b_20141121_195608[0]}x$
	\item $x+\sage{b_20141121_195608[1]}x$
	\item $x+\sage{b_20141121_195608[2]}x$
	\item $x+\sage{b_20141121_195608[3]}x$
	\item $x+\sage{b_20141121_195608[4]}x$
	\item $x+\sage{b_20141121_195608[5]}x$
	\item $\sage{a_20141121_195608[0]}x+x$
	\item $\sage{a_20141121_195608[1]}x+x$
	\item $\sage{a_20141121_195608[2]}x+x$
	\item $\sage{a_20141121_195608[3]}x+x$
	\item $\sage{a_20141121_195608[4]}x+x$
	\item $\sage{a_20141121_195608[5]}x+x$
\end{enumerate}

\end{multicols}	
Answers: \ref{ans20141121-195608}
\end{exercise}


%-=-=-= DEFINITION
\begin{definition}[Definition of Subtraction (DOS)]\index{Definition of Subtraction}
\begin{subequations}
\begin{align}
a-b &= a+\neg b \label{eq:dos1} \\
a+\neg b &= a-b \label{eq:dos2}
\end{align}
\end{subequations}
\end{definition}

%-=-=-= NOTATION
\begin{notation}[Operation of Negation (ONeg)]\index{Operation!Operation of Negation}

\begin{subequations}
\begin{align}
-a &= \neg a \label{eq:oneg1} \\
\neg a &= -a \label{eq:oneg2} 
\end{align}
\end{subequations}

I have used a different symbol, $\neg$, as the prefix negation operator only to differentiate it from the minus sign infix operator symbol, $-$, which is also used as the infix operator for the dyadic operation of subtraction.  I will refer to this change of symbol as ONeg.  This is used only as a teaching tool and should not be confused with the logic negation operator.  Another advantage of using this symbol is that it reduces the number of delimiters used in an expression for example, $\neg a$ versus $(-a)$.

\begin{itemize}
	\item Negative five: $-5$
	\item Negative five: $\neg 5$
	\item Four minus five: $4-5$
	\item Four minus negative five: $4--5$
	\item Four minus negative five: $4-(-5)$
	\item Four minus negative five: $4-\neg 5$
	\item Negative four minus five: $-4-5$
	\item Negative four minus five: $\neg 4-5$
\end{itemize}
	
\end{notation}

%-=-=-= EXAMPLE
\begin{example}[id:20141121-190857] \label{20141121-190857}\index{Example!20141121-190857} \hfill \\

Simplify $8x-6x$

\soln

\solnsteps
\begin{align*}
8x + \neg 6x && \text{DOS} \eqref{eq:dos1} \\
(8+\neg 6)x && \text{DPF} \eqref{eq:dpf1} \\
2x && \text{OOA} \eqref{eq:ooa} 
\end{align*}

\soln

\lesssteps
\begin{align*}
8x + \neg 6x && \text{DOS} \eqref{eq:dos1} \\
2x && \text{OOA} \eqref{eq:ooa} 
\end{align*}

% TODO Possibly add the OOS less steps example

\end{example}


%-=-=-= EXAMPLE
\begin{example}[id:20141121-193636] \label{20141121-193636}\index{Example!20141121-193636} \hfill \\

Simplify $3x-5x$

\soln

\solnsteps
\begin{align*}
3x + \neg 5x && \text{DOS} \eqref{eq:dos1} \\
(3+ \neg 5)x && \text{DPF} \eqref{eq:dpf1} \\
\neg 2x && \text{OOA} \eqref{eq:ooa} \\
-2x && \text{ONeg} \eqref{eq:oneg2} 
\end{align*}

\soln

\lesssteps
\begin{align*}
3x + \neg 5x && \text{DOS} \eqref{eq:dos1} \\
\neg 2x && \text{OOA} \eqref{eq:ooa} \\
-2x && \text{ONeg} \eqref{eq:oneg2} 
\end{align*}

\soln

\lesssteps
\begin{align*}
-2x && \text{OOA} \eqref{eq:ooa} 
\end{align*}

\end{example}

%-=-=-= EXAMPLE
\begin{example}[id:20141106-150622] \label{20141106-150622} \index{Example!20141106-150622} \hfill \\

Simplify $13x-x$

\soln

\solnsteps
\begin{align*}
13x-1x && \text{MId} \eqref{eq:mid1} \\
13x+\neg 1x && \text{DOS} \eqref{eq:dos1} \\
(13+\neg 1)x && \text{DPF} \eqref{eq:dpf1} \\
12x && \text{OOA} \eqref{eq:ooa} 
\end{align*}

\soln

\lesssteps
\begin{align*}
13x + \neg x && \text{DOS} \eqref{eq:dos1} \\
12 x && \text{OOA} \eqref{eq:ooa} 
\end{align*}
\end{example}


\begin{exercise}[id:20141121-191224]\label{20141121-191224}\index{Example!20141121-191224} \hfill \\

Simplify the following univariate polynomial expressions

\begin{multicols}{3}

\begin{enumerate}
	\item $\sage{a_20141121_191224[0]}x-\sage{b_20141121_191224[0]}x$
	\item $\sage{a_20141121_191224[1]}x-\sage{b_20141121_191224[1]}x$
	\item $\sage{a_20141121_191224[2]}x-\sage{b_20141121_191224[2]}x$
	\item $\sage{a_20141121_191224[3]}x-\sage{b_20141121_191224[3]}x$
	\item $\sage{a_20141121_191224[4]}x-\sage{b_20141121_191224[4]}x$
	\item $\sage{a_20141121_191224[5]}x-\sage{b_20141121_191224[5]}x$
	\item $\sage{a_20141121_191224[6]}x-\sage{b_20141121_191224[6]}x$
	\item $\sage{a_20141121_191224[7]}x-\sage{b_20141121_191224[7]}x$
	\item $\sage{a_20141121_191224[8]}x-\sage{b_20141121_191224[8]}x$
	\item $\sage{a_20141121_191224[9]}x-\sage{b_20141121_191224[9]}x$
	\item $\sage{a_20141121_191224[10]}x-\sage{b_20141121_191224[10]}x$
	\item $\sage{a_20141121_191224[11]}x-\sage{b_20141121_191224[11]}x$
\end{enumerate}

\end{multicols}	
Answers: \ref{ans20141121-191224}
\end{exercise}


It is possible for a univariate monomial to have more than two terms in its non-canonical form.  The associative property of addition will be used to help simplify these expressions.

%-=-=-= DEFINITION
\begin{definition}[Associative Property of Addition (APA)]\index{Property!Associative Property of Addition}
\begin{subequations}
\begin{align}
a+b+c &= (a+b)+c \label{eq:apa1} \\
a+b+c &= a+(b+c) \label{eq:apa2}
\end{align}
\end{subequations}
\end{definition}

%-=-=-= EXAMPLE
\begin{example}[id:20141121-184652] \label{20141121-184652}\index{Example!20141121-184652} \hfill \\

Simplify the expression $3x+7x+8x$

\soln
\solnsteps
\begin{align*}
(3x+7x)+8x && \text{APA} \eqref{eq:apa1} \\
(3+7)x+8x && \text{DPF} \eqref{eq:dpf1} \\
10x+8x && \text{OOA} \eqref{eq:ooa} \\
(10+8)x && \text{DPF} \eqref{eq:dpf1} \\
18x && \text{OOA} \eqref{eq:ooa} 
\end{align*}

\soln

\lesssteps
\begin{align*}
(3x+7x)+8x && \text{APA} \eqref{eq:apa1} \\
10x+8x && \text{OOA} \eqref{eq:ooa} \\
18x && \text{OOA} \eqref{eq:ooa} 
\end{align*}

You might have noticed that this expression could be simplified in one step by adding the coefficient of the three monomials $3x$, $7x$ and $8x$, which have the like term $x$.\\

\soln

\lesssteps
\begin{align*}
18x && \text{OOA} \eqref{eq:ooa} 
\end{align*}

\end{example}

%-=-=-= EXAMPLE
\begin{example}[id:20141106-152020] \label{20141106-152020} \index{Example!20141106-152020} \hfill \\

Simplify $4x-2x-x$

\soln

\solnsteps
\begin{align*}
4x-2x-1x && \text{MId} \eqref{eq:mid1} \\
4x+\neg 2x + \neg 1 x && \text{DOS} \eqref{eq:dos1} \\
(4+\neg 2)x + \neg 1x && \text{DPF} \eqref{eq:dpf1} \\ 
2x+\neg 1x && \text{OOA} \eqref{eq:ooa} \\
(2+\neg 1)x && \text{DPF} \eqref{eq:dpf1} \\ 
1x && \text{OOA} \eqref{eq:ooa} \\
x && \text{MId} \eqref{eq:mid2}  
\end{align*}

\soln

\lesssteps
\begin{align*}
4x + \neg 2x + \neg x && \text{DOS} \eqref{eq:dos1} \\
x && \text{OOA} \eqref{eq:ooa} 
\end{align*}
\end{example}

%-=-=-= EXERCISE
\begin{exercise}[id:20141123-091420] \label{20141123-091420}\index{Exercise!20141123-091420} \hfill \\

Simplify the following univariate polynomial expressions

\begin{multicols}{3}

\begin{enumerate}
	\item $\sage{a_20141123_091420[0]}x+\sage{b_20141123_091420[0]}x+\sage{c_20141123_091420[0]}x$
	\item $\sage{a_20141123_091420[1]}x+\sage{b_20141123_091420[1]}x+\sage{c_20141123_091420[1]}x$
	\item $\sage{a_20141123_091420[2]}x+\sage{b_20141123_091420[2]}x+\sage{c_20141123_091420[2]}x$
	\item $\sage{a_20141123_091420[3]}x+\sage{b_20141123_091420[3]}x+\sage{c_20141123_091420[3]}x$
	\item $\sage{a_20141123_091420[4]}x+\sage{b_20141123_091420[4]}x+\sage{c_20141123_091420[4]}x$
	\item $\sage{a_20141123_091420[5]}x+\sage{b_20141123_091420[5]}x+\sage{c_20141123_091420[5]}x$
	\item $\sage{a_20141123_091420[6]}x+\sage{b_20141123_091420[6]}x+\sage{c_20141123_091420[6]}x$
	\item $\sage{a_20141123_091420[7]}x+\sage{b_20141123_091420[7]}x+\sage{c_20141123_091420[7]}x$
	\item $\sage{a_20141123_091420[8]}x+\sage{b_20141123_091420[8]}x+\sage{c_20141123_091420[8]}x$
	\item $\sage{a_20141123_091420[9]}x+\sage{b_20141123_091420[9]}x+\sage{c_20141123_091420[9]}x$
	\item $\sage{a_20141123_091420[10]}x+\sage{b_20141123_091420[10]}x+\sage{c_20141123_091420[10]}x$
	\item $\sage{a_20141123_091420[11]}x+\sage{b_20141123_091420[11]}x+\sage{c_20141123_091420[11]}x$
\end{enumerate}

\end{multicols}	
Answers: \ref{ans20141123-091420}
\end{exercise}

%-=-=-= EXAMPLE
\begin{example}[id:20141108-194431] \label{20141108-194431} \index{Example!20141108-194431} \hfill \\

Simplify $- 3\cdot 7x - 2x \cdot 4$

\soln

\solnsteps
\begin{align*}
\neg 3 \cdot 7x - 2x \cdot 4 && \text{ONeg} \eqref{eq:oneg1} \\
\neg 3 \cdot 7x +\neg 2x \cdot 4 && \text{DOS} \eqref{eq:dos1} \\
\neg 3 \cdot 7 \cdot x + \neg 2 \cdot x \cdot 4 && \text{JTC} \eqref{eq:jtc} \\
\neg 3 \cdot 7 \cdot x + \neg 2 \cdot 4 \cdot x && \text{CPM} \eqref{eq:cpm} \\
(\neg 3 \cdot 7) \cdot x + (\neg 2 \cdot 4) \cdot x && \text{APM} \eqref{eq:apm1} \\
\neg 21 \cdot x + \neg 8 \cdot x && \text{OOM} \eqref{eq:oom} \\
\neg 21 x + \neg 8 x && \text{CTJ} \eqref{eq:ctj} \\
(\neg 21 + \neg 8 )x && \text{DPF} \eqref{eq:dpf1} \\ 
\neg 29x && \text{OOA} \eqref{eq:ooa} \\
-29x && \text{ONeg} \eqref{eq:oneg2} 
\end{align*}

\soln

\lesssteps
\begin{align*}
-3 \cdot 7x + \neg 2x \cdot 4 && \text{DOS} \eqref{eq:dos1} \\
\neg 3 \cdot 7 \cdot x + \neg 2 \cdot 4 \cdot x && \text{CPM} \eqref{eq:cpm} \\
\neg 21x + \neg 8x && \text{OOM} \eqref{eq:oom} \\
-29x && \text{OOA} \eqref{eq:ooa} 
\end{align*}

\end{example}

%-=-=-= EXAMPLE
\begin{example}[id:20141108-194156] \label{20141108-194156} \index{Example!20141108-194156} \hfill \\

Simplify $3 \cdot 5x + 3x \cdot 4$

\soln

\solnsteps
\begin{align*}
3 \cdot 5 \cdot x + 3 \cdot x \cdot 4 && \text{JTC} \eqref{eq:jtc} \\
3 \cdot 5 \cdot x + 3 \cdot 4 \cdot x && \text{CPM} \eqref{eq:cpm} \\
(3 \cdot 5) \cdot x + (3 \cdot 4) \cdot x && \text{APM} \eqref{eq:apm1} \\
15 \cdot x + 12 \cdot x && \text{OOM} \eqref{eq:oom} \\
15x +12x && \text{CTJ} \eqref{eq:ctj} \\
(15+12)x && \text{DPF} \eqref{eq:dpf1} \\
27x && \text{OOA} \eqref{eq:ooa} 
\end{align*}

\soln

\lesssteps
\begin{align*}
3 \cdot 5 \cdot x + 3 \cdot 4 \cdot x && \text{CPM} \eqref{eq:cpm} \\
15x + 12x && \text{OOM} \eqref{eq:oom} \\
27x && \text{OOA} \eqref{eq:ooa}  
\end{align*}

\end{example}

%-=-=-= EXAMPLE
\begin{example}[id:20141108-173613] \label{20141108-173613} \index{Example!20141108-173613} \hfill \\

Simplify $8x \cdot 5$

\soln

\solnsteps
\begin{align*}
8 \cdot x \cdot 5 && \text{JTC} \eqref{eq:jtc} \\
8 \cdot 5 \cdot x && \text{CPM} \eqref{eq:cpm} \\
(8 \cdot 5) \cdot x && \text{APM} \eqref{eq:apm1} \\
40 \cdot x && \text{OOM} \eqref{eq:oom} \\
40x && \text{CTJ} \eqref{eq:ctj} 
\end{align*}
\end{example}

%-=-=-=-=-=-=-=-=-=-=-=-=-=-=-=-=
%	SUBSECTION: 
%-=-=-=-=-=-=-=-=-=-=-=-=-=-=-=-=
\subsection*{Binomials}

%-=-=-= DEFINITION
\begin{definition}[Commutative Property of Addition (CPA)]\index{Property!Commutative Property of Addition}

\begin{align}
\alert{a}b &= b\alert{a} \label{eq:cpa} 
\end{align}
\end{definition}

%-=-=-= DEFINITION
\begin{definition}[Distributive Property Expanding (DPE)]\index{Property!Distributive Property Factoring}
\begin{subequations}
\begin{align}
\alert{a}(b+c) &= \alert{a}b + \alert{a}c \label{eq:dpe1} \\
 (b+c)\alert{a} &= b\alert{a} + c\alert{a}  \label{eq:dpe2}
\end{align}
\end{subequations}
\end{definition}

%-=-=-= EXAMPLE
\begin{example}[id:20141109-090809] \label{20141109-090809} \index{Example!20141109-090809} \hfill \\

Simplify by expanding $5(x+4)$

\soln

\solnsteps
\begin{align*}
5(1x+4) && \text{MId} \eqref{eq:mid1} \\ 
5 \cdot 1x + 5 \cdot 5 && \text{DPF} \eqref{eq:dpf1} \\ 
5 \cdot 1 \cdot x + 5 \cdot 5 && \text{JTC} \eqref{eq:jtc} \\
5 \cdot x + 25  && \text{OOM} \eqref{eq:oom} \\
5x + 25  && \text{CTJ} \eqref{eq:ctj} 
\end{align*}

\soln

\lesssteps
\begin{align*}
5x + 20 && \text{DPE} \eqref{eq:dpe1} 
\end{align*}

\end{example}

%-=-=-= EXAMPLE
\begin{example}[id:20141109-091015] \label{20141109-091015} \index{Example!20141109-091015} \hfill \\

Simplify by expanding $5(3x-9)$

\soln

\solnsteps
\begin{align*}
5(3x+ \neg 9)  && \text{DOS} \eqref{eq:dos1} \\
5 \cdot 3x + 5 \cdot \neg 9  && \text{DPE} \eqref{eq:dpe1} \\  
5 \cdot 3 \cdot x + 5 \cdot \neg 9  && \text{JTC} \eqref{eq:jtc} \\
15 \cdot x + \neg 45  && \text{OOM} \eqref{eq:oom} \\
15x + \neg 45  && \text{CTJ} \eqref{eq:ctj} \\
15x - 45  && \text{DOS} \eqref{eq:dos2}  
\end{align*}

\soln

\lesssteps
\begin{align*}
5(3x+\neg 9) && \text{DOS} \eqref{eq:dos1} \\
15x+ \neg 40 && \text{DPE} \eqref{eq:dpe1} \\
15x-40 && \text{DOS} \eqref{eq:dos2}  
\end{align*}

\end{example}

%-=-=-= EXAMPLE
\begin{example}[id:20141109-092448] \label{20141109-092448} \index{Example!20141109-092448} \hfill \\

Simplify by expanding $-(5x+7)$

\soln

\solnsteps
\begin{align*}
\neg 1(5x+7) && \text{MId} \eqref{eq:mid1} \\ 
\neg 1 \cdot 5x + \neg 1 \cdot 7  && \text{DPE} \eqref{eq:dpe1} \\
\neg 1 \cdot 5 \cdot x + \neg 1 \cdot 7  && \text{JTC} \eqref{eq:jtc} \\
\neg 5 \cdot x + \neg 7 && \text{OOM} \eqref{eq:oom} \\
\neg 5 x + \neg 7  && \text{CTJ} \eqref{eq:ctj} \\
\neg 5x - 7  && \text{DOS} \eqref{eq:dos2} \\
-5x-7  && \text{ONeg} \eqref{eq:oneg2} 
\end{align*}

\soln

\lesssteps
\begin{align*}
-5x-7 && \text{DPE} \eqref{eq:dpe1} 
\end{align*}

\end{example}

%-=-=-= EXAMPLE
\begin{example}[id:20141109-092651] \label{20141109-092651} \index{Example!20141109-092651} \hfill \\

Simplify by expanding $- 13(7x-9)$

\soln

\solnsteps
\begin{align*}
\neg 13 (7x+\neg 9)  && \text{DOS} \eqref{eq:dos1} \\
\neg 13 \cdot 7x + \neg 13 \cdot \neg 9 && \text{DPE} \eqref{eq:dpe1} \\  
\neg 13 \cdot 7 \cdot x + \neg 13 \cdot \neg 9  && \text{JTC} \eqref{eq:jtc} \\
\neg 91 \cdot x + 117  && \text{OOM} \eqref{eq:oom} \\
\neg 91 x + 117  && \text{CTJ} \eqref{eq:ctj} \\
-91x + 117 && \text{ONeg} \eqref{eq:oneg2} 
\end{align*}

\soln

\lesssteps
\begin{align*}
-13(7x+ \neg 9) && \text{DOS} \eqref{eq:dos1} \\
- 91 x + 117 && \text{DPE} \eqref{eq:dpe1} 
\end{align*}

\end{example}

%-=-=-= EXAMPLE
\begin{example}[id:20141109-092910] \label{20141109-092910} \index{Example!20141109-092910} \hfill \\

Simplify by expanding $a(x+b)$ , where $a, b \in \integer$

\soln

\solnsteps
\begin{align*}
a(1x+b)  && \text{MId} \eqref{eq:mid1} \\
a \cdot 1x + a \cdot b  && \text{DPE} \eqref{eq:dpe1} \\ 
a \cdot 1 \cdot x + a \cdot b  && \text{JTC} \eqref{eq:jtc} \\
1 \cdot a \cdot x + a \cdot b  && \text{CPM} \eqref{eq:cpm} \\ 
1ax + ab  && \text{JTC} \eqref{eq:jtc} \\
ax+ab  && \text{MId} \eqref{eq:mid2} 
\end{align*}

\soln

\lesssteps
\begin{align*}
ax +ab && \text{DPE} \eqref{eq:dpe1} 
\end{align*}
\end{example}

%-=-=-= EXAMPLE
\begin{example}[id:20141109-093220] \label{20141109-093220} \index{Example!20141109-093220} \hfill \\

Simplify by expanding $5(x+2)+4$

\soln

\solnsteps
\begin{align*}
5(1x+2)+4 && \text{MId} \eqref{eq:mid1} \\
5 \cdot 1x + 5 \cdot 2 + 4 && \text{DPE} \eqref{eq:dpe1} \\ 
5 \cdot 1 \cdot x + 5 \cdot 2 + 4 && \text{JTC} \eqref{eq:jtc} \\
5 \cdot x +10+4 && \text{OOM} \eqref{eq:oom} \\
5x+10+4 && \text{CTJ} \eqref{eq:ctj} \\
5x+14 && \text{OOA} \eqref{eq:ooa} \\
\end{align*}

\soln

\lesssteps
\begin{align*}
5x + 10 + 4 && \text{DPE} \eqref{eq:dpe1} \\
5x + 14 && \text{OOA} \eqref{eq:ooa} 
\end{align*}
\end{example}

%-=-=-= EXAMPLE
\begin{example}[id:20141109-093419] \label{20141109-093419} \index{Example!20141109-093419} \hfill \\

Simplify by expanding $7x +5(4x+8)$ 

\soln

\solnsteps
\begin{align*}
7x + 5 \cdot 4x + 5 \cdot 8 && \text{DPE} \eqref{eq:dpe1} \\ 
7 \cdot x + 5 \cdot 4 \cdot x + 5 \cdot 8 && \text{JTC} \eqref{eq:jtc} \\
7 \cdot x + 20 \cdot x + 40 && \text{OOM} \eqref{eq:oom} \\
7x + 20x + 40 && \text{CTJ} \eqref{eq:ctj} \\
(7+20)x + 40 && \text{DPF} \eqref{eq:dpf1} \\ 
27x+40 && \text{OOA} \eqref{eq:ooa} 
\end{align*}
\end{example}

%-=-=-= EXAMPLE
\begin{example}[id:20141109-094928] \label{20141109-094928} \index{Example!20141109-094928} \hfill \\

Simplify by expanding $4(3x+4)+x+6$

\soln

\solnsteps
\begin{align*}
4(3x+4)+1x+6 && \text{MId} \eqref{eq:mid1} \\
4 \cdot 3x + 4 \cdot 4 + 1x + 6 && \text{DPE} \eqref{eq:dpe1} \\
4 \cdot 3 \cdot x + 4 \cdot 4 + 1 \cdot x + 6 && \text{JTC} \eqref{eq:jtc} \\
12 \cdot x + 16 + 1 \cdot x + 6 && \text{OOM} \eqref{eq:oom} \\
12x+16+1x+6 && \text{CTJ} \eqref{eq:ctj} \\
12+1x+16+6 && \text{CPA} \eqref{eq:cpa} \\
(12+1)x+16+6 && \text{DPF} \eqref{eq:dpf1} \\ 
13x+22 && \text{OOA} \eqref{eq:ooa} 
\end{align*}

\soln

\lesssteps
\begin{align*}
12x + 16 + x+ 6 && \text{DPE} \eqref{eq:dpe1} \\
12x + x + 16 + 6 && \text{CPA} \eqref{eq:cpa} \\
13x+22 && \text{OOA} \eqref{eq:ooa} 
\end{align*}

\end{example}

%-=-=-= EXAMPLE
\begin{example}[id:20141109-095151] \label{20141109-095151} \index{Example!20141109-095151} \hfill \\

Simplify by expanding $5(x-4)+3x-5$

\soln

\solnsteps
\begin{align*}
5(1x-4)+3x-5 && \text{MId} \eqref{eq:mid1} \\
5(1x + \neg 4)+3x + \neg 5 && \text{DOS} \eqref{eq:dos1} \\
5 \cdot 1x + 5 \cdot \neg 4 + 3x + \neg 5 && \text{DPE} \eqref{eq:dpe1} \\ 
5 \cdot 1 \cdot x + 5 \cdot \neg 4 + 3 \cdot x + \neg 5 && \text{JTC} \eqref{eq:jtc} \\
5 \cdot x + \neg 20 + 3 \cdot x + \neg 5 && \text{OOM} \eqref{eq:oom} \\
5x + \neg 20 + 3x + \neg 5 && \text{JTC} \eqref{eq:jtc} \\
5x + 3x + \neg 20 + \neg 5 && \text{CPA} \eqref{eq:cpa} \\
(5+3)x + \neg 20 + \neg 5 && \text{DPF} \eqref{eq:dpf1} \\
8x + \neg 25 && \text{OOA} \eqref{eq:ooa} \\
8x-25 && \text{DOS} \eqref{eq:dos2} 
\end{align*}

\soln

\lesssteps
\begin{align*}
5(x+ \neg 4)+3x + \neg 5 && \text{DOS} \eqref{eq:dos1} \\
5x + \neg 20 + 3x + \neg 5 && \text{DPE} \eqref{eq:dpe1} \\
5x + 3x + \neg 20 + \neg 5 && \text{CPA} \eqref{eq:cpa} \\
8x + \neg 25 && \text{OOA} \eqref{eq:ooa} \\
8x-25 && \text{DOS} \eqref{eq:dos2} 
\end{align*}

\end{example}

%-=-=-= EXAMPLE
\begin{example}[id:20141109-095536] \label{20141109-095536} \index{Example!20141109-095536} \hfill \\

Simplify by expanding $8x-5-4(x-3)$

\soln

\solnsteps
\begin{align*}
8x-5-4(1x-3) && \text{MId} \eqref{eq:mid1} \\
8x + \neg 5 + \neg 4(1x + \neg 3) && \text{DOS} \eqref{eq:dos1} \\
8x + \neg 5 + \neg 4 \cdot 1x + \neg 4 \cdot \neg 3 && \text{DPE} \eqref{eq:dpe1} \\ 
8 \cdot x + \neg 5 + \neg 4 \cdot 1 \cdot x + \neg 4 \cdot \neg 3 && \text{JTC} \eqref{eq:jtc} \\
8 \cdot x + \neg 5 + \neg 4 \cdot x + 12 && \text{OOM} \eqref{eq:oom} \\
8x + \neg 5 + \neg 4x + 12 && \text{CTJ} \eqref{eq:ctj} \\
8x + \neg 4x + \neg 5 + 12 && \text{CPA} \eqref{eq:cpa} \\
(8+\neg 4)x + \neg 5 + 12 && \text{DPF} \eqref{eq:dpf1} \\
4x + 7 && \text{OOA} \eqref{eq:ooa}
\end{align*}

\soln

\lesssteps
\begin{align*}
8x + \neg 5 + \neg 4(x+\neg 3) && \text{DOS} \eqref{eq:dos1} \\
8x + \neg 5 + \neg 4x + 12 && \text{DPE} \eqref{eq:dpe1} \\ 
8x + \neg 4x + \neg 5 + 12 && \text{CPA} \eqref{eq:cpa} \\
4x+7 && \text{OOA} \eqref{eq:ooa} 
\end{align*}

\end{example}

%-=-=-= EXAMPLE
\begin{example}[id:20141109-095842] \label{20141109-095842} \index{Example!20141109-095842} \hfill \\

Simplify by expanding $5(x+3)+3(x+2)$

\soln

\solnsteps
\begin{align*}
5 \cdot x + 5 \cdot 3 + 3 \cdot x + 3 \cdot 2 && \text{DPE} \eqref{eq:dpe1} \\ 
5 \cdot x + 15 + 3 \cdot x + 6 && \text{OOM} \eqref{eq:oom} \\
5x + 15 + 3x+6 && \text{CTJ} \eqref{eq:ctj} \\
5x +3x+ 15 + 6 && \text{CPA} \eqref{eq:cpa} \\
(3+5)x+15+6 && \text{DPF} \eqref{eq:dpf1} \\
8x+21 && \text{OOA} \eqref{eq:ooa} 
\end{align*}

\soln

\lesssteps
\begin{align*}
5x + 15 +3x+6 && \text{DPE} \eqref{eq:dpe1} \\
5x+3x+15+6 && \text{CPA} \eqref{eq:cpa} \\
8x+21 && \text{OOA} \eqref{eq:ooa} 
\end{align*}

\end{example}

%-=-=-=-=-=-=-=-=-=-=-=-=-=-=-=-=-=-=-=-=-=-=-=-=
%	SECTION: 
%-=-=-=-=-=-=-=-=-=-=-=-=-=-=-=-=-=-=-=-=-=-=-=-=
\section{Degree 2 Univariate Polynomials}\index{Degree 2 Univariate Polynomials}

%-=-=-=-=-=-=-=-=-=-=-=-=-=-=-=-=
%	SUBSECTION: 
%-=-=-=-=-=-=-=-=-=-=-=-=-=-=-=-=
\subsection*{Monomials}

%-=-=-= EXAMPLE
\begin{example}[id:20141106-151138] \label{20141106-151138} \index{Example!20141106-151138} \hfill \\

Simplify $4x^2+12x^2$

\soln

\solnsteps
\begin{align*}
(4+12)x^2 && \text{DPF} \eqref{eq:dpf1} \\
16x^2 && \text{OOA} \eqref{eq:ooa} 
\end{align*}

\soln

\lesssteps
\begin{align*}
16x^2 && \text{OOA} \eqref{eq:ooa} 
\end{align*}

\end{example}

%-=-=-= EXAMPLE
\begin{example}[id:20141106-154547] \label{20141106-154547} \index{Example!20141106-154547} \hfill \\

Simplify $x^2-x+x^2+x$ 

\soln

\solnsteps
\begin{align*}
1x^2-1x+1x^2+1x && \text{MId} \eqref{eq:mid1} \\
1x^2+\neg 1 x + 1x^2+1x && \text{DOS} \eqref{eq:dos1} \\
1x^2+1x^2+\neg 1x+1x && \text{CPA} \eqref{eq:cpa} \\
(1+1)x^2+(\neg 1 + 1)x && \text{DPF} \eqref{eq:dpf1} \\ 
2x^2+0x && \text{OOA} \eqref{eq:ooa} \\
2x^2 && \text{MId} \eqref{eq:mid2} 
\end{align*}

\soln

\lesssteps
\begin{align*}
x^2 + \neg x + x^2 + x && \text{DOS} \eqref{eq:dos1} \\
x^2 + x^2 + \neg x + x && \text{CPA} \eqref{eq:cpa} \\
2x^2 && \text{OOA} \eqref{eq:ooa} 
\end{align*}

\end{example}

% TODO Commutative property of multiplication

%-=-=-= EXAMPLE
\begin{example}[id:20141108-194709] \label{20141108-194709} \index{Example!20141108-194709} \hfill \\

Simplify $- 2x4x - x \cdot-x3$

\soln

\solnsteps
\begin{align*}
-2x4x-1x \cdot 1x3 && \text{MId} \eqref{eq:mid1} \\
\neg 2x4x-1x \cdot 1x3  && \text{ONeg} \eqref{eq:oneg1} \\
\neg 2 \cdot x \cdot 4 \cdot x + \neg 1 \cdot x \cdot 1 \cdot x \cdot 3 && \text{JTC} \eqref{eq:jtc} \\
\neg 2 \cdot 4 \cdot x \cdot x + \neg 1 \cdot \neg 1 \cdot 3 \cdot x \cdot x && \text{CPM} \eqref{eq:cpm} \\
\neg 2 \cdot 4 \cdot x^2  + \neg 1 \cdot \neg 1 \cdot 3 \cdot x^2 && \text{PrCBPo} \eqref{eq:prcbpo1} \\
(\neg 2 \cdot 4) \cdot x^2 + (\neg 1 \cdot \neg 1 \cdot 3) x^2 && \text{APM} \eqref{eq:apm1} \\
\neg 8 x^2 + 3 x^2 && \text{OOM} \eqref{eq:oom} \\
(\neg 8 + 3)x^2 && \text{DPF} \eqref{eq:dpf1} \\
\neg 5 x^2 && \text{OOA} \eqref{eq:ooa} \\
-5x^2 && \text{ONeg} \eqref{eq:oneg2} 
\end{align*}

\soln

\lesssteps
\begin{align*}
- 2x4x + \neg x \cdot-x3  && \text{DOS} \eqref{eq:dos1} \\
\neg 2 \cdot 4 \cdot x \cdot x + 3 \cdot \neg x \cdot \neg x && \text{CPM} \eqref{eq:cpm} \\
\neg 2 \cdot 4 \cdot x^2 + 3 \cdot x^2 && \text{PrCBPo} \eqref{eq:prcbpo1} \\
\neg 8x^2 + 3 x^2 && \text{OOM} \eqref{eq:oom} \\
-5x^2 && \text{OOA} \eqref{eq:ooa} 
\end{align*}

\end{example}

\begin{arule}[Product of a Common Base Powers (PrCBPo)]\index{Powers!Product of Common Base Powers}
\begin{subequations}
\begin{align}
	b^m \cdot b^n &= b^{m+n} \label{eq:prcbpo1}\\
	b^{m+n} &= b^m \cdot b^n \label{eq:prcbpo2}
\end{align}
\end{subequations}
\end{arule}

\begin{arule}[Quotient of a Common Base Powers (QCBPo)]\index{Powers!Quotient of Common Base Powers}
\begin{subequations}
\begin{align}
	\dfrac{b^m}{b^n} &= b^{m-n} \label{eq:qcbpoo1}\\
	b^{m-n} &= \dfrac{b^m}{b^n} \label{eq:qcbpo2}
\end{align}
\end{subequations}
\end{arule}

\begin{arule}[Power of a Power (PoPo)]\index{Powers!Power of a Power}
\begin{subequations}
\begin{align}
	\left(b^m\right)^k &= b^{m \cdot k} \label{eq:popo1}\\
	b^{m \cdot k} &= \left( b^m \right)^k \label{eq:popo2}
\end{align}
\end{subequations}
\end{arule}

%-=-=-= EXAMPLE
\begin{example}[id:20141108-191616] \label{20141108-191616} \index{Example!20141108-191616} \hfill \\

Simplify $-5x \cdot 4x$

\soln

\solnsteps
\begin{align*}
\neg 5x \cdot 4x && \text{ONeg} \eqref{eq:oneg1} \\
\neg 5 \cdot x \cdot 4 \cdot x && \text{JTC} \eqref{eq:jtc} \\
\neg 5 \cdot 4 \cdot x \cdot x && \text{CPM} \eqref{eq:cpm} \\
\neg 5 \cdot 4 \cdot x^2 && \text{PrCBPo} \eqref{eq:prcbpo1} \\
(\neg 5 \cdot 4) \cdot x^2 && \text{APM} \eqref{eq:apm1} \\
\neg 20 \cdot x^2 && \text{OOM} \eqref{eq:oom} \\
\neg 20 x^2 && \text{CTJ} \eqref{eq:ctj} \\
-20x^2 && \text{ONeg} \eqref{eq:oneg2} 
\end{align*}

\soln

\lesssteps
\begin{align*}
\neg 5 \cdot 4 \cdot x \cdot x && \text{CPM} \eqref{eq:cpm} \\
\neg 5 \cdot 4 \cdot x^2 && \text{PrCBPo} \eqref{eq:prcbpo1} \\
-20x^2 && \text{OOM} \eqref{eq:oom}  
\end{align*}

\end{example}


%-=-=-=-=-=-=-=-=-=-=-=-=-=-=-=-=
%	SUBSECTION: 
%-=-=-=-=-=-=-=-=-=-=-=-=-=-=-=-=
\subsection*{Binomials}

%-=-=-= EXAMPLE
\begin{example}[id:20141106-152339] \label{20141106-152339} \index{Example!20141106-152339} \hfill \\

Simplify $3x^2+2x+5x^2+4x$

\soln

\solnsteps
\begin{align*}
3x^2+5x^2+2x+4x && \text{CPA} \eqref{eq:cpa} \\
(3+5)x^2+(2+4)x && \text{DPF} \eqref{eq:dpf1} \\ 
8x^2+6x && \text{OOA} \eqref{eq:ooa} 
\end{align*}

\emph{If needed we could continue and express it in the simplified factored form using the distributive property}	

\begin{align*}
(4x+3)2x && \text{DPF} \eqref{eq:dpf1} \\ 
\end{align*}

\soln

\lesssteps
\begin{align*}
3x^2 + 5x^2 + 2x + 4x && \text{CPA} \eqref{eq:cpa} \\
8x^2+6x && \text{OOA} \eqref{eq:ooa} 
\end{align*}
\end{example}

%-=-=-= EXAMPLE
\begin{example}[id:20141107-121834] \label{20141107-121834} \index{Example!20141107-121834} \hfill \\

Simplify $\left(\sqrt{9-x^2}\right)^2$

\soln

\solnsteps
\begin{align*}
\left(\sqrt{9-1x^2}\right)^2 && \text{MId} \eqref{eq:mid1} \\
\left(\sqrt{9+\neg 1x^2}\right)^2 && \text{DOS} \eqref{eq:dos1} \\
\left[\left(9+\neg x^2\right)^{\alert{\frac{1}{2}}} \right]^2 && \text{RTPo} \eqref{eq:rtpo} \\
9+\neg 1 x^2 && \text{PoPo} \eqref{eq:popo1} \\
\neg 1 x^2 + 9 && \text{CPA} \eqref{eq:cpa} \\
\neg x^2 + 9 && \text{MId} \eqref{eq:mid1} \\
-x^2+9 && \text{ONeg} \eqref{eq:oneg2} 
\end{align*}

\soln

\lesssteps
\begin{align*}
9-x^2 && \text{PoPo} \eqref{eq:popo1} 
\end{align*}


\emph{It might be easier to view this using a variable substitution for the radicand, $9-x^2$. Let $k=9+\neg 1x^2$.}

\begin{align*}
\left(\sqrt{k}\right)^2 && \text{MId} \eqref{eq:mid1} \\
\left(\sqrt{k}\right)^2 && \text{DOS} \eqref{eq:dos1} \\
\left[\left(k\right)^{\alert{\frac{1}{2}}} \right]^2 && \text{RTPo} \eqref{eq:rtpo} \\
k && \text{PoPo} \eqref{eq:popo1} \\
9 + \neg 1 x^2  && \text{CPA} \eqref{eq:cpa} \\
\neg 1 x^2 + 9 && \text{CPA} \eqref{eq:cpa} \\
\neg x^2 + 9 && \text{MId} \eqref{eq:mid1} \\
-x^2+9 && \text{ONeg} \eqref{eq:oneg2} 
\end{align*}

\qdepend 

\qdependlist
example \ref{20141105-144223}-20141105-144223


\end{example}

%-=-=-= EXAMPLE
\begin{example}[id:20141209-145211] \label{20141209-145211}\index{Example!20141209-145211} \hfill \\

Simplify $2x(2x+4)+x^2 \cdot 2 \cdot 1 + 0$

\soln

\solnsteps
\begin{align*}
2x \cdot 2x + 2x \cdot 4 + x^2 \cdot 2 \cdot 1 + 0 && \text{DPE} \eqref{eq:dpe1} \\
2 \cdot x \cdot 2 \cdot x + 2 \cdot x \cdot 4 + x^2 \cdot 2 \cdot 1 + 0 && \text{JTC} \eqref{eq:jtc} \\
2 \cdot 2 \cdot x \cdot x + 2 \cdot 4 \cdot x + 2 \cdot 1 \cdot x^2 + 0 && \text{CPM} \eqref{eq:cpm} \\
2 \cdot 2 \cdot x^2 + 2 \cdot 4 \cdot x + 2 \cdot 1 \cdot x^2 + 0 && \text{PrCBPo} \eqref{eq:prcbpo1} \\
4 \cdot x^2 + 8 \cdot x + 2 \cdot x^2 + 0 && \text{OOM} \eqref{eq:oom} \\
4 x^2 +8x + 2x^2 + 0 && \text{CTJ} \eqref{eq:ctj} \\
4x^2 + 2x^2 + 8x && \text{APA} \eqref{eq:apa1} \\
(4+2)x^2 +8x && \text{DPF} \eqref{eq:dpf2} \\
6x^2+8x && \text{OOA} \eqref{eq:ooa}
\end{align*}

\qdepend 

\qdependlist
example \ref{20141209-144203}-20141209-144203
\end{example}

%-=-=-=-=-=-=-=-=-=-=-=-=-=-=-=-=
%	SUBSECTION: 
%-=-=-=-=-=-=-=-=-=-=-=-=-=-=-=-=
\subsection*{Trinomials}

%-=-=-= EXAMPLE
\begin{example}[id:20141109-133008] \label{20141109-133008}  \index{Example!20141109-133008} \hfill \\

Simplify by expanding $(x+5)(x-8)$

\soln

\solnsteps
\begin{align*}
(1x+5)(1x-8)  && \text{MId} \eqref{eq:mid1} \\
(1x+5)(1x+\neg 8)  && \text{DOS} \eqref{eq:dos1} \\
1x(1x+\neg 8) + 5(1x+\neg 8) && \text{DPE} \eqref{eq:dpe2} \\
1x \cdot 1x + 1x \cdot \neg 8 + 5 \cdot 1x + 5 \cdot \neg 8 && \text{DPE} \eqref{eq:dpe1} \\
1 \cdot x \cdot 1 \cdot x + 1 \cdot x \cdot \neg 8 + 5 \cdot 1 \cdot x + 5 \cdot \neg 8 && \text{JTC} \eqref{eq:jtc} \\
1 \cdot 1 \cdot x \cdot x + \neg 8 \cdot 1 \cdot x + 1 \cdot 5 \cdot x + \neg 8 \cdot 5  && \text{CPM} \eqref{eq:cpm} \\
1 \cdot 1 \cdot x^2 + \neg 8 \cdot 1 \cdot x + 1 \cdot 5 \cdot x + \neg 8 \cdot 5 && \text{PrCBPo} \eqref{eq:prcbpo1} \\
1 \cdot x^2 + \neg 8 \cdot x + 5 \cdot x  + \neg 40  && \text{OOM} \eqref{eq:oom} \\
1  x^2 +  \neg 8 x +5x + \neg 40  && \text{CTJ} \eqref{eq:ctj} \\
1 x^2 + \neg 3 x + \neg 40  && \text{OOA} \eqref{eq:ooa} \\
1 x^2 - 3x - 40  && \text{DOS} \eqref{eq:dos2} \\
x^2-3x-40  && \text{MId} \eqref{eq:mid1} 
\end{align*}

\soln

\lesssteps
\begin{align*}
(x+5)(x+\neg 8) && \text{DOS} \eqref{eq:dos1} \\
x(x+\neg 8)+ 5(x+\neg 8) && \text{DPE} \eqref{eq:dpe2} \\
x^2 + \neg 8x + 5x + \neg 40 && \text{DPE} \eqref{eq:dpe1} \\ 
x^2-3x-40 && \text{OOA} \eqref{eq:ooa} 
\end{align*}

\end{example}

%-=-=-= EXAMPLE
\begin{example}[id:20141109-133316] \label{20141109-133316} \index{Example!20141109-133316} \hfill \\

Simplify by expanding $(x+a)(x+b)$, where $ a, b \in \integer$

\soln

\solnsteps
\begin{align*}
(1x+a)(1x+b)  && \text{MId} \eqref{eq:mid1} \\
1x(1x+b) + a(1x+b)  && \text{DPE} \eqref{eq:dpe2} \\
1x \cdot 1x + 1x \cdot b + a \cdot 1x + a \cdot b  && \text{DPE} \eqref{eq:dpe1} \\ 
1 \cdot x \cdot 1 \cdot x + 1 \cdot x \cdot b + a \cdot 1 \cdot x + a \cdot b  && \text{JTC} \eqref{eq:jtc} \\
1 \cdot 1 \cdot x \cdot x + 1 \cdot b \cdot x + 1 \cdot a \cdot x + a \cdot b  && \text{CPM} \eqref{eq:cpm} \\
1 \cdot 1 \cdot x^2 + 1 \cdot b \cdot x + 1 \cdot a \cdot x + a \cdot b  && \text{PrCBPo} \eqref{eq:prcbpo1} \\
1 \cdot x^2 + 1 \cdot b \cdot x + 1 \cdot a \cdot x + a \cdot b  && \text{OOM} \eqref{eq:oom} \\
1 x^2 + 1 b x + 1 a x + a b  && \text{CTJ} \eqref{eq:ctj} \\
1 x^2 + (1b + 1a)x + ab && \text{DPF} \eqref{eq:dpf1} \\
x^2+(b+a)x+ab  && \text{MId} \eqref{eq:mid2} 
\end{align*}

\soln

\lesssteps
\begin{align*}
x(x+b)+a(x+b) && \text{DPE} \eqref{eq:dpe2} \\ 
x^2 + (b+a)x + ab && \text{DPE} \eqref{eq:dpe1} 
\end{align*}

\end{example}

%-=-=-= EXAMPLE
\begin{example}[id:20141109-140659] \label{20141109-140659} \index{Example!20141109-140659} \hfill \\

Simplify by expanding $(2x+3)(5x+13)$

\soln

\solnsteps
\begin{align*}
2x(5x+13) + 3(5x+13) && \text{DPE} \eqref{eq:dpe2} \\	
2x \cdot 5x + 2x \cdot 13 + 3 \cdot 5x + 3 \cdot 13 && \text{DPE} \eqref{eq:dpe1} \\
2 \cdot x \cdot 5 \cdot x + 2 \cdot x \cdot 13 + 3 \cdot 5 \cdot x + 3 \cdot 13 && \text{JTC} \eqref{eq:jtc} \\
2 \cdot 5 \cdot x \cdot x + 2 \cdot 13 \cdot x + 5 \cdot 3 \cdot x + 3 \cdot 13  && \text{CPM} \eqref{eq:cpm} \\
2 \cdot 5 \cdot x^2 + 2 \cdot 13 \cdot x + 5 \cdot 3 \cdot x + 3 \cdot 13  && \text{PrCBPo} \eqref{eq:prcbpo1} \\
10 \cdot x^2 + 26 \cdot x + 16 \cdot x + 39  && \text{OOM} \eqref{eq:oom} \\
10x^2 + 26x + 15x + 39  && \text{CTJ} \eqref{eq:ctj} \\
10x^2 + 41x + 39  && \text{OOA} \eqref{eq:ooa} 
\end{align*}

\soln

\lesssteps
\begin{align*}
2x(5x+13) + 3(5x+13) && \text{DPE} \eqref{eq:dpe2} \\
10x^2+26x+15x+39 && \text{DPE} \eqref{eq:dpe1} \\
10x^2+41x+39 && \text{OOA} \eqref{eq:ooa} 
\end{align*}
\end{example}

%-=-=-= EXAMPLE
\begin{example}[id:20141109-141019] \label{20141109-141019} \index{Example!20141109-141019} \hfill \\

Simplify by expanding $(-3x-5)(7x+8)$

\soln

\solnsteps
\begin{align*}
(\neg 3x - 5)(7x+8)  && \text{ONeg} \eqref{eq:oneg1} \\
(\neg 3x + \neg 5)(7x+8)  && \text{DOS} \eqref{eq:dos1} \\
\neg 3x(7x+8) + \neg 5(7x+8) && \text{DPE} \eqref{eq:dpe2} \\
\neg 3x \cdot 7x + \neg 3x \cdot 8 + \neg 5 \cdot 7x + \neg 5 \cdot 8 && \text{DPE} \eqref{eq:dpe1} \\
\neg 3 \cdot x \cdot 7 \cdot x + \neg 3 \cdot x \cdot 8 + \neg 5 \cdot 7 \cdot x + \neg 5 \cdot 8 && \text{JTC} \eqref{eq:jtc} \\
\neg 3 \cdot 7 \cdot x \cdot x + \neg 3 \cdot 8 \cdot x + \neg 5 \cdot 7 \cdot x + \neg 5 \cdot 8  && \text{CPM} \eqref{eq:cpm} \\
\neg 3 \cdot 7 \cdot x^2 + \neg 3 \cdot 8 \cdot x + \neg 5 \cdot 7 \cdot x + \neg 5 \cdot 8  && \text{PrCBPo} \eqref{eq:prcbpo1} \\
\neg 21 \cdot x^2 + \neg 24 \cdot x + \neg 35 \cdot x + \neg 40  && \text{OOM} \eqref{eq:oom} \\
\neg 21  x^2 + \neg 24  x + \neg 35  x + \neg 40  && \text{CTJ} \eqref{eq:ctj} \\
\neg 21 x^2 + \neg 59 x + \neg 40  && \text{OOA} \eqref{eq:ooa} \\
\neg 21 x^2 - 59 x - 40  && \text{DOS} \eqref{eq:dos2} \\
- 21 x^2 - 59 x - 40  && \text{ONeg} \eqref{eq:oneg2} 
\end{align*}

\soln

\lesssteps
\begin{align*}
(-3x+\neg 5)(7x+8) && \text{DOS} \eqref{eq:dos1} \\
\neg 3x(7x+8) + \neg 5(7x+8) && \text{DPE} \eqref{eq:dpe2} \\
\neg 21  x^2 + \neg 24  x + \neg 35  x + \neg 40  && \text{CTJ} \eqref{eq:dpe1} \\
-21x^2-59x-40 && \text{OOA} \eqref{eq:ooa} 
\end{align*}
\end{example}

%-=-=-= EXAMPLE
\begin{example}[id:20141109-141347] \label{20141109-141347} \index{Example!20141109-141347} \hfill \\

Simplify by expanding  $(ax+b)(cx+d)$, where $a, b, c, d \in \integer$

\soln

\solnsteps
\begin{align*}
ax(cx+d)+b(cx+d)  && \text{DPE} \eqref{eq:dpe2} \\
ax \cdot cx + ax \cdot d + b \cdot cx + b \cdot d && \text{DPE} \eqref{eq:dpe1} \\
a \cdot x \cdot c \cdot x + a \cdot x \cdot d + b \cdot c \cdot x + b \cdot d   && \text{JTC} \eqref{eq:jtc} \\
a \cdot c \cdot x \cdot x + a \cdot d \cdot x + b \cdot c \cdot x + b \cdot d  && \text{CPM} \eqref{eq:cpm} \\
a \cdot c \cdot x^2 + a \cdot d \cdot x + b \cdot c \cdot x + b \cdot d  && \text{PrCBPo} \eqref{eq:prcbpo1} \\
acx^2 + adx + bcx + bd  && \text{CTJ} \eqref{eq:ctj} \\
acx^2 + (ad+bc)x + bd  && \text{DPF} \eqref{eq:dpf1} 
\end{align*}

\soln

\lesssteps
\begin{align*}
ax(cx+d)+b(cx+d) && \text{DPE} \eqref{eq:dpe2} \\
acx^2 + (ad+bc)x + bd && \text{DPE} \eqref{eq:dpe1} 
\end{align*}

\end{example}


\begin{example}[id:20141105-161225]\label{20141105-161225} \index{Example!20141105-161225} \hfill \\

Simplify $\left(2-\dfrac{x}{2} \right)^2$ by expanding.

\soln

\solnsteps
\begin{align*}
\left(2-\dfrac{1x}{2} \right)^2 && \text{MId} \eqref{eq:mid1} \\
\left(2+ \neg \dfrac{1x}{2} \right)^2 && \text{DOS} \eqref{eq:dos1} \\
\left(2+ \neg \dfrac{1x}{2} \right)\left(2+ \neg \dfrac{1x}{2} \right) && \text{PoTF} \eqref{eq:potf} \\
2\left(2+ \neg \dfrac{1x}{2} \right) + \neg \dfrac{1x}{2}\left(2+ \neg \dfrac{1x}{2} \right) && \text{DPE} \eqref{eq:dpe2} \\
2 \cdot 2 + 2 \cdot \neg \dfrac{1x}{2} + \neg \dfrac{1x}{2} \cdot 2 + \neg \dfrac{1x}{2} \cdot \neg \dfrac{1x}{2} && \text{DPE} \eqref{eq:dpe1} \\
2 \cdot 2 + 2 \cdot \neg \dfrac{1}{2} \cdot x + \neg \dfrac{1}{2} \cdot x \cdot 2 + \neg \dfrac{1}{2} \cdot x \cdot \neg \dfrac{1}{2} \cdot x && \text{JTC} \eqref{eq:jtc} \\
2 \cdot 2 + 2 \cdot \neg \dfrac{1}{2} \cdot x + \neg \dfrac{1}{2} \cdot 2 \cdot x + \neg \dfrac{1}{2} \cdot \neg \dfrac{1}{2} \cdot x \cdot x && \text{CPM} \eqref{eq:cpm} \\
2 \cdot 2 + 2 \cdot \neg \dfrac{1}{2} \cdot x + \neg \dfrac{1}{2} \cdot 2 \cdot x + \neg \dfrac{1}{2} \cdot \neg \dfrac{1}{2} \cdot x^2 && \text{PrCBPo} \eqref{eq:prcbpo1} \\
4 + \neg 1 \cdot x + \neg 1 \cdot x +  \dfrac{1}{4} \cdot x^2 && \text{OOM} \eqref{eq:oom} \\
4 + \neg 1 x + \neg 1 x +  \dfrac{1}{4} x^2 && \text{CTJ} \eqref{eq:ctj} \\
\dfrac{1}{4} x^2 + \neg 1 x + \neg 1 x + 4 && \text{CPA} \eqref{eq:cpa} \\
\dfrac{1}{4} x^2 + \neg 2 x + 4 && \text{OOA} \eqref{eq:ooa} \\
\dfrac{1}{4} x^2 - 2 x + 4 && \text{DOS} \eqref{eq:dos2} 
\end{align*}

\soln

\lesssteps
\begin{align*}
7x + 20x + 40 && \text{DPE} \eqref{eq:dpe1} \\ 
27x + 40 && \text{OOA} \eqref{eq:ooa} 
\end{align*}
\end{example}



%-=-=-=-=-=-=-=-=-=-=-=-=-=-=-=-=-=-=-=-=-=-=-=-=
%	SECTION: 
%-=-=-=-=-=-=-=-=-=-=-=-=-=-=-=-=-=-=-=-=-=-=-=-=
\section{Degree 3 Univariate Polynomials}\index{Degree 3 Univariate Polynomials}

%-=-=-=-=-=-=-=-=-=-=-=-=-=-=-=-=
%	SUBSECTION: 
%-=-=-=-=-=-=-=-=-=-=-=-=-=-=-=-=
\subsection*{Monomials}

%-=-=-=-=-=-=-=-=-=-=-=-=-=-=-=-=
%	SUBSECTION: 
%-=-=-=-=-=-=-=-=-=-=-=-=-=-=-=-=
\subsection*{Binomials}

%-=-=-=-=-=-=-=-=-=-=-=-=-=-=-=-=
%	SUBSECTION: 
%-=-=-=-=-=-=-=-=-=-=-=-=-=-=-=-=
\subsection*{Trinomials}

%-=-=-=-=-=-=-=-=-=-=-=-=-=-=-=-=
%	SUBSECTION: 
%-=-=-=-=-=-=-=-=-=-=-=-=-=-=-=-=
\subsection*{Polynomials}

%-=-=-=-=-=-=-=-=-=-=-=-=-=-=-=-=-=-=-=-=-=-=-=-=
%	SECTION: 
%-=-=-=-=-=-=-=-=-=-=-=-=-=-=-=-=-=-=-=-=-=-=-=-=
\section{Degree $n$ Univariate Polynomials}\index{Degree $n$ Univariate Polynomials}

%-=-=-=-=-=-=-=-=-=-=-=-=-=-=-=-=
%	SUBSECTION: 
%-=-=-=-=-=-=-=-=-=-=-=-=-=-=-=-=
\subsection*{Monomials}



%-=-=-=-=-=-=-=-=-=-=-=-=-=-=-=-=
%	SUBSECTION: 
%-=-=-=-=-=-=-=-=-=-=-=-=-=-=-=-=
\subsection*{Binomials}

%-=-=-=-=-=-=-=-=-=-=-=-=-=-=-=-=
%	SUBSECTION: 
%-=-=-=-=-=-=-=-=-=-=-=-=-=-=-=-=
\subsection*{Trinomials}

%-=-=-=-=-=-=-=-=-=-=-=-=-=-=-=-=
%	SUBSECTION: 
%-=-=-=-=-=-=-=-=-=-=-=-=-=-=-=-=
\subsection*{Polynomials}
%!TEX root = /Users/markholson/Dropbox/+Projects/LatexFiles/MathNotebook/20150312-163541-rs2.2N-MarksMathNotebook
%-=-=-=-=-=-=-=-=-=-=-=-=-=-=-=-=-=-=-=-=-=-=-=-=
%
%	CHAPTER 
%
%-=-=-=-=-=-=-=-=-=-=-=-=-=-=-=-=-=-=-=-=-=-=-=-=

\chapterimage{chapter_head_2.pdf} % Chapter heading image

\chapter{Multivariate Polynomial Expressions}

%-=-=-=-=-=-=-=-=-=-=-=-=-=-=-=-=-=-=-=-=-=-=-=-=
%	SECTION: OPERATION OF ADDITION
%-=-=-=-=-=-=-=-=-=-=-=-=-=-=-=-=-=-=-=-=-=-=-=-=
\section*{Classification of Multivariate Polynomial Expressions}\index{Classification of Multivariate Polynomial Expressions}

%-=-=-=-=-=-=-=-=-=-=-=-=-=-=-=-=-=-=-=-=-=-=-=-=
%	SECTION: 
%-=-=-=-=-=-=-=-=-=-=-=-=-=-=-=-=-=-=-=-=-=-=-=-=
\section*{Degree -1 Multivariate Polynomials}\index{Degree -1 Multivariate Polynomials}

%-=-=-=-=-=-=-=-=-=-=-=-=-=-=-=-=
%	SUBSECTION: 
%-=-=-=-=-=-=-=-=-=-=-=-=-=-=-=-=
\subsection*{Monomials}

%-=-=-=-=-=-=-=-=-=-=-=-=-=-=-=-=-=-=-=-=-=-=-=-=
%	SECTION: 
%-=-=-=-=-=-=-=-=-=-=-=-=-=-=-=-=-=-=-=-=-=-=-=-=
\section*{Degree 0 Multivariate Polynomials}\index{Degree 0 Multivariate Polynomials}

%-=-=-=-=-=-=-=-=-=-=-=-=-=-=-=-=
%	SUBSECTION: 
%-=-=-=-=-=-=-=-=-=-=-=-=-=-=-=-=
\subsection*{Monomials}

%-=-=-=-=-=-=-=-=-=-=-=-=-=-=-=-=-=-=-=-=-=-=-=-=
%	SECTION: 
%-=-=-=-=-=-=-=-=-=-=-=-=-=-=-=-=-=-=-=-=-=-=-=-=
\section*{Degree 1 Multivariate Polynomials}\index{Degree 1 Multivariate Polynomials}

%-=-=-=-=-=-=-=-=-=-=-=-=-=-=-=-=
%	SUBSECTION: 
%-=-=-=-=-=-=-=-=-=-=-=-=-=-=-=-=
\subsection*{Monomials}

%-=-=-=-=-=-=-=-=-=-=-=-=-=-=-=-=
%	SUBSECTION: 
%-=-=-=-=-=-=-=-=-=-=-=-=-=-=-=-=
\subsection*{Binomials}

%-=-=-=-=-=-=-=-=-=-=-=-=-=-=-=-=-=-=-=-=-=-=-=-=
%	SECTION: 
%-=-=-=-=-=-=-=-=-=-=-=-=-=-=-=-=-=-=-=-=-=-=-=-=
\section*{Degree 2 Multivariate Polynomials}\index{Degree 2 Multivariate Polynomials}

%-=-=-=-=-=-=-=-=-=-=-=-=-=-=-=-=
%	SUBSECTION: 
%-=-=-=-=-=-=-=-=-=-=-=-=-=-=-=-=
\subsection*{Monomials}

%-=-=-= EXAMPLE
\begin{example}[id:20141106-150953] \label{20141106-150953} \index{Example!20141106-150953} \hfill \\

Simplify $8xy+5xy$

\soln

\solnsteps
\begin{align*}
(8+5)xy && \text{DPF} \eqref{eq:dpf1} \\
13xy && \text{OOA} \eqref{eq:ooa} 
\end{align*}

\soln

\lesssteps
\begin{align*}
13xy && \text{OOA} \eqref{eq:ooa} 
\end{align*}

\end{example}

%-=-=-= EXAMPLE
\begin{example}[id:20141108-193835] \label{20141108-193835} \index{Example!20141108-193835} \hfill \\

Simplify $3xy + 4yx$

\soln

\solnsteps
\begin{align*}
3 \cdot x \cdot y + 4 \cdot y \cdot x && \text{JTC} \eqref{eq:jtc} \\
3 \cdot x \cdot y + 4 \cdot x \cdot y && \text{CPM} \eqref{eq:cpm} \\
3xy+4xy && \text{CTJ} \eqref{eq:ctj} \\
(3+4)xy && \text{DPF} \eqref{eq:dpf1} \\ 
7xy && \text{OOA} \eqref{eq:ooa} 
\end{align*}

\soln

\lesssteps
\begin{align*}
3xy + 4xy && \text{CPM} \eqref{eq:cpm} \\
7xy && \text{OOA} \eqref{eq:ooa} 
\end{align*}

\end{example}


%-=-=-= EXAMPLE
\begin{example}[id:20141108-180100] \label{20141108-180100} \index{Example!20141108-180100} \hfill \\

Simplify $8x \cdot 2y$

\soln

\solnsteps
\begin{align*}
8 \cdot x \cdot 2 \cdot y && \text{JTC} \eqref{eq:jtc} \\
8 \cdot 2 \cdot x \cdot y && \text{CPM} \eqref{eq:cpm} \\
(8 \cdot 2) \cdot x \cdot y && \text{APM} \eqref{eq:apm1} \\
16 \cdot x \cdot y && \text{OOM} \eqref{eq:oom} \\
16 xy && \text{CTJ} \eqref{eq:ctj}  
\end{align*}

\soln

\lesssteps
\begin{align*}
8 \cdot 2 \cdot x \cdot y && \text{CPM} \eqref{eq:cpm} \\
16xy && \text{OOM} \eqref{eq:oom} 
\end{align*}

\end{example}

%-=-=-= EXAMPLE
\begin{example}[id:20141108-180608] \label{20141108-180608} \index{Example!20141108-180608} \hfill \\

Simplify $3x \cdot y$

\soln

\solnsteps
\begin{align*}
3x \cdot 1y && \text{MId} \eqref{eq:mid1} \\
3 \cdot x \cdot 1 \cdot y && \text{JTC} \eqref{eq:jtc} \\
3 \cdot 1 \cdot x \cdot y && \text{CPM} \eqref{eq:cpm} \\
(3 \cdot 1) \cdot x \cdot y && \text{CPM} \eqref{eq:cpm} \\
3 \cdot x \cdot y && \text{OOM} \eqref{eq:oom} \\
3xy && \text{CTJ} \eqref{eq:ctj} 
\end{align*}

\soln

\lesssteps
\begin{align*}
3xy && \text{CTJ} \eqref{eq:ctj}  
\end{align*}

\end{example}

%-=-=-=-=-=-=-=-=-=-=-=-=-=-=-=-=
%	SUBSECTION: 
%-=-=-=-=-=-=-=-=-=-=-=-=-=-=-=-=
\subsection*{Binomials}

%-=-=-= EXAMPLE
\begin{example}[id:20141106-151622] \label{20141106-151622} \index{Example!20141106-151622} \hfill \\

Simplify $5xy^2-3xy$

\soln

\solnsteps
\begin{align*}
5xy^2+\neg 3xy && \text{DOS} \eqref{eq:dos1} \\
(5y+\neg 3)xy && \text{DPF} \eqref{eq:dpf1} \\
(5y-3)xy && \text{DOS} \eqref{eq:dos2} 
\end{align*}

\soln

\lesssteps
\begin{align*}
5xy^2-3xy 
\end{align*}

\emph{Depending on the context of the problem, both $(5y-3)xy$, the factored form and $5xy^2-3xy$, the expanded form, can both be considered simplified.  }

\end{example}

%-=-=-= EXAMPLE
\begin{example}[id:20141108-180858] \label{20141108-180858} \index{Example!20141108-180858} \hfill \\

Simplify $\dfrac{x}{4} \cdot \dfrac{y}{5}$

\soln

\solnsteps
\begin{align*}
\dfrac{1x}{4} \cdot \dfrac{1y}{5} && \text{MId} \eqref{eq:mid1} \\ 
\dfrac{1}{4} \cdot x \cdot \dfrac{1}{5} \cdot y && \text{JTC} \eqref{eq:jtc} \\
\dfrac{1}{4} \cdot \dfrac{1}{5} \cdot x \cdot y && \text{CPM} \eqref{eq:cpm} \\
\dfrac{1}{20} \cdot x \cdot y && \text{OOM} \eqref{eq:oom} \\
\dfrac{1xy}{20} && \text{CTJ} \eqref{eq:ctj} \\
\dfrac{xy}{20} && \text{MId} \eqref{eq:mid2} 
\end{align*}

\soln

\lesssteps
\begin{align*}
\dfrac{xy}{20} && \text{OOM} \eqref{eq:oom} 
\end{align*}
\end{example}


%-=-=-=-=-=-=-=-=-=-=-=-=-=-=-=-=
%	SUBSECTION: 
%-=-=-=-=-=-=-=-=-=-=-=-=-=-=-=-=
\subsection*{Trinomials}

%-=-=-=-=-=-=-=-=-=-=-=-=-=-=-=-=-=-=-=-=-=-=-=-=
%	SECTION: 
%-=-=-=-=-=-=-=-=-=-=-=-=-=-=-=-=-=-=-=-=-=-=-=-=
\section*{Degree 3 Multivariate Polynomials}\index{Degree 3 Multivariate Polynomials}

%-=-=-=-=-=-=-=-=-=-=-=-=-=-=-=-=
%	SUBSECTION: 
%-=-=-=-=-=-=-=-=-=-=-=-=-=-=-=-=
\subsection*{Monomials}

%-=-=-= EXAMPLE
\begin{example}[id:20141106-151409] \label{20141106-151409} \index{Example!20141106-151409} \hfill \\

Simplify $7xy^2-xy^2$

\soln

\solnsteps
\begin{align*}
7xy^2-1xy^2 && \text{MId} \eqref{eq:mid1} \\
7xy^2+\neg 1xy^2 && \text{DOS} \eqref{eq:dos1} \\
(7+\neg 1)xy^2 && \text{DPF} \eqref{eq:dpf1} \\
6xy^2 && \text{OOA} \eqref{eq:ooa} 
\end{align*}

\soln

\lesssteps
\begin{align*}
7xy^2 + \neg xy^2 && \text{DOS} \eqref{eq:dos1} \\
6xy^2 && \text{OOA} \eqref{eq:ooa} 
\end{align*}

\end{example}

%-=-=-=-=-=-=-=-=-=-=-=-=-=-=-=-=
%	SUBSECTION: 
%-=-=-=-=-=-=-=-=-=-=-=-=-=-=-=-=
\subsection*{Binomials}

%-=-=-=-=-=-=-=-=-=-=-=-=-=-=-=-=
%	SUBSECTION: 
%-=-=-=-=-=-=-=-=-=-=-=-=-=-=-=-=
\subsection*{Trinomials}

%-=-=-=-=-=-=-=-=-=-=-=-=-=-=-=-=
%	SUBSECTION: 
%-=-=-=-=-=-=-=-=-=-=-=-=-=-=-=-=
\subsection*{Polynomials}

%-=-=-=-=-=-=-=-=-=-=-=-=-=-=-=-=-=-=-=-=-=-=-=-=
%	SECTION: 
%-=-=-=-=-=-=-=-=-=-=-=-=-=-=-=-=-=-=-=-=-=-=-=-=
\section*{Degree $n$ Multivariate Polynomials}\index{Degree $n$ Multivariate Polynomials}

%-=-=-=-=-=-=-=-=-=-=-=-=-=-=-=-=
%	SUBSECTION: 
%-=-=-=-=-=-=-=-=-=-=-=-=-=-=-=-=
\subsection*{Monomials}

\begin{arule}[Power of a Product (PoPr)]\index{Powers!Power of a Product}
\begin{subequations}
\begin{align}
	(a \cdot b)^k &= a^k \cdot b^k \label{eq:popr1}\\
	a^k \cdot b^k &= (a \cdot b)^k \label{eq:popr2}
\end{align}
\end{subequations}
\end{arule}

%-=-=-= EXAMPLE
\begin{example}[id:20141108-193145] \label{20141108-193145} \index{Example!20141108-193145} \hfill \\

Simplify $12xy \cdot \left(-\frac{1}{2}\right)yx$

\soln

\solnsteps
\begin{align*}
12x1y \cdot \left(-\frac{1}{2}\right)y1x && \text{MId} \eqref{eq:mid1} \\
12x1y \cdot \neg \frac{1}{2}y1x && \text{ONeg} \eqref{eq:oneg1} \\
12 \cdot x \cdot 1 \cdot y \cdot \neg \frac{1}{2} \cdot y \cdot 1 \cdot x  && \text{JTC} \eqref{eq:jtc} \\
12 \cdot 1 \cdot \neg \frac{1}{2} \cdot 1 \cdot x \cdot x \cdot y \cdot y && \text{CPM} \eqref{eq:cpm} \\
12 \cdot 1 \cdot \neg \frac{1}{2} \cdot 1 \cdot x^2 \cdot y^2 && \text{PrCBPo} \eqref{eq:prcbpo1} \\
\left(12 \cdot 1 \cdot \neg \frac{1}{2} \cdot 1  \right) \cdot x^2 \cdot y^2 && \text{APM} \eqref{eq:apm1} \\
\neg 6 \cdot x^2 \cdot y^2 && \text{OOM} \eqref{eq:oom} \\
\neg 6 x^2y^2 && \text{CTJ} \eqref{eq:ctj} \\
-6x^2y^2 && \text{ONeg} \eqref{eq:oneg2}   
\end{align*}

\soln

\lesssteps
\begin{align*}
12 \cdot \neg \frac{1}{2} \cdot x \cdot x \cdot y \cdot y && \text{CPM} \eqref{eq:cpm} \\
12 \cdot \neg \frac{1}{2} \cdot x^2 \cdot y^2 && \text{PrCBPo} \eqref{eq:prcbpo1} \\
-6x^2y^2 && \text{OOM} \eqref{eq:oom}  
\end{align*}

\end{example}

%-=-=-= EXAMPLE
\begin{example}[id:20141108-192234] \label{20141108-192234} \index{Example!20141108-192234} \hfill \\

Simplify $ac \cdot ab$

\soln

\solnsteps
\begin{align*}
1a1c \cdot 1a1b && \text{MId} \eqref{eq:mid1} \\
1 \cdot a \cdot 1 \cdot c \cdot 1 \cdot a \cdot 1 \cdot b && \text{JTC} \eqref{eq:jtc} \\
1 \cdot 1 \cdot 1 \cdot 1 \cdot a \cdot a \cdot b \cdot c && \text{CPM} \eqref{eq:cpm} \\
1 \cdot 1 \cdot 1 \cdot 1 \cdot a^2 \cdot b \cdot c && \text{PrCBPo} \eqref{eq:prcbpo1} \\
(1 \cdot 1 \cdot 1 \cdot 1) \cdot a^2 \cdot b \cdot c && \text{APM} \eqref{eq:apm1} \\
1 \cdot a^2 \cdot b \cdot c && \text{OOM} \eqref{eq:oom} \\
1 a^2bc && \text{CTJ} \eqref{eq:ctj} \\
a^2bc && \text{MId} \eqref{eq:mid2} 
\end{align*}

\soln

\lesssteps
\begin{align*}
a \cdot a \cdot b \cdot c && \text{CPM} \eqref{eq:cpm} \\
a^2bc && \text{PrCBPo} \eqref{eq:prcbpo1} 
\end{align*}

\end{example}

%-=-=-= EXAMPLE
\begin{example}[id:20141108-192611] \label{20141108-192611} \index{Example!20141108-192611} \hfill \\

Simplify $(-ac) \cdot (-bc)$

\soln

\solnsteps
\begin{align*}
(-1a1c)\cdot (-1b1c) && \text{MId} \eqref{eq:mid1} \\
\neg 1a1c \cdot \neg 1b1c && \text{ONeg} \eqref{eq:oneg1} \\
\neg 1 \cdot a \cdot 1 \cdot c \cdot \neg 1 \cdot b \cdot 1 \cdot c && \text{JTC} \eqref{eq:jtc} \\
\neg 1 \cdot \neg 1 \cdot 1 \cdot 1 \cdot a \cdot b \cdot c \cdot c && \text{CPM} \eqref{eq:cpm} \\
\neg 1 \cdot \neg 1 \cdot 1 \cdot 1 \cdot a \cdot b \cdot c^2 && \text{PrCBPo} \eqref{eq:prcbpo1} \\
(\neg 1 \cdot \neg 1 \cdot 1 \cdot 1) \cdot a \cdot b \cdot c^2 && \text{APM} \eqref{eq:apm1} \\
1 \cdot a \cdot b \cdot c^2 && \text{OOM} \eqref{eq:oom} \\
1abc^2 && \text{CTJ} \eqref{eq:ctj} \\
abc^2 && \text{MId} \eqref{eq:mid2} 
\end{align*}

\soln

\lesssteps
\begin{align*}
\neg a \cdot \neg b \cdot c \cdot c && \text{CPM} \eqref{eq:cpm} \\
abc^2 && \text{PrCBPo} \eqref{eq:prcbpo1} \\  
\end{align*}

\end{example}

%-=-=-=-=-=-=-=-=-=-=-=-=-=-=-=-=
%	SUBSECTION: 
%-=-=-=-=-=-=-=-=-=-=-=-=-=-=-=-=
\subsection*{Binomials}

%-=-=-=-=-=-=-=-=-=-=-=-=-=-=-=-=
%	SUBSECTION: 
%-=-=-=-=-=-=-=-=-=-=-=-=-=-=-=-=
\subsection*{Trinomials}

%-=-=-=-=-=-=-=-=-=-=-=-=-=-=-=-=
%	SUBSECTION: 
%-=-=-=-=-=-=-=-=-=-=-=-=-=-=-=-=
\subsection*{Polynomials}
%!TEX root = /Users/markholson/Dropbox/+Projects/LatexFiles/MathNotebook/20150312-163541-rs2.2N-MarksMathNotebook.tex
%-=-=-=-=-=-=-=-=-=-=-=-=-=-=-=-=-=-=-=-=-=-=-=-=
%
%	CHAPTER 
%
%-=-=-=-=-=-=-=-=-=-=-=-=-=-=-=-=-=-=-=-=-=-=-=-=

\chapterimage{chapter_head_2.pdf} % Chapter heading image

\chapter{Equations}

%-=-=-=-=-=-=-=-=-=-=-=-=-=-=-=-=-=-=-=-=-=-=-=-=
%	SECTION: Equality
%-=-=-=-=-=-=-=-=-=-=-=-=-=-=-=-=-=-=-=-=-=-=-=-=

\section{Equality}\index{Equality}

%-=-=-= PROPERTY
\begin{property}[Reflexive Property of Eqaulity (RPE)]\index{Reflexive Property of Equality}
\begin{subequations}
\begin{align}
a &= a \label{eq:rpe} 
\end{align}
\end{subequations}
\end{property}

%-=-=-= PROPERTY
\begin{property}[Substitution Property of Equality (SPE)]\index{Substitution Property of Equality} \hfill \\

Given $a=b$, then 
\begin{align}
E(a) &= E(b) \label{eq:spe} 
\end{align}
$E(x)$ represents any expression.
\end{property}

%-=-=-= PROPERTY
\begin{property}[Symmetric Property of Equality (SyPE)]\index{Symmetric Property of Equality}
\begin{subequations}
\begin{align}
a =b \quad \text{then} \quad b=a \label{eq:sype} 
\end{align}
\end{subequations}
\end{property}

%-=-=-= PROPERTY
\begin{property}[Transitive Property of Equality (TPE)]\index{Transitive Property of Equality}
\begin{subequations}
\begin{align}
\text{if} \quad a =b \quad \text{and} \quad b=c \quad \text{then} \quad a =c \label{eq:tpe} 
\end{align}
\end{subequations}
\end{property}





%-=-=-=-=-=-=-=-=-=-=-=-=-=-=-=-=-=-=-=-=-=-=-=-=
%	SECTION: Evaluating Functions
%-=-=-=-=-=-=-=-=-=-=-=-=-=-=-=-=-=-=-=-=-=-=-=-=

\section{Solving Linear Equations}\label{Solving Linear Equations}

%-=-=-= EXAMPLE
\begin{example}[id:20141206-102142] \label{20141206-102142}\index{Example!20141206-102142} \hfill \\

Solve the equation $x+a=b$ for $x$

\soln

\solnsteps
\begin{align*}
\farg{x+a} + \neg a  &= \farg{b} + \neg a  && \text{SPE} \eqref{eq:spe} + \text{AI} \eqref{eq:ai} \\
x + (a+\neg a) &= b + \neg a  && \text{APA} \eqref{eq:apa2} \\ 
x + 0 &= b + \neg a  && \text{OOA} \eqref{eq:ooa} \\
x &= b + \neg a  && \text{AId} \eqref{eq:aid2} \\
x &= b-a  && \text{DOS} \eqref{eq:dos2} 
\end{align*}
\end{example}


%-=-=-= EXAMPLE
\begin{example}[id:20141111-222931] \label{20141111-222931}\index{Example!20141111-222931} \hfill \\

Solve the equations $x+8=0$

\soln

\solnsteps
\begin{align*}
\farg{x+8}+\neg 8 &= \farg{0}+\neg 8 && \text{SPE} \eqref{eq:spe} + \text{AI} \eqref{eq:ai} \\
x+(8+\neg 8) &= 0 + \neg 8 && \text{APA} \eqref{eq:apa2} \\
x+0 &= \neg 8 && \text{OOA} \eqref{eq:ooa} \\
x &= \neg 8 && \text{AId} \eqref{eq:aid2} \\
x &= -8 && \text{ONeg} \eqref{eq:oneg2} 
\end{align*}

\soln

\lesssteps
\begin{align*}
\farg{x+8}+\neg 8 &= \farg{0}+\neg 8 && \text{SPE} \eqref{eq:spe} + \text{AI} \eqref{eq:ai} \\
x &= - 8 && \text{OOA} \eqref{eq:ooa} 
\end{align*}

\qdepend 

\qdependlist

example \ref{20141111-190212}-20141111-190212

\end{example}

%-=-=-= EXAMPLE
\begin{example}[id:20141206-101632] \label{20141206-101632}\index{Example!20141206-101632} \hfill \\

Solve the equation $x+4=7$

\soln

\solnsteps
\begin{align*}
\farg{x} +4 + \neg 4  &= \farg{7} + \neg 4  && \text{SPE} \eqref{eq:spe} + \text{AI} \eqref{eq:ai} \\
x + (4+\neg 4) &= 7 + \neg 4  && \text{APA} \eqref{eq:apa2} \\
x + 0 &= 3  && \text{OOA} \eqref{eq:ooa} \\
x  &= 3  && \text{AId} \eqref{eq:aid2}  
\end{align*}

\soln

\lesssteps

\begin{align*}
\farg{x} +4 + \neg 4  &= \farg{7} + \neg 4  && \text{SPE} \eqref{eq:spe} + \text{AI} \eqref{eq:ai} \\
x &= 3  && \text{OOA} \eqref{eq:ooa} 
\end{align*}
\end{example}

%-=-=-= EXAMPLE
\begin{example}[id:20141206-101107] \label{20141206-101107}\index{Example!20141206-101107} \hfill \\

Solve the equation $x-8=15$ for $x$

\soln

\solnsteps
\begin{align*}
x+ \neg 8 &= 15  && \text{DOS} \eqref{eq:dos1} \\
\farg{x+ \neg 8} + 8  &= \farg{15} + 8  && \text{SPE} \eqref{eq:spe} + \text{AI} \eqref{eq:ai} \\
x + (\neg 8+ 8) &= 15 + 8  && \text{APA} \eqref{eq:apa2} \\
x + 0 &= 23  && \text{OOA} \eqref{eq:ooa} \\
x  &= 23  && \text{AId} \eqref{eq:aid2} 
\end{align*}

\soln

\lesssteps
\begin{align*}
\farg{x+ \neg 8} + 8  &= \farg{15} + 8  && \text{SPE} \eqref{eq:spe} + \text{AI} \eqref{eq:ai} \\
x &= 23  && \text{OOA} \eqref{eq:ooa} \\
\end{align*}
\end{example}

%-=-=-= EXAMPLE
\begin{example}[id:20141206-102404] \label{20141206-102404}\index{Example!20141206-102404} \hfill \\

Solve the equation $5x=9$ for $x$.

\soln

\solnsteps
\begin{align*}
\dfrac{1}{5}\farg{5x} &= \dfrac{1}{5}\farg{9} && \text{SPE} \eqref{eq:spe} + \text{MI} \eqref{eq:mi1} \\
\dfrac{1}{5} \cdot \left[ 5 \cdot x \right] &= \dfrac{1}{5} \cdot 9 && \text{JTC} \eqref{eq:jtc} \\
\left(\dfrac{1}{5} \cdot 5 \right) \cdot x &= \dfrac{1}{5} \cdot 9  && \text{APM} \eqref{eq:apm2} \\ 
1\cdot x &= \dfrac{9}{5}  && \text{OOM} \eqref{eq:oom} \\
x &= \dfrac{9}{5}  && \text{MId} \eqref{eq:mid2}  
\end{align*}

\soln

\lesssteps
\begin{align*}
\dfrac{1}{5}\farg{5x} &= \dfrac{1}{5}\farg{9} && \text{SPE} \eqref{eq:spe} + \text{MI} \eqref{eq:mi1} \\ 
x &= \dfrac{9}{5}  && \text{OOM} \eqref{eq:oom} 
\end{align*}
\end{example}

%-=-=-= EXAMPLE
\begin{example}[id:20141206-104404] \label{20141206-104404}\index{Example!20141206-104404} \hfill \\

Solve the equation $ax=b$ for $x$.

\soln

\solnsteps
\begin{align*}
\dfrac{1}{a}\farg{ax} &= \dfrac{1}{a}\farg{b} && \text{SPE} \eqref{eq:spe} + \text{MI} \eqref{eq:mi1} \\
\dfrac{1}{a} \cdot \left( a \cdot x \right) &= \dfrac{1}{a} \cdot b && \text{JTC} \eqref{eq:jtc} \\
\left(\dfrac{1}{a} \cdot a \right) \cdot x &= \dfrac{1}{a} \cdot b  && \text{APM} \eqref{eq:apm2} \\ 
1\cdot x &= \dfrac{b}{a}  && \text{OOM} \eqref{eq:oom} \\
x &= \dfrac{b}{a}  && \text{MId} \eqref{eq:mid2} 
\end{align*}
\end{example}

%-=-=-= EXAMPLE
\begin{example}[id:20141206-102723] \label{20141206-102723}\index{Example!20141206-102723} \hfill \\

Solve the equation $-2x=7$ for $x$

\soln

\solnsteps
\begin{align*}
\neg 2x &= 7  && \text{ONeg} \eqref{eq:oneg1} \\
\neg \dfrac{1}{2}\farg{\neg 2x}  &= \neg \dfrac{1}{2}\farg{7} && \text{SPE} \eqref{eq:spe} + \text{MI} \eqref{eq:mi2} \\
\neg \dfrac{1}{2} \cdot \left( \neg 2 \cdot x \right) &= \neg \dfrac{1}{2} \cdot 7 && \text{JTC} \eqref{eq:jtc} \\
\left(\neg\dfrac{1}{2} \cdot \neg 2 \right) \cdot x &= \neg \dfrac{1}{2} \cdot 7  && \text{APM} \eqref{eq:apm2} \\ 
1\cdot x &= \neg \dfrac{7}{2}  && \text{OOM} \eqref{eq:oom} \\
1\cdot x &= -\dfrac{7}{2}  && \text{ONeg} \eqref{eq:oneg2} \\
x &= -\dfrac{7}{2}  && \text{MId} \eqref{eq:mid2}  
\end{align*}

\soln

\lesssteps
\begin{align*}
\neg \dfrac{1}{2}\farg{\neg 2x}  &= \neg \dfrac{1}{2}\farg{7} && \text{SPE} \eqref{eq:spe} + \text{MI} \eqref{eq:mi2} \\ 
x &= -\dfrac{7}{2}  && \text{OOM} \eqref{eq:oom}
\end{align*}
\end{example}


%-=-=-= EXAMPLE
\begin{example}[id:20141111-215726] \label{20141111-215726}\index{Example!20141111-215726} \hfill \\

Solve the equation $2x+5=0$ for $x$

\soln

\solnsteps
\begin{align*}
\farg{2x+5}+\neg 5 &= \farg{0} + \neg 5 && \text{SPE} \eqref{eq:spe} + \text{AI} \eqref{eq:ai} \\
2x+(5+\neg 5) &= 0+\neg 5 && \text{APA} \eqref{eq:apa1} \\
2x + 0 &= \neg 5 && \text{OOA} \eqref{eq:ooa} \\
2x &= \neg 5 && \text{AId} \eqref{eq:aid1} \\
\dfrac{1}{2}\farg{2x} &= \dfrac{1}{2}\farg{\neg 5} && \text{SPE} \eqref{eq:spe} + \text{MI} \eqref{eq:mi1} \\
\dfrac{1}{2} \cdot 2 \cdot x &= \dfrac{1}{2} \cdot \neg 5 && \text{JTC} \eqref{eq:jtc} \\
\left(\dfrac{1}{2}\cdot 2 \right) \cdot x &= \dfrac{1}{2} \cdot \neg 5 && \text{APM} \eqref{eq:apm1} \\
1 \cdot x &= \dfrac{\neg 5}{2} && \text{OOM} \eqref{eq:oom} \\
1x &= -\dfrac{5}{2} && \text{ONeg} \eqref{eq:oneg2} \\
x &= -\dfrac{5}{2} && \text{MId} \eqref{eq:mid2} 
\end{align*}

\soln

\lesssteps
\begin{align}
\farg{2x+5}+\neg 5 &= \farg{0} + \neg 5 && \text{SPE} \eqref{eq:spe} + \text{AI} \eqref{eq:ai} \\
2x &= \neg 5 && \text{OOA} \eqref{eq:ooa} \\
\dfrac{1}{2}\farg{2x} &= \dfrac{1}{2}\farg{\neg 5} && \text{SPE} \eqref{eq:spe} + \text{MI} \eqref{eq:mi1} \\
x &= -\dfrac{5}{2} && \text{OOM} \eqref{eq:oom} 
\end{align}

\qdepend 

\qdependlist

example \ref{20141111-192213}-20141111-192213


\end{example}

%-=-=-=-=-=-=-=-=-=-=-=-=-=-=-=-=-=-=-=-=-=-=-=-=
%	SECTION: Evaluating Functions
%-=-=-=-=-=-=-=-=-=-=-=-=-=-=-=-=-=-=-=-=-=-=-=-=

\section{Solving Quadratic Equations}\label{Solving Quadratic Equations}

%-=-=-= EXAMPLE
\begin{example}[id:20141107-131748] \label{20141107-131748} \index{Example!20141107-131748} \hfill \\

Solve the equation $2-x^2=0$ for $x$

\soln

\solnsteps
\begin{align*}
2-1x^2 &= 0 && \text{MId} \eqref{eq:mid1} \\
2+\neg 1 x^2 &= 0 && \text{DOS} \eqref{eq:dos1} \\
\left[\alert{2+\neg 1x^2} \right] + 1x^2 &= \left[\alert{0}\right]+ 1x^2 && \text{SPE} \eqref{eq:spe} + \text{AI} \eqref{eq:ai} \\
2 + \left(\neg 1x^2 + 1x^2 \right) &= 0 + 1x^2 && \text{APA} \eqref{eq:apa1} \\
2 + 0 &= 0 + 1x^2 && \text{OOA} \eqref{eq:ooa} \\
2 &= 1x^2 && \text{AId} \eqref{eq:aid1} \\
2 &= x^2 && \text{MId} \eqref{eq:mid2} \\
\alert{\pm} \left[\alert{2} \right]^{\frac{1}{2}} &= \left[\alert{x^2}\right]^{\frac{1}{2}} \\ % TODO Power Inverse
\pm 2^{\frac{1}{2}} &= x && \text{PoPo} \eqref{eq:popo1} \\
\pm \sqrt{2} &= x && \text{PoTR} \eqref{eq:potr} \\
x &= \pm \sqrt{2} && \text{SyPE} \eqref{eq:sype} 
\end{align*}
\end{example}

\subsection{Completing the Square}\label{Completing the Square}

Completing the square is an algebraic algorithm used to find the solutions of quadratic equations of the form, $ax^2+bx+c=0$.   Essentially, we want to manipulate this equation such that $x=\text{some values}$.  To gain some understanding of how this algorithm works, we will consider each step individually.  Let's begin with a quadratic equation in the general form: $ax^2+bx+c=0$.\\

Since we are trying to manipulate the equation  $ax^2+bx+c=0$ such that $x=\text{some value}$, we first want the coefficient factor $a$ to be equal to 1.  This is done by multiplying both expressions by the multiplicative inverse of $a$ (step 1) followed by simplifying both expressions.


\begin{align*}
\dfrac{1}{a} \farg{ax^2 + bx +c} & = \dfrac{1}{a} \farg{0} && \text{SPE} \eqref{eq:spe} + \text{MI} \eqref{eq:mi1} \\
\dfrac{1}{a} \cdot ax^2 + \dfrac{1}{a} \cdot bx + \dfrac{1}{a} \cdot c  & = \dfrac{1}{a} \left[ 0\right] && \text{DPE} \eqref{eq:dpe1} \\
\dfrac{1}{a} \cdot a \cdot x^2 + \dfrac{1}{a} \cdot b \cdot x + \dfrac{1}{a} \cdot c  & = \dfrac{1}{a} \left[ 0\right] && \text{JTC} \eqref{eq:jtc} \\
x^2 + \dfrac{b}{a} \cdot x + \dfrac{c}{a} & = 0 && \text{OOM} \eqref{eq:oom} \\
x^2 +\dfrac{b}{a}x + \dfrac{c}{a}  & = 0 && \text{CTJ} \eqref{eq:ctj} 
\end{align*}

We now have three terms in the left hand expression where the first two terms have at least one variable factor, $x$.  The problem is that the third term is a constant and we want $x=\text{some value}$.  This text step is to get rid of the $\dfrac{c}{a}$ term, which can be done by using the additive inverse followed by simplifying both expressions.\\

\begin{align*}
\left[x^2 +\dfrac{b}{a}x + \dfrac{c}{a} \right]  + \neg \dfrac{c}{a}  & = \left[ 0 \right]   + \neg \dfrac{c}{a} && && \text{SPE} \eqref{eq:spe} + \text{AI} \eqref{eq:ai} \\
x^2 +\dfrac{b}{a}x + \dfrac{c}{a} + \neg \dfrac{c}{a}  & = 0  + \neg \dfrac{c}{a} && && \text{APA} \eqref{eq:apa1} \\
x^2 +\dfrac{b}{a}x  & = \neg \dfrac{c}{a} && \text{OOA} \eqref{eq:ooa} \\
\end{align*}



The next step is called completing the square.  The idea is to add a \textit{NeW} constant, $\color{stockholmPink}k$, to the left-hand expression, $x^2 +\dfrac{b}{a}x+\color{stockholmPink}k$,  such that the quadratic expression can then be factored as two identical factors, $({\color{stockholmBlue}x}+{\color{stockholmPink}m})({\color{stockholmBlue}x}+{\color{stockholmPink}m}) = ({\color{stockholmBlue}x}+{\color{stockholmPink}m})^2$, where  ${{\color{stockholmPink}k}=\color{stockholmPink}m} \cdot {\color{stockholmPink}m}$.  Notice that since we are adding a constant term, $\color{stockholmPink}k$, to the left-hand expression, then we must also add this constant, ${\color{stockholmPink}k}$, to the right-hand expression, $x^2 +\dfrac{b}{a}x+{\color{stockholmPink}k} = \neg \dfrac{c}{a}+{\color{stockholmPink}k}$.  To determine the values of both ${\color{stockholmPink}m}$ and ${\color{stockholmPink}k}$ we should refer to the the organization of the two factors that make up the product of a quadratic expression, $x^2 +\dfrac{b}{a}x+{\color{stockholmPink}k}= ({\color{stockholmBlue}x}+{\color{stockholmPink}m})^2$.\\

%\begin{figure}[htbp]
%\begin{center}
%\begin{tikzpicture}[description/.style={fill=white, inner sep=2pt}]
%	\matrix (m) [matrix of math nodes, row sep=3em,
%	column sep=2.5em, text height=1.5ex, text depth=0.25ex]
%	{\textcolor{stockholmBlue}{x} 	&&	\textcolor{stockholmPink}{m} &	\textcolor{stockholmPink}{m}\textcolor{stockholmBlue}{x} \\
%	 \textcolor{stockholmBlue}{x}	&&	\textcolor{stockholmPink}{m}  & \textcolor{stockholmPink}{m}\textcolor{stockholmBlue}{x} \\
%	 \underbrace{\textcolor{stockholmBlue}{x \cdot x}}_{\text{1st Term}} && \underbrace{\textcolor{stockholmPink}{k}}_{\text{Last Term}} & \underbrace{\dfrac{b}{a}x}_{\text{Middle Terms}}\\};
%	\path[->, font=\scriptsize]
%	(m-1-1) 	edgenode[description] {$\times$} (m-2-3)
%			edge  node[description] {$\times$} (m-2-1)
%	(m-2-1)	edge  node[description] {$\times$} (m-1-3)
%	(m-1-3)	edge  (m-1-4)
%	(m-2-3)	edge (m-2-4)
%	(m-1-4)	edge node[description] {$+$}(m-2-4)
%	(m-2-1)	edge  (m-3-1)
%	(m-2-3)	edge  (m-3-3)
%	(m-2-4)	edge  (m-3-4)			
%	(m-1-3)	edge  node[description] {$\times$} (m-2-3);
%\end{tikzpicture}

%\caption{The Organization of the Distributive Property}
%\end{center}
%\end{figure}

Since both factors of this new quadratic expression are the same, both terms that make up the middle term,  must also be the same.  We know that ${\color{stockholmPink}m}{\color{stockholmBlue}x} + {\color{stockholmPink}m}{\color{stockholmBlue}x}=\dfrac{b}{a}x$, so we should be able to determine the value of ${\color{stockholmPink}m}$ from this equation.    If we can determine the value of ${\color{stockholmPink}m}$, then we can determine the value of  ${\color{stockholmPink}k}$.

\begin{align*}
\dfrac{b}{a}x 	& = \textcolor{stockholmPink}{m}\textcolor{stockholmBlue}{x} + \textcolor{stockholmPink}{m}\textcolor{stockholmBlue}{x}\\
		& = 2mx && \text{OOA} \eqref{eq:ooa} \\
\end{align*}

Solving for $m$,

\begin{align*}
2mx & = \dfrac{b}{a}x \\
2 \cdot m \cdot x &= \dfrac{b}{a} \cdot x  && \text{JTC} \eqref{eq:jtc} \\
\dfrac{1}{2} \left[2 \cdot m \cdot x \right] &= \dfrac{1}{2} \left[\dfrac{b}{a} \cdot x \right]  && \text{SPE} \eqref{eq:spe} + \text{MI} \eqref{eq:mi1} \\
m \cdot  x &= \dfrac{b}{2a} \cdot x  && \text{OOM} \eqref{eq:oom} \\
\left[m \cdot x \right]\dfrac{1}{x}  &= \left[\dfrac{b}{2a}  \cdot x\right] \dfrac{1}{x}  && \text{SPE} \eqref{eq:spe} + \text{MI} \eqref{eq:mi1} \\
m &= \dfrac{b}{2a}  && \text{OOM} \eqref{eq:oom} 
\end{align*}


%\begin{figure}[htbp]
%\begin{center}
%\begin{tikzpicture}[description/.style={fill=white, inner sep=2pt}]
%	\matrix (m) [matrix of math nodes, row sep=3em,
%	column sep=2.5em, text height=1.5ex, text depth=0.25ex]
%	{\textcolor{stockholmBlue}{x} 	&&	\textcolor{stockholmPink}{m} &	\dfrac{b}{2a} \cdot x   \\
%	 \textcolor{stockholmBlue}{x}	&&	\textcolor{stockholmPink}{m}  &  \dfrac{b}{2a} \cdot x  \\
%	 \underbrace{\textcolor{stockholmBlue}{x \cdot x}}_{\text{1st Term}} && \underbrace{\textcolor{stockholmPink}{k}}_{\text{Last Term}} & \underbrace{\dfrac{b}{a} x}_{\text{Middle Terms}}\\};
%	\path[->, font=\scriptsize]
%	(m-1-1) 	edge node[description] {$\times$} (m-2-3)
%			edge  node[description] {$\times$} (m-2-1)
%	(m-2-1)	edge  node[description] {$\times$} (m-1-3)
%	(m-1-3)	edge  (m-1-4)
%	(m-2-3)	edge (m-2-4)
%	(m-1-4)	edge node[description] {$+$}(m-2-4)
%	(m-2-1)	edge  (m-3-1)
%	(m-2-3)	edge  (m-3-3)
%	(m-2-4)	edge  (m-3-4)			
%	(m-1-3)	edge  node[description] {$\times$} (m-2-3);
%\end{tikzpicture}

%\caption{We now have the two terms, $\dfrac{b}{2a}x$ that make up the sum of the middle term $\dfrac{b}{a}x$}
%\end{center}
%\end{figure}

%\newpage
%\begin{figure}[htbp]
%\begin{center}
%\begin{tikzpicture}[description/.style={fill=white, inner sep=2pt}]
%	\matrix (m) [matrix of math nodes, row sep=3em,
%	column sep=2.5em, text height=1.5ex, text depth=0.25ex]
%	{\textcolor{stockholmBlue}{x} 	&&	\textcolor{stockholmPink}{b/2a}	& x \cdot b/2a \\
%	 \textcolor{stockholmBlue}{x}	&&	\textcolor{stockholmPink}{b/2a}  & x \cdot b/2a\\
%	 \underbrace{\textcolor{stockholmBlue}{x \cdot x}}_{\text{1st Term}} && \underbrace{\textcolor{stockholmPink}{(b/2a)^2}}_{\text{Last Term}} & %\underbrace{bx/a}_{\text{Middle Terms}}\\};
%	\path[->, font=\scriptsize]
%	(m-1-1) 	edge node[description] {$\times$} (m-2-3)
%			edge  node[description] {$\times$} (m-2-1)
%	(m-2-1)	edge  node[description] {$\times$} (m-1-3)
%	(m-1-3)	edge  (m-1-4)
%	(m-2-3)	edge (m-2-4)
%	(m-1-4)	edge node[description] {$+$}(m-2-4)
%	(m-2-1)	edge  (m-3-1)
%	(m-2-3)	edge  (m-3-3)
%	(m-2-4)	edge  (m-3-4)			
%	(m-1-3)	edge  node[description] {$\times$} (m-2-3);
%\end{tikzpicture}

%\caption{Doing reverse Cross-Check, we find the second term of both factors to be $\dfrac{b}{2a}$ and the  third term to be $\left( \dfrac{b}{2a} \right)^2$.}
%\end{center}
%\end{figure}

\begin{align*}
\left[x^2+\dfrac{b}{a}x \right]+\left( \dfrac{b}{2a} \right)^2 &= \left[\neg \dfrac{c}{a}\right] + \left( \dfrac{b}{2a} \right)^2 && \text{SPE} \eqref{eq:spe} \\
x^2+\dfrac{b}{a}x +\left( \dfrac{b}{2a} \right)^2 &= \neg \dfrac{c}{a}+ \left( \dfrac{b}{2a} \right)^2 && \text{APA} \eqref{eq:apa2} \\
\left( x+ \dfrac{b}{2a} \right) \left( x+ \dfrac{b}{2a} \right)& = \neg \dfrac{c}{a} + \left( \dfrac{b}{2a} \right)^2 && \text{DPF} \eqref{eq:dpf2} \\
\left( x+ \dfrac{b}{2a} \right)^2 &=\neg \dfrac{c}{a} + \left( \dfrac{b}{2a} \right)^2 && \text{PoTF} \eqref{eq:potf} \\
\left( x+ \dfrac{b}{2a} \right)^2 &=\neg \dfrac{c}{a} +  \dfrac{\left(b \right)^2}{\left(2a \right)^2}  && \text{PoQPo} \eqref{eq:poqpo1} \\
\left( x+ \dfrac{b}{2a} \right)^2 &=\neg \dfrac{c}{a} +  \dfrac{b^2}{2^2a^2} && \text{PoPrPo}\\
\left( x+ \dfrac{b}{2a} \right)^2 &=\neg \dfrac{c}{a} +  \dfrac{b^2}{4a^2} && \text{OOE}\\
\left( x+ \dfrac{b}{2a} \right)^2 &=\neg \dfrac{c}{a} \cdot {\color{stockholmPink}\dfrac{4a}{4a}} +  \dfrac{b^2}{4a^2}  && \text{MId}\\
\left( x+ \dfrac{b}{2a} \right)^2 &=\neg \dfrac{c \cdot 4 \cdot a}{a \cdot 4 \cdot a} +  \dfrac{b^2}{4a^2}  && \text{JTC}\\
\left( x+ \dfrac{b}{2a} \right)^2 &=\neg \dfrac{4 \cdot a \cdot c}{4 \cdot a \cdot a} +  \dfrac{b^2}{4a^2}  && \text{CPM}\\
\left( x+ \dfrac{b}{2a} \right)^2 &=\neg \dfrac{4 \cdot a \cdot c}{4 \cdot a^2} +  \dfrac{b^2}{4a^2}  && \text{PrCBPo}\\
\left( x+ \dfrac{b}{2a} \right)^2 &=\neg \dfrac{4 a c}{4a^2} +  \dfrac{b^2}{4a^2}  && \text{CTJ}\\
\left( x+ \dfrac{b}{2a} \right)^2 &= \dfrac{\neg 4ac+b^2}{4a^2} && \text{CD}\\
\left( x+ \dfrac{b}{2a} \right)^2 &= \dfrac{b^2+ \neg 4ac}{4a^2} && \text{CPA}\\
\displaybreak
\left[\left( x+ \dfrac{b}{2a} \right)^2\right]^{\frac{1}{2}} &=\pm \left[\dfrac{b^2+ \neg 4ac}{4a^2}\right]^{\frac{1}{2}} && \text{PoI}\\
 x+ \dfrac{b}{2a}  &=\pm \left[\dfrac{b^2+ \neg 4ac}{4a^2}\right]^{\frac{1}{2}} && \text{PoPrPo}\\
x+ \dfrac{b}{2a}  &=\pm \dfrac{\left[b^2+ \neg 4ac \right]^{\frac{1}{2}}}{\left[4a^2\right]^{\frac{1}{2}}} && \text{PoQPo}\\
x+ \dfrac{b}{2a}  &=\pm \dfrac{\left[b^2+ \neg 4ac \right]^{\frac{1}{2}}}{4^{\frac{1}{2}} a} && \text{PoPrPo}\\
x+ \dfrac{b}{2a}  &=\pm \dfrac{\left[b^2+ \neg 4ac \right]^{\frac{1}{2}}}{2a} && \text{OOE}\\
\left[x+ \dfrac{b}{2a}\right] +  \neg   \dfrac{b}{2a} &= \left[\pm\dfrac{\left[b^2+ \neg 4ac \right]^{\frac{1}{2}}}{2a}\right] + \neg \dfrac{b}{2a} && \text{AI}\\
x+ \dfrac{b}{2a} +  \neg   \dfrac{b}{2a} &= \pm\dfrac{\left[b^2+ \neg 4ac \right]^{\frac{1}{2}}}{2a} + \neg \dfrac{b}{2a} && \text{APA}\\
x+ \dfrac{b}{2a} + \neg   \dfrac{b}{2a} &=\neg   \dfrac{b}{2a} \pm\dfrac{\left[b^2+ \neg 4ac \right]^{\frac{1}{2}}}{2a} && \text{CPA}\\
x &=\neg   \dfrac{b}{2a}  \pm\dfrac{\left[b^2+ \neg 4ac \right]^{\frac{1}{2}}}{2a} && \text{OOA}\\
x &= \dfrac{\neg  b  \pm {\left[b^2+ \neg 4ac \right]^{\frac{1}{2}}}}{2a} && \text{CD}\\
x &= \dfrac{\neg  b  \pm {\left[b^2 - 4ac \right]^{\frac{1}{2}}}}{2a} && \text{DOS}\\
x &= \dfrac{\neg  b  \pm {\sqrt{b^2 - 4ac}}}{2a} && \text{ETR}\\
x &= \dfrac{-  b  \pm {\sqrt{b^2 - 4ac}}}{2a} && \text{ONeg}
\end{align*}

%!TEX root = /Users/markholson/Dropbox/+Projects/LatexFiles/Book/main.tex
%-=-=-=-=-=-=-=-=-=-=-=-=-=-=-=-=-=-=-=-=-=-=-=-=
%
%	CHAPTER 
%
%-=-=-=-=-=-=-=-=-=-=-=-=-=-=-=-=-=-=-=-=-=-=-=-=

\chapterimage{chapter_head_2.pdf} % Chapter heading image

\chapter{Functions}

%-=-=-=-=-=-=-=-=-=-=-=-=-=-=-=-=-=-=-=-=-=-=-=-=
%	SECTION: Evaluating Functions
%-=-=-=-=-=-=-=-=-=-=-=-=-=-=-=-=-=-=-=-=-=-=-=-=

\section{Evaluating Functions}\label{Evaluating Functions}

\begin{example}[id:20141106-083703]\label{20141106-083703}\index{Example!20141106-083703} \hfill \\
Given $f(x)= \dfrac{1}{12}x^3-x^2+4x$, evaluate $f(3)-f(2)$.

\soln

\solnsteps
\begin{align*}
f(3)-f(2) & = \left(\dfrac{\farg{3}^3}{12} - \farg{3}^2 + 4\farg{3} \right) - \left(\dfrac{\farg{2}^3}{12} - \farg{2}^2 + 4\farg{2} \right) \\
&= \left(\dfrac{(3)^3}{12} - (3)^2 + 4(3) \right) - 1 \left(\dfrac{(2)^3}{12} - (2)^2 + 4(2) \right) && \text{MId} \eqref{eq:mid1} \\
&= \left(\dfrac{(3)^3}{12} + \neg  (3)^2 + 4(3)\right) + \neg 1 \left(\dfrac{(2)^3}{12} + \neg (2)^2 + 4(2) \right) && \text{DOS} \eqref{eq:dos1} \\
&= \left(\dfrac{27}{12} + \neg 9 + 4(3) \right)+ \neg 1 \left(\dfrac{8}{12} + \neg 4 + 4(2) \right) && \text{OOE} \eqref{eq:ooe} \\
&= \left(\dfrac{27}{12} + \neg 9 + 12 \right) + \neg 1 \left(\dfrac{8}{12} + \neg 4 + 8 \right) && \text{OOM} \eqref{eq:oom} \\
&= \left(\dfrac{27}{12} + 3 \right) + \neg 1 \left(\dfrac{8}{12} + 4 \right) && \text{OOA} \eqref{eq:ooa} \\
&= \left(\dfrac{27+36}{12} \right) + \neg 1 \left(\dfrac{8+48}{12}\right) && \text{FOOA} \eqref{eq:fooa1} \\
&= \left(\dfrac{63}{12} \right) + \neg 1 \left(\dfrac{56}{12}\right) && \text{OOA} \eqref{eq:ooa} \\
&= \left(\dfrac{63}{12} \right) + \neg \dfrac{56}{12} && \text{OOM} \eqref{eq:oom} \\
&= \dfrac{7}{12} && \text{OOA} \eqref{eq:ooa} 
\end{align*}

\qdepend 

\qdependlist

example \ref{20141105-144223}-20141105-144223


\end{example}

%-=-=-=-=-=-=-=-=-=-=-=-=-=-=-=-=-=-=-=-=-=-=-=-=
%	SECTION: Rational Functions
%-=-=-=-=-=-=-=-=-=-=-=-=-=-=-=-=-=-=-=-=-=-=-=-=

\section{Quadratic Functions}\label{Quadratic Functions}

\subsection{Completing the Square}

\begin{definition}[Completing The Square]
Completing the square is the process used to convert a quadratic polynomial function \(f(x)=ax^2+bx+c\) to the form

\[ f(x)=a \left(x+\dfrac{b}{2a} \right)^2 + \left(c-\dfrac{b^2}{4a} \right) \]

We can simplify this form by defining \(B = \dfrac{b}{2a} \) and \(C=c-\dfrac{b^2}{4a} \), which gives us 

\[ f(x)=a \left(x+B \right) + C \] 

\hfill \cite{mathworld:completethesquare}
\end{definition}

\begin{align*}
f(x) 	&= ax^2+bx+c \\
		&= a \left[x^2+ \dfrac{b}{a}x + \dfrac{c}{a} \right] && \text{DPF} \eqref{eq:dpf2} \\
		&= a \left[x^2 + \dfrac{b}{a}x + \cRed{k} + \cRed{\neg k} + \dfrac{c}{a} \right] && \text{AId} \eqref{eq:aid2} \\
		&= a \left[ \left(x^2 + \dfrac{b}{a}x + \cRed{k} \right)+ \left(\cRed{-k}+ \dfrac{c}{a}  \right) \right] && \text{APA} \eqref{eq:apa2} \\
\end{align*}

\begin{figure}[h!]
\centering
\begin{tikzpicture}[scale=1, auto]

% Place nodes

\node[firstterm](11){$x$}; \node[factoradd,right=of 11](plus1){}; \node[secondterm, right=of plus1](12){$\frac{b}{2a}$};
\node[firstterm, below=of 11](21){$x$}; \node[factoradd,right=of 21](plus2){}; \node[secondterm, right=of plus2](22){$\frac{b}{2a}$};

\node[multiply, below=of 21](31){$x^2$};
\node[multiply, below=of 22](32){$\cRed{k}=\frac{b^2}{4a^2}$};

\node[multiply, right=of 12](13){$\frac{b}{2a}x$};
\node[multiply, right=of 22](23){$\frac{b}{2a}x$};

\node[add, below=of 23](33){$\frac{b}{2}x$};

\path [line](11) edge[bend right=30]node[color=black, midway, left]{$\times$}(21);
\path [line](12) edge[bend left=30]node[color=black,, midway, right]{$\times$}(22);
\path [line](21)--(31);

\path [line](21) edge[bend left=30]node[color=black, pos=0.4, below]{$\times$}(12);
\path [line](11) edge[bend right=30](22);
\path [line](22)--(32);

\path [line](12)--(13);
\path [line](22)--(23);

\path [line](13)--node[color=black, midway, right]{$+$}(23);
\path [line](23)--(33);

\end{tikzpicture}
\caption{Factoring Organizer used to find the value of $\cRed{k}$}
\end{figure}

\begin{align*}
f(x) 	&=  a \left[ \left(x^2 + \dfrac{b}{a}x + \cRed{\dfrac{b^2}{4a^2}} \right)+ \left(\cRed{\dfrac{-b^2}{4a^2}}+ \dfrac{c}{a}  \right) \right]\\
		&= a \left[ \left(x+ \dfrac{b}{2a} \right)^2 + \cRed{\dfrac{-b^2}{4a^2}}+ \dfrac{c}{a} \right]  && \text{DPF} \eqref{eq:dpf2} \\
		&= a \left[  \left(x+ \dfrac{b}{2a} \right)^2 + \dfrac{-b^2}{4a^2}+ \dfrac{4ac}{4a^2} \right] \\ % TODO Common Denominator Step
		&= a \left[  \left(x+ \dfrac{b}{2a} \right)^2 + \dfrac{4ac-b^2}{4a^2} \right]  && \text{OOA} \eqref{eq:ooa} \\
		&= a \left(x+ \dfrac{b}{2a} \right)^2 + \dfrac{4ac-b^2}{4a} && \text{DPE} \eqref{eq:dpe1} \\ 
		&= a \left(x+ \dfrac{b}{2a} \right)^2 + \left(c-\dfrac{b^2}{4a} \right) % TODO Reduce fraction step 
\end{align*}

%-=-=-=-=-=-=-=-=-=-=-=-=-=-=-=-=-=-=-=-=-=-=-=-=
%	SECTION: Rational Functions
%-=-=-=-=-=-=-=-=-=-=-=-=-=-=-=-=-=-=-=-=-=-=-=-=

\section{Rational Functions}\label{Rational Functions}

%-=-=-= DEFINITION
\begin{definition}[Asympote]\index{Asympote}

An asymptote is a line or curve that apporoaches a given curve arbitrarily close.\cite{mathworld:asymptote} \\

A vertical asymptote is a vertical line $x_{va}=c$, that is used to visualize the values of $x$ for which the function is not defined. 

\end{definition}

\begin{figure}[ht]
\begin{center}
	\begin{tikzpicture}
	\begin{axis}[
            domain=-6:2,
            ymax=4,
            ymin=-4,
            samples=100,
            axis lines =middle, xlabel=$x$, ylabel=$y$,
            every axis y label/.style={at=(current axis.above origin),anchor=south},
            every axis x label/.style={at=(current axis.right of origin),anchor=west},
			restrict y to domain=-20:20
          ]
          \addplot [dashed, stockholmPink, smooth] plot coordinates {(-2,4) (-2,-4)}; %% {.451};

          \addplot [very thick, stockholmBlue, smooth] {1/(x+2)};

          \node at (axis cs:-5.7,3.8) [anchor=west] {\color{stockholmPink}Vertical Asymptote};  

        \end{axis}
\end{tikzpicture}

\end{center}
\caption{Vertical Asymptote}
\label{figure:rectangularhyperbola}
\end{figure}

%-=-=-= EXAMPLE
\begin{example}[id:20141111-190212] \label{20141111-190212}\index{Example!20141111-190212} \hfill \\

Find the vertical asymptote of the function $R(x)=\dfrac{7}{x+8}$

\soln

\solnsteps

We are interested in the values of $x$ for which the denominator of $R(x)$ has a value of zero.

\begin{align*}
x_{va}+8 &= 0 \\
x_{va} &=-8  &&\text{solving} 
\end{align*}

\begin{center}
	\begin{tikzpicture}
	\begin{axis}[
            domain=-20:10,
            ymax=10,
            ymin=-10,
            samples=100,
            axis lines =middle, xlabel=$x$, ylabel=$y$,
            every axis y label/.style={at=(current axis.above origin),anchor=south},
            every axis x label/.style={at=(current axis.right of origin),anchor=west},
			restrict y to domain=-20:20
          ]
          \addplot [dashed, stockholmPink, smooth] plot coordinates {(-8,10) (-8,-10)}; %% {.451};

          \addplot [very thick, stockholmBlue, smooth] {7/(x+8)};

          \node at (axis cs:-17.7,5.8) [anchor=west] {\color{stockholmPink}$x_{va}=-8$};  

        \end{axis}
\end{tikzpicture}

\end{center}
\end{example}

%-=-=-= EXAMPLE
\begin{example}[id:20141111-192213] \label{20141111-192213}\index{Example!20141111-192213} \hfill \\

Find the vertical asymptote of the function $R(x)=\dfrac{x+4}{2x+5}$

\soln

\solnsteps

We are interested in the values of $x$ for which the denominator of $R(x)$ has a value of zero.

\begin{align*}
2x_{va}+5 &= 0 \\
x_{va} &= -\dfrac{5}{2}  &&\text{solving \ref{20141111-215726}} \\
\end{align*}

\begin{center}
	\begin{tikzpicture}
	\begin{axis}[
            domain=-6:2,
            ymax=5,
            ymin=-5,
            samples=100,
            axis lines =middle, xlabel=$x$, ylabel=$y$,
            every axis y label/.style={at=(current axis.above origin),anchor=south},
            every axis x label/.style={at=(current axis.right of origin),anchor=west},
			restrict y to domain=-5:5
          ]
          \addplot [dashed, stockholmPink, smooth] plot coordinates {(-5/2,4) (-5/2,-4)}; %% {.451};

          \addplot [very thick, stockholmBlue, smooth] {(x+4)/(2*x+5)};

          \node at (axis cs:-4.7,2.8) [anchor=west] {\color{stockholmPink}$x_{va}=-\frac{5}{2}$};  

        \end{axis}
\end{tikzpicture}

\end{center}
\end{example}








%!TEX root = /Users/markholson/Dropbox/+Projects/LatexFiles/Book/main.tex
%-=-=-=-=-=-=-=-=-=-=-=-=-=-=-=-=-=-=-=-=-=-=-=-=
%
%	CHAPTER 
%
%-=-=-=-=-=-=-=-=-=-=-=-=-=-=-=-=-=-=-=-=-=-=-=-=

\chapterimage{chapter_head_2.pdf} % Chapter heading image

\chapter{Vectors}

%-=-=-=-=-=-=-=-=-=-=-=-=-=-=-=-=-=-=-=-=-=-=-=-=
%	SECTION: 
%-=-=-=-=-=-=-=-=-=-=-=-=-=-=-=-=-=-=-=-=-=-=-=-=
\section{Component \& Geometric Vectors}


\section{Dot Product}

%-=-=-= DEFINITION
\begin{definition}[Dot Product in $\mathbb{R}^2$ (VDPr)]\index{Dot Product $\mathbb{R}^2$}
The dot product of vectors $\vec{u}=(u_x,u_y)$ and $\vec{v}=(v_x, v_y)$ is 

\begin{align}
\vec{u} \cdot \vec{v} 	&= u_x v_x + u_y v_y  \label{eq:vdpr1}
\end{align}

The dot product of vectors $\vec{u}=u_x \ihat + u_y \jhat$ and $\vec{v}=v_x \ihat + v_y \jhat$ is 

\begin{align}
\vec{u} \cdot \vec{v} 	&= u_xv_x \ihat + u_yv_y \jhat \label{eq:vdpr2}
\end{align}

\end{definition}

%-=-=-= EXAMPLE
\begin{example}[id:20141214-195907] \label{20141214-195907}\index{Example!20141214-195907} \hfill \\

Find the dot product of the vectors $\vec{u}=(-5,12)$ and $\vec{v}=(-7,- 2)$

\soln

\solnsteps
\begin{align*}
\vec{u} \cdot \vec{v} &= u_x v_x + u_y v_y && \text{VDPr} \eqref{eq:vdpr1} \\
&=\farg{-5}\farg{-7} + \farg{12}\farg{-2} && \text{SPE} \eqref{eq:spe} \\
&= 35 + -24  && \text{OOM} \eqref{eq:oom} \\
&= 11
\end{align*}
\end{example}

%-=-=-= EXAMPLE
\begin{example}[id:20141214-194701] \label{20141214-194701}\index{Example!20141214-194701} \hfill \\

Find the dot product of the vectors $\vec{u}=-10\ihat - 10 \jhat$ and $\vec{v}=-9\ihat + 7 \jhat$

\soln
\solnsteps
\begin{align*}
\vec{u} \cdot \vec{v} & = \cBlue{u_x}\cRed{v_x} \ihat + \cBlue{u_y}\cRed{v_y} \jhat && \text{VDPr} \eqref{eq:vdpr2} \\
&= \farg{-10}\farg{-9}+ \farg{-10}\farg{7} && \text{SPE} \eqref{eq:spe} \\
&= 90+\neg 70 && \text{OOM} \eqref{eq:oom} \\
&= 20 && \text{OOA} \eqref{eq:ooa} 
\end{align*}


\end{example}
%!TEX root = /Users/markholson/Dropbox/+Projects/LatexFiles/MathNotebook/20150312-163541-rs2.2N-MarksMathNotebook.tex
%-=-=-=-=-=-=-=-=-=-=-=-=-=-=-=-=-=-=-=-=-=-=-=-=
%
%	CHAPTER 
%
%-=-=-=-=-=-=-=-=-=-=-=-=-=-=-=-=-=-=-=-=-=-=-=-=

\chapterimage{chapter_head_2.pdf} % Chapter heading image

\chapter{Differentiation}

%-=-=-=-=-=-=-=-=-=-=-=-=-=-=-=-=-=-=-=-=-=-=-=-=
%	SECTION: 
%-=-=-=-=-=-=-=-=-=-=-=-=-=-=-=-=-=-=-=-=-=-=-=-=
\section{Limit of the Difference Quotient}

%-=-=-= DEFINITION
\begin{definition}[Derivative]\index{Derivative}

The derivative of a function $f(x)$ with respect to the variable $x$ is defined as 
\begin{align}
f'(\cBlue{x}) & \equiv \displaystyle \lim_{\cRed{\Delta x} \to 0} \underbrace{\dfrac{f(\cBlue{x}+\cRed{\Delta x})-f(\cBlue{x})}{\cRed{\Delta x}}}_{\text{Difference Quotient}} \label{eq:dbfp1} 
\end{align}

\end{definition}
\begin{figure}[h!]
\centering
\begin{tikzpicture}
	\begin{axis}[
            domain=0:2, range=0:6,ymax=6,ymin=0,
            axis lines =left, xlabel=$x$, ylabel=$y$,
            every axis y label/.style={at=(current axis.above origin),anchor=south},
            every axis x label/.style={at=(current axis.right of origin),anchor=west},
            xtick={1,1.666}, ytick={1,3},
            xticklabels={$x$,$x+\Delta x$}, yticklabels={$f(x)$,$f(x+\Delta x)$},
            axis on top,
          ]         
          \addplot [stockholmPink!15, domain=(0:2)] {-3.348+4.348*x};
          \addplot [stockholmPink!32, domain=(0:2)] {-2.704+3.704*x};
          \addplot [stockholmPink!49, domain=(0:2)] {-1.994+2.994*x};         
          \addplot [stockholmPink!66, domain=(0:2)] {-1.326+2.326*x}; 
          \addplot [stockholmPink!83, domain=(0:2)] {-0.666+1.666*x};
	  \addplot [stockholmDarkGrey,dashed] plot coordinates {(1,0) (1,1)};
          \addplot [stockholmDarkGrey,dashed] plot coordinates {(0,1) (1,1)};
          \addplot [stockholmDarkGrey,dashed] plot coordinates {(0,3) (1.666,3)};
          \addplot [stockholmDarkGrey,dashed] plot coordinates {(1.666,0) (1.666,3)};
          \addplot [very thick,stockholmBlue, smooth,domain=(0:1.833)] {-1/(x-2)};
          \addplot[color=stockholmDarkGrey,fill=stockholmDarkGrey,only marks,mark=*] coordinates{(1.666,3)};  %% closed hole          
          \addplot[color=stockholmDarkGrey,fill=stockholmDarkGrey,only marks,mark=*] coordinates{(1,1)};  %% closed hole          
          \addplot [very thick,stockholmPink, smooth,domain=(0:2)] {x};
        \end{axis}
\end{tikzpicture}
\caption{\cite{mooculus:textbook}}
\end{figure}
%-=-=-= EXAMPLE
\begin{example}[id:20141219-212546] \label{20141219-212546}\index{Example!20141219-212546} \hfill \\

Differentiate the function $f(x)=5$

\soln

\solnsteps
\begin{align*}
f(x) & = 5x^0 && \text{PoID} \eqref{eq:poid1} \\
f'(x) &= \displaystyle \lim_{\Delta x \to 0} \dfrac{5\farg{x+\Delta x}^0-5\farg{x}^0}{\Delta x} && \text{SPE} \eqref{eq:spe} \& \text{DBFP} \eqref{eq:dbfp1}\\
f'(x) &= \displaystyle \lim_{\Delta x \to 0} \dfrac {5(1)-5(1)}{\Delta x} && \text{PoID} \eqref{eq:poid2} \\
f'(x) &= \displaystyle \lim_{\Delta x \to 0} 0 && \text{OOM} \eqref{eq:oom} \\
f'(x) &= 0 % TODO Evaluate the limit
\end{align*}

\end{example}

%-=-=-=-=-=-=-=-=-=-=-=-=-=-=-=-=-=-=-=-=-=-=-=-=
%	SECTION: 
%-=-=-=-=-=-=-=-=-=-=-=-=-=-=-=-=-=-=-=-=-=-=-=-=
\section{Derivative of a Monomial Functions}

\begin{arule}[Derivative of a Constant (DC)]\index{Derivatives! Derivative of a Constant}
\begin{align}
	\left[c \right]' &= 0 \label{eq:dc1} \\
	\dfrac{\dd}{\dx} \left[c \right]&=0  \label{eq:dc2}
\end{align}
\end{arule}

\begin{arule}[Derivative of a Constant Multiple (DCM)]\index{Derivatives! Derivative of a Constant Multiple}
\begin{align}
	\left[c f(x) \right]' &= c \left[f(x)\right]' \label{eq:dcm1} \\
	\dfrac{\dd}{\dx}\left[c f(x) \right] &= c \dfrac{\dd}{\dx}\left[f(x)\right]  \label{eq:dcm2}
\end{align}
\end{arule}

\begin{arule}[Derivative of a Power (DPo)]\index{Derivatives! Derivative of a Power}
\begin{align}
	\left[x^n \right]'&= n x^{n-1} \label{eq:dpo1} \\
	\dfrac{\dd}{\dx} \left[x^n \right] &= n x^{n-1}  \label{eq:dpo2}
\end{align}
\end{arule}

%-=-=-= EXAMPLE
\begin{example}[id:20141124-153017] \label{20141124-153017}\index{Example!20141124-153017} \hfill \\

Differentiate $f(x)=-3$

\soln

\solnsteps
\begin{align*}
f'(x) &= [-3]' && \text{SPE} \eqref{eq:spe} \\
f'(x) &= 0 && \text{DC} \eqref{eq:dc1} 
\end{align*}

\qdepend

\qdependlist

example \ref{20141124-152503}-20141124-152503


\end{example}


%-=-=-= EXAMPLE
\begin{example}[id:20141124-141850] \label{20141124-141850}\index{Example!20141124-141850} \hfill \\

Differentiate $f(x)=x^2$

\soln

\solnsteps
\begin{align*}
f'(x) &= \left[x^2\right]' && \text{SPE} \eqref{eq:spe} \\
f'(x) &= \cRed{2}x^{\cRed{2-1}} && \text{DPo} \eqref{eq:dpo1} \\
f'(x) &= 2x^1 && \text{OOA} \eqref{eq:ooa} \\
f'(x) &= 2x && \text{MId} \eqref{eq:mid2} 
\end{align*}

\soln

\lesssteps
\begin{align*}
f'(x) & = 2x && \text{DPo} \eqref{eq:dpo1}  
\end{align*}

\qdepend

\qdependlist

example \ref{20141124-152503}-20141124-152503



\end{example}

%-=-=-=-=-=-=-=-=-=-=-=-=-=-=-=-=-=-=-=-=-=-=-=-=
%	SECTION: 
%-=-=-=-=-=-=-=-=-=-=-=-=-=-=-=-=-=-=-=-=-=-=-=-=
\section{Derivative of Polynomial Functions}

\begin{arule}[Derivative of a Sum (DS)]\index{Derivatives! Derivative of a Sum}
\begin{align}
	[f(x) + g(x)]' &= f'(x) + g'(x) \label{eq:ds1} \\
	\dfrac{\dd}{\dx} \left[f(x) + g(x) \right] &= \dfrac{\dd}{\dx} \left[f(x) \right] + \dfrac{\dd}{\dx} \left[g(x) \right]  \label{eq:ds2}
\end{align}
\end{arule}

%-=-=-= EXAMPLE
\begin{example}[id:20141124-152503] \label{20141124-152503}\index{Example!20141124-152503} \hfill \\

Differentiate $f(x)=x^2-3$

\soln

\solnsteps
\begin{align*}
f(x) &= x^2 + \neg 3 && \text{DOS} \eqref{eq:dos1} \\
f'(x) &= \farg{x^2 + \neg 3}' && \text{SPE} \eqref{eq:spe} \\
f'(x) &= \farg{x^2}' + \farg{\neg 3}' && \text{DS} \eqref{eq:ds1} \\
f'(x) &= \farg{x^2}' + 0 && \text{DC} \eqref{eq:dc1} \\
f'(x) &= \farg{x^2}' && \text{AId} \eqref{eq:aid1} \\
f'(x) &= 2x && \text{DPo} \eqref{eq:dpo1}  \text{\, goto \,}  \ref{20141124-141850} 
\end{align*}

\soln

\lesssteps

\begin{align*}
f(x) & = 2x^2 && \text{DPo} \eqref{eq:dpo1} \& \text{DC} \eqref{eq:dpo1} 
\end{align*}

\qdepend

\qdependlist

example \ref{20141124-205219}-20141124-205219

\end{example}

%-=-=-= EXAMPLE
\begin{example}[id:20141128-151834] \label{20141128-151834}\index{Example!20141128-151834} \hfill \\

Differentiate $f(x)=3x^2-6x+4$

\soln

\solnsteps
\begin{align*}
f(x) &= 3x^2 + \neg 6x + 4 && \text{DOS} \eqref{eq:dos1} \\
f'(x) &= \farg{3x^2 + \neg 6x + 4}' && \text{SPE} \eqref{eq:spe} \\
f'(x) &= \farg{3x^2}' + \farg{6x}' + \farg{4}' && \text{DS} \eqref{eq:ds1} \\
f'(x) &= \farg{3x^2}' + \farg{6x}' + 0 && \text{DC} \eqref{eq:dc1} \\
f'(x) &= \farg{3x^2}' + \farg{6x}' && \text{AId} \eqref{eq:aid1} \\
f'(x) &= 3\farg{x^2}' + 6\farg{x}' && \text{DCM} \eqref{eq:dcm1} \\
f'(x) &= 3 \left(2x \right) + 6 \left(1\right) && \text{DPo} \eqref{eq:dpo1} \\
f'(x) &= 6x+6 && \text{OOM} \eqref{eq:oom} 
\end{align*}

\soln
\lesssteps
\begin{align*}
f'(x) & = 6x+6 && \text{DS} \eqref{eq:ds1}  
\end{align*}
\end{example}


\begin{arule}[Derivative of a Product (DPr)]\index{Derivatives! Derivative of a Product}
\begin{align}
	\left[f(x)g(x) \right]' &= \cRed{f'(x)}g(x) + f(x)\cRed{g'(x)} \label{eq:dpr1} \\
	\dfrac{\dd}{\dx} \left[f(x)g(x) \right] &= \cRed{\dfrac{\dd}{\dx} \left[f(x) \right]}g(x) + f(x)\cRed{\dfrac{\dd}{\dx} \left[g(x)\right]}  \label{eq:dpr2}
\end{align}
\end{arule}

%-=-=-= EXAMPLE
\begin{example}[id:20141209-144203] \label{20141209-144203}\index{Example!20141209-144203} \hfill \\

Differentiate $f(x)=x^2(2x+4)$

\soln

\solnsteps
\begin{align*}
f'(x)&= \farg{x^2(2x+4)}' && \text{SPE} \eqref{eq:spe} \\
f'(x)&= \farg{x^2}'(2x+4)+x^2\farg{2x+4}' && \text{DPr} \eqref{eq:dpr1} \\
f'(x)&= \farg{x^2}'(2x+4)+x^2\farg{2x}+\farg{4}' && \text{DS} \eqref{eq:ds1} \\
f'(x)&= \farg{x^2}'(2x+4)+x^2 \cdot 2 \farg{x}+\farg{4}' && \text{DCM} \eqref{eq:dcm1} \\
f'(x)&= 2x(2x+4)+ x^2\cdot 2 \cdot 1 + \farg{4}' && \text{DPo} \eqref{eq:dpo1} \\
f'(x)&= 2x(2x+4)+ x^2\cdot 2 \cdot 1 + 0 && \text{DC} \eqref{eq:dc1} \\
f'(x)&= 6x^2+8x && \text{simplify} \text{\, goto \,} \, \ref{20141209-145211}\\
\end{align*}
\end{example}


%-=-=-= EXAMPLE
\begin{example}[id:20141209-142321] \label{20141209-142321}\index{Example!20141209-142321} \hfill \\

Differentiate $f(x)=x^2 \cos(x)$

\soln

\solnsteps
\begin{align*}
f'(x)&= \farg{x^2 \cos(x)}' && \text{SPE} \eqref{eq:spe} \\
f'(x)&= \farg{x^2}' \cos(x) + x^2 \farg{\cos (x)}' && \text{DPr} \eqref{eq:dpr1} \\
f'(x)&= 2x \cos (x) + x^2 \farg{\cos (x)}' && \text{DPo} \eqref{eq:dpo1} \\
f'(x)&= 2x \cos (x) + x^2 (-1 \sin(x)) && \text{DCos} \eqref{eq:dcos1} \\
f'(x)&= 2x \cos (x) - x^2 \sin x && \text{OOM} \eqref{eq:oom} 
\end{align*}
\end{example}

%-=-=-= EXAMPLE
\begin{example}[id:20141209-151354] \label{20141209-151354}\index{Example!20141209-151354} \hfill \\

Differentiate $f(x)=\sin(x)\sin(x)$

\soln

\solnsteps
\begin{align*}
f'(x) &= \farg{\sin(x)\sin(x)} && \text{SPE} \eqref{eq:spe} \\
f'(x) &= \farg{\sin(x)}' \sin(x) + \sin (x) \farg{\sin(x)} && \text{DPr} \eqref{eq:dpr1} \\
f'(x) &= \cos(x) \sin(x) + \sin (x) \cos(x) && \text{DSin} \eqref{eq:dsin1} \\
f'(x) &= \cos(x)\sin(x) + \cos(x)\sin(x) && \text{CPM} \eqref{eq:cpm} \\
f'(x) &= 2\cos(x)\sin(x) && \text{OOA} \eqref{eq:ooa} % TODO Simplifying trigonometric expressions
\end{align*}
\end{example}

%-=-=-=-=-=-=-=-=-=-=-=-=-=-=-=-=-=-=-=-=-=-=-=-=
%	SECTION: 
%-=-=-=-=-=-=-=-=-=-=-=-=-=-=-=-=-=-=-=-=-=-=-=-=
\section{Derivative of Rational Functions}


\begin{arule}[Derivative of a Quotient (DQ)]\index{Derivatives! Derivative of a Quotient}
\begin{align}
	\left[ \dfrac{f(x)}{g(x)} \right]' &= \dfrac{f'(x)g(x)-f(x)g'(x)}{\left[g(x)\right]^2} \label{eq:drf1} \\
	\dfrac{\dd}{\dx} \left[ \dfrac{f(x)}{g(x)} \right] & = \dfrac{\dfrac{\dd}{\dx}\left[f(x)\right]g(x)-f(x) \dfrac{\dd}{\dx} \left[g(x)\right] }{\left[g(x)\right]^2}  \label{eq:drf2}
\end{align}
\end{arule}

%-=-=-=-=-=-=-=-=-=-=-=-=-=-=-=-=-=-=-=-=-=-=-=-=
%	SECTION: 
%-=-=-=-=-=-=-=-=-=-=-=-=-=-=-=-=-=-=-=-=-=-=-=-=
\section{Derivative of Trigonometric Functions}

\begin{arule}[Derivative of Sine (DSin)]\index{Derivatives! Derivative of Sine}
\begin{align}
	\left[ \sin(x) \right]' &= \cos (x) \label{eq:dsin1} \\
	\dfrac{\dd}{\dx} \left[ \sin(x) \right] & = \cos (x)  \label{eq:dsin2}
\end{align}
\end{arule}


\begin{arule}[Derivative of Cosine (DCos)]\index{Derivatives! Derivative of Cosine}
\begin{align}
	\left[ \cos(x) \right]' &= -\sin (x) \label{eq:dcos1} \\
	\dfrac{\dd}{\dx} \left[ \cos(x) \right] & = -\sin (x)  \label{eq:dcos2}
\end{align}
\end{arule}

%-=-=-=-=-=-=-=-=-=-=-=-=-=-=-=-=-=-=-=-=-=-=-=-=
%	SECTION: 
%-=-=-=-=-=-=-=-=-=-=-=-=-=-=-=-=-=-=-=-=-=-=-=-=
\section{Derivative of Exponential Functions}

%-=-=-=-=-=-=-=-=-=-=-=-=-=-=-=-=-=-=-=-=-=-=-=-=
%	SECTION: 
%-=-=-=-=-=-=-=-=-=-=-=-=-=-=-=-=-=-=-=-=-=-=-=-=
\section{Derivative of Logarithmic Functions}

\begin{arule}[Derivative of a Natural Logarithm (DNL)]\index{Derivatives! Derivative of a Natural Logarithm}
\begin{align}
	\left[\ln x \right]' &= \frac{1}{x} \label{eq:dnl1} \\
	\dfrac{\dd}{\dx} \left[\ln x \right] &= \frac{1}{x}  \label{eq:dnl2}
\end{align}
\end{arule}




%-=-=-=-=-=-=-=-=-=-=-=-=-=-=-=-=-=-=-=-=-=-=-=-=
%	SECTION: 
%-=-=-=-=-=-=-=-=-=-=-=-=-=-=-=-=-=-=-=-=-=-=-=-=
\section{Derivative of Composite Functions}


\begin{arule}[Derivative of a Composite Function (DComp)]\index{Derivatives! Derivative of a Composite Function}
\begin{align}
	\left[f \left( \cRed{g(x)} \right) \right]' &= \left[\cRed{g(x)} \right]' \left[f\left(\cRed{g(x)}\right)\right]' \label{eq:dcf1} \\
	\dfrac{\dd}{\dx} \left[f \left( \cRed{g(x)} \right) \right] &= \dfrac{\dd}{\dx} \left[\cRed{g(x)} \right] \dfrac{\dd}{\dx} \left[f\left(\cRed{g(x)}\right)\right]  \label{eq:dcf2}
\end{align}
\end{arule}

%-=-=-= EXAMPLE
\begin{example}[id:20141124-203850] \label{20141124-203850}\index{Example!20141124-203850} \hfill \\

Differentiate $y=\ln (3x)$

\soln

\solnsteps

After identifying that $y=\ln (\cRed{3x})$ is a composite function, we let $\cRed{u=3x}$ and thus we get a new function $y=\ln (\cRed{u})$.\\

Using chain rule, 
\begin{align*}
\dydx &= \cBlue{\dydu} \cdot \cPurple{\dudx} && \text{DComp} \eqref{eq:dcf2}
\end{align*}

We need to find the factors $\cBlue{\dydu}$ and $\cPurple{\dudx}$.

\begin{alignat*}{2}
		y &= \ln (\cRed{u})				&\quad	 		\cRed{u}		&= 3x \\
	\cBlue{\dydu} &= \dfrac{1}{u}		&\quad 	 \cPurple{\dudx}	 &= 3 
\end{alignat*}

\begin{align*}
\dydx &= \cBlue{\dydu} \cdot \cPurple{\dudx} && \text{DComp} \eqref{eq:dcf2} \\
\dydx &= \dfrac{1}{\cRed{u}} \cdot \dudx && \text{DNL} \eqref{eq:dnl1} \\
\dydx &= \dfrac{1}{\cRed{u}} \cdot 3 && \text{DPo} \eqref{eq:dpo1} \\
\dydx &= \dfrac{1}{\cRed{3x}} \cdot 3 \\	% TODO variable subsitution 
\dydx &= \dfrac{3}{3x} && \text{OOM} \eqref{eq:oom} \\
\dydx &= \dfrac{1}{x} % TODO reduce fraction
\end{align*}

\qdepend

\qdependlist

example \ref{20141124-205219}-20141124-205219

\end{example}


%-=-=-= EXAMPLE
\begin{example}[id:20141128-160248] \label{20141128-160248}\index{Example!20141128-160248} \hfill \\

Differentiate $y=\ln(3x^2-6x+4)$

\soln

\solnsteps
After identifying that $y=\ln (\cRed{3x^2-6x+4})$ is a composite function, we let \\
$\cRed{u=3x^2-6x+4}$ and thus we get a new function $y=\ln (\cRed{u})$.\\

Using chain rule, 
\begin{align*}
\dydx &= \cBlue{\dydu} \cdot \cPurple{\dudx} && \text{DComp} \eqref{eq:dcf2}
\end{align*}

We need to find the factors $\cBlue{\dydu}$ and $\cPurple{\dudx}$.

\begin{alignat*}{3}
		y &= \ln (\cRed{u})				&\quad	 			\cRed{u} 	&= 3x^2-6x+4 	&\quad\\
	\cBlue{\dydu} &= \dfrac{1}{u}		&\quad 	 \cPurple{\dudx}	 &= 6x-6 		&\quad \text{\, goto \,} \, \ref{20141128-151834}
\end{alignat*}

\begin{align*}
\dydx &= \cBlue{\dydu} \cdot \cPurple{\dudx} && \text{DComp} \eqref{eq:dcf2} \\
\dydx &= \dfrac{1}{\cRed{u}} \cdot \dudx && \text{DNL} \eqref{eq:dnl1} \\
\dydx &= \dfrac{1}{\cRed{u}} \cdot (6x-6) && \text{DPo} \eqref{eq:dpo1} \\
\dydx &= \dfrac{1}{\cRed{3x^2-6x+4}} \cdot (6x-6) \\	% TODO variable subsitution 
\dydx &= \dfrac{6x-6}{3x^2-6x+4} && \text{OOM} \eqref{eq:oom} 
\end{align*}
\end{example}

%-=-=-= EXAMPLE
\begin{example}[id:20141128-155506] \label{20141128-155506}\index{Example!20141128-155506} \hfill \\

Differentiate $y=\ln(\cos x)$

\soln

\solnsteps
After identifying that $y=\ln (\cRed{\cos x})$ is a composite function, we let $\cRed{u=\cos x}$ and thus we get a new function $y=\ln (\cRed{u})$.\\

Using chain rule, 
\begin{align*}
\dydx &= \cBlue{\dydu} \cdot \cPurple{\dudx} && \text{DComp} \eqref{eq:dcf2}
\end{align*}

We need to find the factors $\cBlue{\dydu}$ and $\cPurple{\dudx}$.

\begin{alignat*}{2}
		y &= \ln (\cRed{u})				&\quad	 		\cRed{u}		&= \cos x \\
	\cBlue{\dydu} &= \dfrac{1}{u}		&\quad 	 \cPurple{\dudx}	 &= - \sin x 
\end{alignat*}

\begin{align*}
\dydx &= \cBlue{\dydu} \cdot \cPurple{\dudx} && \text{DComp} \eqref{eq:dcf2} \\
\dydx &= \dfrac{1}{\cRed{u}} \cdot \dudx && \text{DNL} \eqref{eq:dnl1} \\
\dydx &= \dfrac{1}{\cRed{u}} \cdot - \sin x && \text{DPo} \eqref{eq:dpo1} \\
\dydx &= \dfrac{1}{\cRed{\cos x}} \cdot - \sin x \\	% TODO variable subsitution 
\dydx &= \dfrac{-\sin x}{\cos x} && \text{OOM} \eqref{eq:oom} \\
\dydx &= -\tan x % TODO reduce fraction
\end{align*}
\end{example}

%-=-=-= EXAMPLE
\begin{example}[id:20141124-205219] \label{20141124-205219}\index{Example!20141124-205219} \hfill \\

Differentiate $y=\left(x^2-1 \right) \ln(3x)$

\soln

\solnsteps
\begin{align*}
y' &= \farg{x^2-3}' \cdot \ln x + \left(x^2-3\right) \cdot \farg{\ln 3x}' && \text{DPr} \eqref{eq:dpr1} \\
y' &= 2x \cdot \ln x + \left(x^2-3\right) \cdot \farg{\ln 3x}' && \text{differentiate} \text{\, goto \,} \, \ref{20141124-152503} \\
y' &= 2x \cdot \ln x + \left(x^2-3\right) \cdot \dfrac{1}{x}  &&\text{differentiate} \text{\, goto \,} \, \ref{20141124-203850}\\
y' &= 2x \cdot \ln x + \dfrac{x^2-3}{x} && \text{OOM} \eqref{eq:oom} \\
y' &= 2x^2 \ln x +\dfrac{x^2-3}{x} && \text{JTC} \eqref{eq:jtc} \\
y' &= \dfrac{2x^2 \ln x + \left(x^2-3\right)}{x} && \text{OOA} \eqref{eq:ooa} 
\end{align*}
\end{example}




%!TEX root = /Users/markholson/Dropbox/+Projects/LatexFiles/Book/main.tex
%-=-=-=-=-=-=-=-=-=-=-=-=-=-=-=-=-=-=-=-=-=-=-=-=
%
%	CHAPTER 
%
%-=-=-=-=-=-=-=-=-=-=-=-=-=-=-=-=-=-=-=-=-=-=-=-=

\chapterimage{chapter_head_2.pdf} % Chapter heading image

\chapter{Integration}

%-=-=-=-=-=-=-=-=-=-=-=-=-=-=-=-=-=-=-=-=-=-=-=-=
%	SECTION: Integration Rules
%-=-=-=-=-=-=-=-=-=-=-=-=-=-=-=-=-=-=-=-=-=-=-=-=

\section{Integration Rules}

\begin{arule}[Integral of a Constant (IC)]\index{Integrals! Integral of a Constant}
\begin{align}
	\displaystyle \int \! c \, \mathrm{d} x &= cx + k \label{eq:ic}
\end{align}
\end{arule}

\begin{arule}[Integral of a Constant Multiple (ICM)]\index{Integrals! Integral of a Constant Multiple}
\begin{align}
	\displaystyle \int \! \alert{c}x \, \mathrm{d} x &= \alert{c} \displaystyle \int \! x \, \mathrm{d} x \label{eq:icm}
\end{align}
\end{arule}

\begin{arule}[Integral of a Power (IPo)]\index{Integrals! Integral of a Power}
\begin{align}
	\displaystyle \int \! x^n \, \mathrm{d} x &=  \dfrac{x^{n+1}}{n+1}  && \text{for} \, x \ne -1 \label{eq:ipo}
\end{align}
\end{arule}

\begin{arule}[Integral of a Sum (IS)]\index{Integrals! Integral of a Sum}
\begin{align}
	\displaystyle \int \! f(x)+g(x) \, \mathrm{d} x &=  \displaystyle \int \! f(x) \, \mathrm{d} x + \displaystyle \int \! g(x) \, \mathrm{d} x  \label{eq:is}
\end{align}
\end{arule}

\begin{example}[id:20141105-172103] \label{20141105-172103} \index{Example!20141105-172103} \hfill \\
	
Integrate $\displaystyle \int \! \dfrac{1}{4}x^2-2x+4 \, \mathrm{d} x$ 

\soln

\solnsteps
\begin{align*}
\displaystyle \int \! \dfrac{1}{4}x^2-2x+4 \, \mathrm{d} x &= \displaystyle \int \! \dfrac{1}{4}x^2 + \neg 2x+4 \, \mathrm{d} x && \text{DOS} \eqref{eq:dos1} \\
&= \displaystyle \int \! \dfrac{1}{4}x^2 \, \mathrm{d} x + \displaystyle \int \! \neg 2x \, \mathrm{d} x + \displaystyle \int \! 4 \, \mathrm{d} x && \text{IS} \eqref{eq:is} \\
&= \dfrac{1}{4}\displaystyle \int \! x^2 \, \mathrm{d} x + \neg 2\displaystyle \int \! x \, \mathrm{d} x + \displaystyle \int \! 4 \, \mathrm{d} x && \text{ICM} \eqref{eq:icm} \\
&= \dfrac{1}{4}\displaystyle \int \! x^2 \, \mathrm{d} x + \neg 2\displaystyle \int \! x \, \mathrm{d} x + 4(x) + c_1 && \text{IC} \eqref{eq:ic} \\
&= \dfrac{1}{4} \left( \dfrac{x^3}{3} \right) + c_2 + \neg 2 \left(\dfrac{x^2}{2} \right) + c_3 + 4(x) + c_1 && \text{IPo} \eqref{eq:ipo} \\
&= \dfrac{x^3}{12} + c_2 - x^2 + c_3 + 4x + c_1 && \text{OOM} \eqref{eq:oom} \\
&= \dfrac{x^3}{12} - x^2 + 4x + \underbrace{c_1 + c_2 + c_3}_{\text{Just a constant}} && \text{CPA} \eqref{eq:cpa} \\
&= \dfrac{x^3}{12} - x^2 + 4x + C && \text{OOA} \eqref{eq:ooa} 
\end{align*}

\qdepend

\qdependlist

example \ref{20141105-144223}-20141105-144223


\end{example}

\begin{theorem}[Fundamental Theorem of Calculus (FThmC)]\index{Theorems!Fundamental Theorem of Calculus}
If $f$ is continuous on the closed interval $[a, b]$ and $F$ is the indefinite integral of $f$ on $[a, b]$, then

\begin{subequations}
\begin{align}
\displaystyle \int_{a}^{b} \! f(x) \, \mathrm{d} x 	&= F(b)-F(a) \label{eq:fthmc}\\
													&= \left[F(x)\right]_{a}^b 
\end{align}
\end{subequations}
\end{theorem}

%-=-=-=-=-=-=-=-=-=-=-=-=-=-=-=-=-=-=-=-=-=-=-=-=
%	SECTION: Volumes of Revolution
%-=-=-=-=-=-=-=-=-=-=-=-=-=-=-=-=-=-=-=-=-=-=-=-=

\section{Volumes of Revolution}

\begin{definition}[Volume of Revolution: Method of Disks]\index{Method of Disks}
Let $f$ be a nonnegative and continuous function on the closed interval $[a,b]$, then the solid of revolution obtained by rotating the curve $f(x)$ about the $x$-axis from $x=a$ to $x=b$ has volume given by

\begin{align}
V & =\pi \displaystyle \int_{a}^{b} \! [f(x)]^2 \, \mathrm{d} x \label{eq:vrd}
\end{align}	
\end{definition}

%-=-=-= EXAMPLE
\begin{example}[id:20141105-144223]\label{20141105-144223} \index{Example!20141105-144223} \hfill \\
	
Find the volume of the solid formed by rotating the region enclosed by $y=2-\dfrac{x}{2}$ and $y=0$ through $2\pi$ about the $x$-axis on the closed interval $[2, 3]$. 

\soln

\solnsteps
\begin{align*}
V & =\pi \displaystyle \int_{a}^{b} \! [f(x)]^2 \, \mathrm{d} x &&\text{VRD}\eqref{eq:vrd} \\
  & =\pi \displaystyle \int_{2}^{3} \! \farg{2-\dfrac{x}{2}}^2 \, \mathrm{d} x && \text{SPE} \eqref{eq:spe} \\
  &= \pi \displaystyle \int_{2}^{3} \! \dfrac{1}{4} x^2 + \neg 2 x + 4 \, \mathrm{d} x && \text{simplify:} \text{\, goto \,} \, \ref{20141105-161225} \\
  &= \pi \left[ \dfrac{x^3}{12} - x^2 + 4x\right]_{2}^{3} && \text{integrate} \text{\, goto \,} \, \ref{20141105-172103} \\
  &= \pi \left[\dfrac{\farg{3}^3}{12} - \farg{3}^2 + 4\farg{3}- \left(\dfrac{\farg{2}^3}{12} - \farg{2}^2 + 4\farg{2} \right) \right] && \text{FThmC} \eqref{eq:fthmc} \\
  &= \pi \left(\dfrac{7}{12}\right)  &&\text{evaluate: example} \, \ref{20141106-083703} \\
  &= \dfrac{7}{12} \pi && \text{CPM} \eqref{eq:cpm}
\end{align*}

% TODO Add tikz plot

\end{example}

%-=-=-= EXAMPLE
\begin{example}[id:20141106-114907] \label{20141106-114907} \index{Example!20141106-114907} \hfill \\
	
Find the volume of the solid formed by rotating the region enclosed by $y=\sqrt{9-x^2}$ and $y=0$ through $2\pi$ about the $x$-axis on the closed interval $[1, 3]$.

\soln

\solnsteps
\begin{align*}
V & = \pi \displaystyle \int_{a}^{b} \! [f(x)]^2 \, \mathrm{d} x &&\text{VRD}\eqref{eq:vrd} \\
  & = \pi \displaystyle \int_{1}^{3} \! \farg{\sqrt{9-x^2}}^2 \, \mathrm{d} x && \text{SPE} \eqref{eq:spe} \\
  & = \pi \displaystyle \int_{1}^{3} \! 9-x^2 \, \mathrm{d} x && \text{simplify:} \text{\, goto \,} \, \ref{20141107-121834}\\
  & = \pi \left[9x-\dfrac{x^3}{3} \right]_{1}^{3}  &&\text{integrate} \\ % TODO Create integration example
  & = \pi \left[9\farg{3}-\dfrac{\farg{3}^3}{3} - \left(9\farg{1}-\dfrac{\farg{1}^3}{3} \right) \right] && \text{FThmC} \eqref{eq:fthmc} \\
  & = \pi \left(\dfrac{28}{3} \right)  &&\text{evaluate} \\
  & = \dfrac{28}{3}\pi && \text{CPM} \eqref{eq:cpm}  
\end{align*}

% TODO Add tikz plot

\end{example}

%-=-=-= EXAMPLE
\begin{example}[id:20141106-122528] \label{20141106-122528} \index{Example!20141106-122528} \hfill \\
	
Find the volume of the solid formed by rotating the region enclosed by $y=\sqrt{\sin x}$ and $y=0$ through $2\pi$ about the $x$-axis on the closed interval $\left[\frac{\pi}{2}, \pi \right]$.

\soln

\solnsteps
\begin{align*}
V & = \pi \displaystyle \int_{a}^{b} \! [f(x)]^2 \, \mathrm{d} x &&\text{VRD}\eqref{eq:vrd} \\
  & = \pi \displaystyle \int_{\pi/2}^{\pi} \! \farg{\sqrt{\sin} \theta}^2 \, \mathrm{d} \theta && \text{SPE} \eqref{eq:spe} \\
  & = \pi \displaystyle \int_{\pi/2}^{\pi} \! \sin \theta  \, \mathrm{d} \theta && \text{PoPo} \eqref{eq:popo1} \\
  & = \pi \left[ -1 \cos \theta \right]_{\pi/2}^{\pi}  &&\text{integrate} \\
  & = \pi \left[-1\cos \farg{\pi}- \left(-1\cos \farg{\dfrac{\pi}{2}} \right) \right] && \text{FThmC} \eqref{eq:fthmc} \\
  & = \pi (1)  &&\text{evaluate} \\
  & = \pi && \text{MId} \eqref{eq:mid2}  
\end{align*}

% TODO Add tikz plot

\end{example}

%-=-=-= EXAMPLE
\begin{example}[id:20141108-083108] \label{20141108-083108} \index{Example!20141108-083108} \hfill \\

Find the volume of the solid formed by rotating the region enclosed by $y=\sqrt{2-x^2}$ and $y=0$ through $2\pi$ about the $x$-axis.

\soln
\solnsteps


Since no closed interval is stated, we will need to find where the curves $y=\sqrt{2-x^2}$ and $y=0$ intersect to determine the closed interval.  Solving this system of equations (see example \ref{20141107-131748}), we find that $x=\pm \sqrt{2}$.

\begin{align*}
V & = \pi \displaystyle \int_{a}^{b} \! [f(x)]^2 \, \mathrm{d} x &&\text{VRD}\eqref{eq:vrd} \\
&= \pi \displaystyle \int_{-\sqrt{2}}^{\sqrt{2}} \! \farg{\sqrt{2-x^2}}^2 \, \mathrm{d} x \\ % TODO Justify step
&= \pi \displaystyle \int_{-\sqrt{2}}^{\sqrt{2}} \! 2-x^2 \, \mathrm{d} x && \text{PoPo} \eqref{eq:popo1} \\
&= \pi \left[2x-\dfrac{x^3}{3} \right]_{-\sqrt{2}}^{\sqrt{2}}  &&\text{integrate} \\ % TODO add example
&= \pi \left[2\farg{\sqrt{2}}-\dfrac{\farg{\sqrt{2}}^3}{3}-\left(2\farg{-\sqrt{2}}-\dfrac{\farg{-\sqrt{2}}^3}{3} \right) \right] && \text{FThmC} \eqref{eq:fthmc} \\
&= \pi \left(\dfrac{8}{2} \sqrt{2}\right)  &&\text{simplify} \text{\, goto \,} \, \ref{20141107-121834}\\
&= \dfrac{8}{2} \pi \sqrt{2} && \text{CPM} \eqref{eq:cpm} 
\end{align*}

% TODO Add tikz plot

\end{example}

%-=-=-= EXAMPLE
\begin{example}[id:20141125-084541] \label{20141125-084541}\index{Example!20141125-084541} \hfill \\

Find the volume of the solid formed by rotating the region enclosed by $y=\sqrt{2-x^2}$ and $y=0$ through $2\pi$ about the $x$-axis.

\soln
\solnsteps

Since no closed interval is stated, we will need to find where the curves $y=4-x^2$ and $y=0$ intersect to determine the closed interval.  Solving this system of equations (see the similar example \ref{20141107-131748}), we find that $x=\pm 2$. \\


Since the parabolar $y=4-x^2$ is symmetric about the $y$-axis, we can find the volume on the closed interval [0,2] and multiply this by a factor of two to include the volume on the closed interval [-2,0].
 
\soln

\solnsteps
\begin{align*}
V & = 2 \cdot \pi \displaystyle \int_{a}^{b} \! [f(x)]^2 \, \mathrm{d} x &&\text{VRD}\eqref{eq:vrd} \\
&= 2 \cdot \pi \displaystyle \int_{0}^{2} \! \farg{4-x^2}^2 \, \mathrm{d} x \\ % TODO Justify step
&= 2 \cdot \pi \displaystyle \int_{0}^{2} \! x^4-8x^2+16 \, \mathrm{d} x  \\ % TODO Justify step
&= 2 \cdot \pi \left[\dfrac{x^5}{5}-\dfrac{8x^3}{3}+16x \right]_{0}^{2} \\ % TODO Justify step
&= 2 \cdot \pi \left( \dfrac{\farg{2}^5}{5}-\dfrac{8\farg{2}^3}{3}+16\farg{2} \right) && \text{FThmC} \eqref{eq:fthmc} \\
&= \dfrac{512 \pi}{15}
\end{align*}
\end{example}

%!TEX root = /Users/markholson/Dropbox/+Projects/LatexFiles/Book/main.tex
%-=-=-=-=-=-=-=-=-=-=-=-=-=-=-=-=-=-=-=-=-=-=-=-=
%
%	CHAPTER 
%
%-=-=-=-=-=-=-=-=-=-=-=-=-=-=-=-=-=-=-=-=-=-=-=-=

\chapterimage{chapter_head_2.pdf} % Chapter heading image

\chapter{Answers}

\subsection{id:20141120-211843}\label{ans20141120-211843}\index{ans20141120-211843}
\begin{multicols}{4}

\begin{enumerate}
	\item $\sage{a_20141120_211843[0]*x + b_20141120_211843[0]*x}$
	\item $\sage{a_20141120_211843[1]*x + b_20141120_211843[1]*x}$
	\item $\sage{a_20141120_211843[2]*x + b_20141120_211843[2]*x}$
	\item $\sage{a_20141120_211843[3]*x + b_20141120_211843[3]*x}$
	\item $\sage{a_20141120_211843[4]*x + b_20141120_211843[4]*x}$
	\item $\sage{a_20141120_211843[5]*x + b_20141120_211843[5]*x}$
	\item $\sage{a_20141120_211843[6]*x + b_20141120_211843[6]*x}$
	\item $\sage{a_20141120_211843[7]*x + b_20141120_211843[7]*x}$
	\item $\sage{a_20141120_211843[8]*x + b_20141120_211843[8]*x}$
	\item $\sage{a_20141120_211843[9]*x + b_20141120_211843[9]*x}$
	\item $\sage{a_20141120_211843[10]*x + b_20141120_211843[10]*x}$
	\item $\sage{a_20141120_211843[11]*x + b_20141120_211843[11]*x}$
\end{enumerate}
\end{multicols}	

View questions from exercise: \ref{20141120-211843}.

\subsection{id:20141121-195608}\label{ans20141121-195608}\index{ans20141121-195608}
\begin{multicols}{4}

\begin{enumerate}
	\item $\sage{x + b_20141121_195608[0]*x}$
	\item $\sage{x + b_20141121_195608[1]*x}$
	\item $\sage{x + b_20141121_195608[2]*x}$
	\item $\sage{x + b_20141121_195608[3]*x}$
	\item $\sage{x + b_20141121_195608[4]*x}$
	\item $\sage{x + b_20141121_195608[5]*x}$
	\item $\sage{a_20141121_195608[0]*x + x}$
	\item $\sage{a_20141121_195608[1]*x + x}$
	\item $\sage{a_20141121_195608[2]*x + x}$
	\item $\sage{a_20141121_195608[3]*x + x}$
	\item $\sage{a_20141121_195608[4]*x + x}$
	\item $\sage{a_20141121_195608[5]*x + x}$
\end{enumerate}
\end{multicols}	

View questions from exercise: \ref{20141121-195608}.

\subsection{id:20141121-191224}\label{ans20141121-191224}\index{ans20141121-191224}
\begin{multicols}{4}

\begin{enumerate}
	\item $\sage{a_20141121_191224[0]*x - b_20141121_191224[0]*x}$
	\item $\sage{a_20141121_191224[1]*x - b_20141121_191224[1]*x}$
	\item $\sage{a_20141121_191224[2]*x - b_20141121_191224[2]*x}$
	\item $\sage{a_20141121_191224[3]*x - b_20141121_191224[3]*x}$
	\item $\sage{a_20141121_191224[4]*x - b_20141121_191224[4]*x}$
	\item $\sage{a_20141121_191224[5]*x - b_20141121_191224[5]*x}$
	\item $\sage{a_20141121_191224[6]*x - b_20141121_191224[6]*x}$
	\item $\sage{a_20141121_191224[7]*x - b_20141121_191224[7]*x}$
	\item $\sage{a_20141121_191224[8]*x - b_20141121_191224[8]*x}$
	\item $\sage{a_20141121_191224[9]*x - b_20141121_191224[9]*x}$
	\item $\sage{a_20141121_191224[10]*x - b_20141121_191224[10]*x}$
	\item $\sage{a_20141121_191224[11]*x - b_20141121_191224[11]*x}$
\end{enumerate}
\end{multicols}	

View questions from exercise: \ref{20141121-191224}.

\subsection{id:20141123-091420}\label{ans20141123-091420}\index{ans20141123-091420}
\begin{multicols}{4}

\begin{enumerate}
	\item $\sage{a_20141123_091420[0]*x + b_20141123_091420[0]*x + c_20141123_091420[0]*x}$
	\item $\sage{a_20141123_091420[1]*x + b_20141123_091420[1]*x + c_20141123_091420[1]*x}$
	\item $\sage{a_20141123_091420[2]*x + b_20141123_091420[2]*x + c_20141123_091420[2]*x}$
	\item $\sage{a_20141123_091420[3]*x + b_20141123_091420[3]*x + c_20141123_091420[3]*x}$
	\item $\sage{a_20141123_091420[4]*x + b_20141123_091420[4]*x + c_20141123_091420[4]*x}$
	\item $\sage{a_20141123_091420[5]*x + b_20141123_091420[5]*x + c_20141123_091420[5]*x}$
	\item $\sage{a_20141123_091420[6]*x + b_20141123_091420[6]*x + c_20141123_091420[6]*x}$
	\item $\sage{a_20141123_091420[7]*x + b_20141123_091420[7]*x + c_20141123_091420[7]*x}$
	\item $\sage{a_20141123_091420[8]*x + b_20141123_091420[8]*x + c_20141123_091420[8]*x}$
	\item $\sage{a_20141123_091420[9]*x + b_20141123_091420[9]*x + c_20141123_091420[9]*x}$
	\item $\sage{a_20141123_091420[10]*x + b_20141123_091420[10]*x + c_20141123_091420[10]*x}$
	\item $\sage{a_20141123_091420[11]*x + b_20141123_091420[11]*x + c_20141123_091420[11]*x}$
\end{enumerate}
\end{multicols}	

View questions from exercise: \ref{20141123-091420}.
%%-=-=-=-=-=-=-=-=-=-=-=-=-=-=-=-=-=-=-=-=-=-=-=-=
%
%	CHAPTER: 
%
%-=-=-=-=-=-=-=-=-=-=-=-=-=-=-=-=-=-=-=-=-=-=-=-=
\chapter{GY2011}
\newpage
\pagestyle{empty}
\begin{landscape}
\section{GY2011 Local Knowledge Requirement Summaries}
\label{chap:LocalKnowledgeRequirements}


\newcolumntype{A}{ >{\centering\arraybackslash} m{5cm} }
\newcolumntype{B}{ >{\centering\arraybackslash} m{6cm} }
\newcolumntype{C}{ >{\centering\arraybackslash} m{6cm} }
\newcolumntype{D}{ >{\centering\arraybackslash} m{6cm} }
\begin{longtable}{ | A | B | C | D |}
% START HEADER
\hline 
\rowcolor[gray]{.9}   \textbf{Category} &
    \textbf{E Local Summary} &
    \textbf{C Local Summary} &
    \textbf{A Local Summary}\\
  \hline  
\endhead  
% END HEADER


%CONCEPT
\textbf{Concept} &
\alert{Understands} concepts, their purpose and connections between them, illustrated by \alert{some} examples and different representations &
\alert{Thoroughly understands} concepts, their purpose and connections between them, illustrated by \alert{some} examples and different representations &
\alert{Thoroughly understands} concepts, their purpose and connections between them, illustrated by \alert{several} examples and different representations
\\

\hline

%PROCEDURE
\textbf{Procedure} &
Carries out one-step procedures, \alert{some} alternative solution methods &
Carries out \alert{multi-step} procedures&
\alert{Fluently interchanges} solution methods
\\

\hline

%PROBLEM SOLVING
\textbf{Problem Solving} &
Solves problems in \alert{familiar} situations, with \alert{simple} interpretations/formulations &
Solves, formulates and analyses problems with \alert{advanced} interpretations/formulations&
Solves, formulates and analyses \alert{complex problems effectively}
\\

\hline

%MODELING
\textbf{Modeling} &
Uses given models &
\alert{Chooses and applies} appropriate models&
\alert{Adapts} mathematical models and \alert{provides alternative} solution strategies
\\

\hline

%REASONING AND EVALUATION
\textbf{Reasoning \& Evaluation} &
\alert{Simple} evaluation and reasoning, differentiates between guesses and well-founded statements &
\alert{Well-founded} reasoning with \alert{nuanced} judgement and evaluation&
\alert{Well-founded} and \alert{nuanced} reasoning, judgement and evaluation. Finds \alert{general rules and relationships} and expresses these using \alert{algebra and symbols}
\\

\hline

%COMMUNICATIONS
\textbf{Communication} &
Can communicate with \alert{some} certainty. Uses symbols, terms and conventions &
Can communicate with \alert{some} certainty. Uses symbols, terms and conventions \alert{according to purpose and context}&
Can communicate with certainty. Symbols, terms and conventions are \alert{well-suited} to purpose and context
\\

\hline

%CONTEXT AND RELEVANCE
\textbf{Context \& Relevance} &
\alert{Some} contextualization through examples &
\alert{Relevant} contextualization through examples &
\alert{Well-founded} and \alert{nuanced} contextualization through examples 
\\

\hline
\caption{Table caption}
\end{longtable}


\end{landscape}
%%------------------------------------------------

\section{Citation}\index{Citation}

This statement requires citation \cite{book_key}; this one is more specific \cite[122]{article_key}.

%------------------------------------------------

\section{Lists}\index{Lists}

Lists are useful to present information in a concise and/or ordered way\footnote{Footnote example...}.

\subsection{Numbered List}\index{Lists!Numbered List}

\begin{enumerate}
\item The first item
\item The second item
\item The third item
\end{enumerate}

\subsection{Bullet Points}\index{Lists!Bullet Points}

\begin{itemize}
\item The first item
\item The second item
\item The third item
\end{itemize}

\subsection{Descriptions and Definitions}\index{Lists!Descriptions and Definitions}

\begin{description}
\item[Name] Description
\item[Word] Definition
\item[Comment] Elaboration
\end{description}

%----------------------------------------------------------------------------------------
%	CHAPTER 2
%----------------------------------------------------------------------------------------

\chapter{In-text Elements}

\lipsum[1-7] % Dummy text

\section{Theorems}\index{Theorems}

This is an example of theorems.

\subsection{Several equations}\index{Theorems!Several Equations}
This is a theorem consisting of several equations.

\begin{theorem}[Name of the theorem]
In $E=\mathbb{R}^n$ all norms are equivalent. It has the properties:
\begin{align}
& \big| ||\mathbf{x}|| - ||\mathbf{y}|| \big|\leq || \mathbf{x}- \mathbf{y}||\\
&  ||\sum_{i=1}^n\mathbf{x}_i||\leq \sum_{i=1}^n||\mathbf{x}_i||\quad\text{where $n$ is a finite integer}
\end{align}
\end{theorem}

\subsection{Single Line}\index{Theorems!Single Line}
This is a theorem consisting of just one line.

\begin{theorem}
A set $\mathcal{D}(G)$ in dense in $L^2(G)$, $|\cdot|_0$. 
\end{theorem}

%------------------------------------------------

\section{Definitions}\index{Definitions}

This is an example of a definition. A definition could be mathematical or it could define a concept.

\begin{definition}[Definition name]
Given a vector space $E$, a norm on $E$ is an application, denoted $||\cdot||$, $E$ in $\mathbb{R}^+=[0,+\infty[$ such that:
\begin{align}
& ||\mathbf{x}||=0\ \Rightarrow\ \mathbf{x}=\mathbf{0}\\
& ||\lambda \mathbf{x}||=|\lambda|\cdot ||\mathbf{x}||\\
& ||\mathbf{x}+\mathbf{y}||\leq ||\mathbf{x}||+||\mathbf{y}||
\end{align}
\end{definition}

%------------------------------------------------

\section{Notations}\index{Notations}

\begin{notation}
Given an open subset $G$ of $\mathbb{R}^n$, the set of functions $\varphi$ are:
\begin{enumerate}
\item Bounded support $G$;
\item Infinitely differentiable;
\end{enumerate}
a vector space is denoted by $\mathcal{D}(G)$. 
\end{notation}

%------------------------------------------------

\section{Remarks}\index{Remarks}

This is an example of a remark.

\begin{remark}
The concepts presented here are now in conventional employment in mathematics. Vector spaces are taken over the field $\mathbb{K}=\mathbb{R}$, however, established properties are easily extended to $\mathbb{K}=\mathbb{C}$.
\end{remark}

%------------------------------------------------

\section{Corollaries}\index{Corollaries}

This is an example of a corollary.

\begin{corollary}[Corollary name]
The concepts presented here are now in conventional employment in mathematics. Vector spaces are taken over the field $\mathbb{K}=\mathbb{R}$, however, established properties are easily extended to $\mathbb{K}=\mathbb{C}$.
\end{corollary}

%------------------------------------------------

\section{Propositions}\index{Propositions}

This is an example of propositions.

\subsection{Several equations}\index{Propositions!Several Equations}

\begin{proposition}[Proposition name]
It has the properties:
\begin{align}
& \big| ||\mathbf{x}|| - ||\mathbf{y}|| \big|\leq || \mathbf{x}- \mathbf{y}||\\
&  ||\sum_{i=1}^n\mathbf{x}_i||\leq \sum_{i=1}^n||\mathbf{x}_i||\quad\text{where $n$ is a finite integer}
\end{align}
\end{proposition}

\subsection{Single Line}\index{Propositions!Single Line}

\begin{proposition} 
Let $f,g\in L^2(G)$; if $\forall \varphi\in\mathcal{D}(G)$, $(f,\varphi)_0=(g,\varphi)_0$ then $f = g$. 
\end{proposition}

%------------------------------------------------

\section{Examples}\index{Examples}

This is an example of examples.

\subsection{Equation and Text}\index{Examples!Equation and Text}

\begin{example}
Let $G=\{x\in\mathbb{R}^2:|x|<3\}$ and denoted by: $x^0=(1,1)$; consider the function:
\begin{equation}
f(x)=\left\{\begin{aligned} & \mathrm{e}^{|x|} & & \text{si $|x-x^0|\leq 1/2$}\\
& 0 & & \text{si $|x-x^0|> 1/2$}\end{aligned}\right.
\end{equation}
The function $f$ has bounded support, we can take $A=\{x\in\mathbb{R}^2:|x-x^0|\leq 1/2+\epsilon\}$ for all $\epsilon\in\intoo{0}{5/2-\sqrt{2}}$.
\end{example}

\subsection{Paragraph of Text}\index{Examples!Paragraph of Text}

\begin{example}[Example name]
\lipsum[2]
\end{example}

%------------------------------------------------

\section{Exercises}\index{Exercises}

This is an example of an exercise.

\begin{exercise}
Simplify the following expressions.
\begin{multicols}{3}
\begin{enumerate}
\item $5x+4$
\columnbreak
\item $5x+4$
\columnbreak
\item $5x+4$
\end{enumerate}
\end{multicols}
\end{exercise}

\begin{exercise}
This is a good place to ask a question to test learning progress or further cement ideas into students' minds.
\end{exercise}

\begin{exercise}
This is a good place to ask a question to test learning progress or further cement ideas into students' minds.
\end{exercise}

%------------------------------------------------

\section{Problems}\index{Problems}

\begin{problem}
What is the average airspeed velocity of an unladen swallow?
\end{problem}

%------------------------------------------------

\section{Vocabulary}\index{Vocabulary}

Define a word to improve a students' vocabulary.

\begin{vocabulary}[Word]
Definition of word.
\end{vocabulary}

%----------------------------------------------------------------------------------------
%	CHAPTER 3
%----------------------------------------------------------------------------------------

\chapterimage{chapter_head_1.pdf} % Chapter heading image

\chapter{Presenting Information}

\section{Table}\index{Table}

\begin{table}[h]
\centering
\begin{tabular}{l l l}
\toprule
\textbf{Treatments} & \textbf{Response 1} & \textbf{Response 2}\\
\midrule
Treatment 1 & 0.0003262 & 0.562 \\
Treatment 2 & 0.0015681 & 0.910 \\
Treatment 3 & 0.0009271 & 0.296 \\
\bottomrule
\end{tabular}
\caption{Table caption}
\end{table}

%------------------------------------------------

\section{Figure}\index{Figure}

\begin{figure}[h]
\centering\includegraphics[scale=0.5]{placeholder}
\caption{Figure caption}
\end{figure}


%----------------------------------------------------------------------------------------
%	BIBLIOGRAPHY
%----------------------------------------------------------------------------------------

\chapter*{Bibliography}
\addcontentsline{toc}{chapter}{\textcolor{ocre}{Bibliography}}
\section*{Website}
\addcontentsline{toc}{section}{Websites}
\printbibliography[heading=bibempty,type=misc]

%----------------------------------------------------------------------------------------
%	INDEX
%----------------------------------------------------------------------------------------

\cleardoublepage
\phantomsection
\setlength{\columnsep}{0.75cm}
\addcontentsline{toc}{chapter}{\textcolor{ocre}{Index}}
\printindex

%----------------------------------------------------------------------------------------

\end{document}
